\documentclass[11pt,fontset=none,a4paper,oneside]{ctexbook}

% ---- 中文字体：使用系统字体，保证 PDF 文本可检索、可复制 ----
\setCJKmainfont{Songti SC}[BoldFont={Songti SC Bold}]
\setCJKsansfont{Heiti SC Medium}
\setCJKfamilyfont{zhhei}{Heiti SC Medium}
\providecommand{\heiti}{\CJKfamily{zhhei}}

%========================================================
%  宏包
%========================================================
\usepackage{amsmath,amssymb,amsthm,mathrsfs}
\usepackage{mathtools}
\usepackage[margin=2.5cm,top=2.6cm,bottom=2.6cm]{geometry}
\usepackage{xcolor}
\usepackage[most]{tcolorbox}
\usepackage{booktabs,longtable,array,tabularx,multirow}
\newcolumntype{L}[1]{>{\raggedright\arraybackslash}p{#1}}
\newcolumntype{Y}{>{\raggedright\arraybackslash}X}
\usepackage{enumitem}
\usepackage{tikz}
\usetikzlibrary{arrows.meta,positioning,calc,fit,backgrounds}
\usepackage{pgfplots}
\pgfplotsset{compat=1.18}
\usepackage{fancyhdr}
\usepackage{hyperref}

%========================================================
%  颜色
%========================================================
\definecolor{mainblue}{RGB}{18,62,120}
\definecolor{softblue}{RGB}{233,241,251}
\definecolor{warnred}{RGB}{150,38,38}
\definecolor{softred}{RGB}{253,237,237}
\definecolor{okgreen}{RGB}{26,104,76}
\definecolor{softgreen}{RGB}{234,246,240}
\definecolor{warmorange}{RGB}{158,84,18}
\definecolor{softorange}{RGB}{255,245,230}
\definecolor{notegray}{RGB}{82,82,88}
\definecolor{softgray}{RGB}{243,243,246}
\definecolor{rulegray}{RGB}{150,155,165}

%========================================================
%  版式
%========================================================
\linespread{1.32}
\setlength{\parskip}{0.35em}
\setlength{\parindent}{2em}

\ctexset{
  chapter = {
    format = \raggedright\huge\heiti\color{mainblue},
    name = {第,章},
    number = \arabic{chapter},
    beforeskip = 0.8em,
    afterskip = 1.6em,
  },
  section = {
    format = \large\heiti\color{mainblue},
    beforeskip = 1.6em,
    afterskip = 0.8em,
  },
  subsection = {
    format = \normalsize\bfseries\color{notegray},
    beforeskip = 1.1em,
    afterskip = 0.5em,
  },
}

\pagestyle{fancy}
\fancyhf{}
\fancyhead[L]{\small\color{rulegray}\nouppercase{\leftmark}}
\fancyhead[R]{\small\color{rulegray}\thepage}
\renewcommand{\headrulewidth}{0.3pt}
\renewcommand{\headrule}{\hbox to\headwidth{\color{rulegray}\leaders\hrule height \headrulewidth\hfill}}
\fancypagestyle{plain}{\fancyhf{}\fancyhead[R]{\small\color{rulegray}\thepage}\renewcommand{\headrulewidth}{0pt}}

%========================================================
%  定理环境
%========================================================
\theoremstyle{definition}
\newtheorem{defn}{定义}[chapter]
\newtheorem{examp}[defn]{例}
\theoremstyle{plain}
\newtheorem{thm}[defn]{定理}
\newtheorem{prop}[defn]{命题}
\newtheorem{lem}[defn]{引理}
\theoremstyle{remark}
\newtheorem{remark}[defn]{说明}
\numberwithin{equation}{chapter}

%========================================================
%  提示盒
%========================================================
\newtcolorbox{ideabox}[1][直觉]{enhanced,breakable,
  colback=softblue,colframe=mainblue!60,colbacktitle=mainblue!78,coltitle=white,
  fonttitle=\small\bfseries,title={#1},boxrule=0pt,arc=1.5pt,
  left=9pt,right=9pt,top=6pt,bottom=6pt,before skip=9pt,after skip=9pt}

\newtcolorbox{warnbox}[1][讲解禁区]{enhanced,breakable,
  colback=softred,colframe=warnred!55,colbacktitle=warnred!80,coltitle=white,
  fonttitle=\small\bfseries,title={#1},boxrule=0pt,arc=1.5pt,
  left=9pt,right=9pt,top=6pt,bottom=6pt,before skip=9pt,after skip=9pt}

\newtcolorbox{quizbox}[1][自测]{enhanced,breakable,
  colback=softgreen,colframe=okgreen!55,colbacktitle=okgreen!80,coltitle=white,
  fonttitle=\small\bfseries,title={#1},boxrule=0pt,arc=1.5pt,
  left=9pt,right=9pt,top=6pt,bottom=6pt,before skip=9pt,after skip=9pt}

\newtcolorbox{scriptbox}[1][脚本对照]{enhanced,breakable,
  colback=softorange,colframe=warmorange!55,colbacktitle=warmorange!80,coltitle=white,
  fonttitle=\small\bfseries,title={#1},boxrule=0pt,arc=1.5pt,
  left=9pt,right=9pt,top=6pt,bottom=6pt,before skip=9pt,after skip=9pt}

\newtcolorbox{notebox}[1][记一笔]{enhanced,breakable,
  colback=softgray,colframe=notegray!40,colbacktitle=notegray!70,coltitle=white,
  fonttitle=\small\bfseries,title={#1},boxrule=0pt,arc=1.5pt,
  left=9pt,right=9pt,top=6pt,bottom=6pt,before skip=9pt,after skip=9pt}

%========================================================
%  记号
%========================================================
\newcommand{\F}{\mathcal{F}}
\newcommand{\R}{\mathbb{R}}
\newcommand{\C}{\mathbb{C}}
\newcommand{\D}{\mathcal{D}}
\newcommand{\Sv}{\mathcal{S}}
\newcommand{\dd}{\mathop{}\!\mathrm{d}}
\newcommand{\PV}{\operatorname{PV}}
\newcommand{\sgn}{\operatorname{sgn}}
\newcommand{\supp}{\operatorname{supp}}
\newcommand{\ip}[2]{\langle #1,#2\rangle}
\newcommand{\Lone}{L^1(\R)}
\newcommand{\Ltwo}{L^2(\R)}
\newcommand{\Lonloc}{L^1_{\mathrm{loc}}(\R)}
\newcommand{\ep}{\varepsilon}
\newcommand{\fh}{\widehat{f}}
\newcommand{\gh}{\widehat{g}}
\newcommand{\ph}{\widehat{\varphi}}
\newcommand{\hf}{\widehat{H}}

\hypersetup{
  colorlinks=true,
  linkcolor=mainblue,
  urlcolor=mainblue,
  citecolor=mainblue,
  bookmarksnumbered=true,
}

%========================================================
\begin{document}

%--------------------------------------------------------
%  封面
%--------------------------------------------------------
\begin{titlepage}
\thispagestyle{empty}
\centering
\vspace*{2.2cm}

{\color{rulegray}\rule{\textwidth}{0.6pt}}

\vspace{1.4cm}

{\Huge\bfseries\color{mainblue} 傅里叶变换的另一面}

\vspace{0.5cm}

{\LARGE\bfseries\color{mainblue} 创作者进阶读本}

\vspace{1.0cm}

{\large 六期系列的知识地图、完整推导与讲解边界}

\vspace{0.4cm}

{\large 从零频率的直觉，到温和分布中的严格定义}

\vspace{1.2cm}

\begin{tcolorbox}[colback=softblue,colframe=mainblue!40,boxrule=0pt,arc=2pt,width=0.86\textwidth]
\centering
\vspace{0.2em}
{\large 贯穿全书的三个结果}\\[0.9em]
$\F[1]=2\pi\delta(\omega)$\\[0.45em]
$\F[t^{\,n}]=2\pi i^{\,n}\delta^{(n)}(\omega)$\\[0.45em]
$\F[H]=\pi\delta(\omega)+\PV\dfrac{1}{i\omega}$
\vspace{0.2em}
\end{tcolorbox}

\vspace{1.0cm}

{\normalsize
\begin{tabular}{rl}
\textbf{配套} & \texttt{background\_knowledge.pdf}（骨架版）\\[0.3em]
\textbf{本册} & 同一套内容的扩写：补动机、补中间步骤、补边界\\[0.3em]
\textbf{用途} & 开拍前通读；写脚本时按章查阅\\[0.3em]
\textbf{版本} & 2026 年 11 月 1 日\\
\end{tabular}}

\vfill

{\small\color{notegray}
全书统一使用角频率约定，配对不含复共轭。\\
不同资料可能采用不同的 $2\pi$ 归一化，引用公式前请先核对约定。}

\vspace{0.8cm}
{\color{rulegray}\rule{\textwidth}{0.6pt}}
\end{titlepage}

%--------------------------------------------------------
%  摘要页
%--------------------------------------------------------
\frontmatter
\thispagestyle{plain}
\vspace*{0.5cm}

\begin{ideabox}[怎么用这本书（先读这一页）]
\small
\textbf{一句话：}这不是一本傅里叶分析教材，而是你拍这六期视频的“后台”。

\vspace{0.4em}
\textbf{它回答三类问题。}
\begin{enumerate}[leftmargin=1.6em,itemsep=0.25em,topsep=0.3em]
\item \textbf{这几期到底在讲什么？} —— 第 0 章。六集的分工、顺序、依赖关系、以及那个把所有内容串起来的一条线。
\item \textbf{那些公式是怎么来的？} —— 第 1--6 章。每一步都写清楚“为什么这一步合法”，而不只是给结果。
\item \textbf{哪些话能说、哪些不能说？} —— 每章末的「讲解禁区」与附录 B。这是最省时间的一部分：它把最容易犯的错直接列出来。
\end{enumerate}

\vspace{0.4em}
\textbf{建议的三遍读法。}
\begin{enumerate}[leftmargin=1.6em,itemsep=0.25em,topsep=0.3em]
\item \textbf{第一遍（两小时）}：只读第 0 章，加上每章开头的“一句话”和“观众的问题链”。目的是拿到全局地图。
\item \textbf{第二遍（分集）}：动手重算每章的推导。看不懂就停在那里，把卡住的那一步标出来——那通常正是视频里最容易糊弄过去、也最容易翻车的地方。
\item \textbf{第三遍（开拍前）}：只看每章的「讲解禁区」、附录 B 的对照表和附录 D 的练习题。
\end{enumerate}

\vspace{0.4em}
\textbf{一条纪律。}三层内容要分清：
\begin{center}
\small
\begin{tabular}{@{}L{2.0cm}L{5.3cm}L{5.4cm}@{}}
\toprule
层次 & 讲的是什么 & 用在哪里 \\
\midrule
直觉 & 为什么这个结论应该是对的 & 视频开场、动画旁白 \\
计算 & 结果一步一步怎么算出来 & 视频主体 \\
证明 & 为什么这些运算合法、条件是什么 & 视频可省略；\textbf{你不能不知道} \\
\bottomrule
\end{tabular}
\end{center}

视频可以省略证明，但不能因为省略了证明，就以为那些计算没有条件。
\end{ideabox}

\vspace{0.3cm}

\vspace{0.2cm}
{\noindent\bfseries\color{mainblue} 全书地图}\par

\vspace{0.15cm}
{\footnotesize\noindent
\renewcommand{\arraystretch}{1.18}
\begin{tabularx}{\textwidth}{@{}L{1.1cm}Y L{2.4cm}@{}}
\toprule
章 & 内容 & 什么时候读 \\
\midrule
0 & 导读：全局地图、三条暗线、约定、四个空间 & 开拍前必读 \\
1 & EP01　常数为什么对应 $\delta$ & 拍 01 前 \\
2 & EP02　多项式与 $\delta$ 的导数 & 拍 02 前 \\
3 & EP03　阶跃、主值与多出来的 $\delta$ & 拍 03 前 \\
4 & EP04　收敛方式与近似恒等核 & 拍 04 前 \\
5 & EP05　测试函数与连续线性泛函 & 拍 05 前 \\
6 & EP06　$\Sv$、$\Sv'$ 与对偶定义 & 拍 06 前 \\
7 & 把六集缝起来：总表与制作落地 & 全部拍完 \\
附 & A 公式总表／B 能说不能说／C 术语／D 练习 & 随时 \\
\bottomrule
\end{tabularx}}

\setcounter{tocdepth}{1}
\tableofcontents

\mainmatter
\setcounter{chapter}{-1}

\IfFileExists{ch0.tex}{\chapter{导读：六期到底在讲什么}

\begin{ideabox}[这一章的作用]
如果你只有两小时，读这一章就够了。它不推导任何公式，只回答四个问题：
\begin{enumerate}[leftmargin=1.6em,itemsep=0.2em,topsep=0.2em]
\item 这六期合起来在回答哪一个问题？
\item 为什么这个问题必须拆成六集，而且必须按这个顺序？
\item 六集之间靠什么串起来？（三条暗线）
\item 每一集的结论，是在哪一集被真正证明的？（欠账与还账）
\end{enumerate}
读完这一章，你应该能不看笔记，用三分钟把整个系列讲给一个朋友听。
\end{ideabox}

\section{这一系列要回答的唯一问题}

傅里叶变换的通常说法是：\textbf{看一个信号里包含哪些频率，各占多少}。这句话在工程上极其好用，因为大多数真实信号都满足一个不起眼的前提——它们“足够快地衰减”，所以那个积分真的收敛。

现在把最普通的一个信号放进去：永远等于 $1$ 的常数。

\begin{itemize}[leftmargin=1.6em,itemsep=0.3em]
\item 从直觉上说，它完全不变，所以它\textbf{只包含零频率}。这个直觉几乎不可能错。
\item 从计算上说，把 $f(t)=1$ 代进定义，得到 $\int_{\R}e^{-i\omega t}\dd t$，这个积分对任何 $\omega$ 都不收敛。
\end{itemize}

直觉与计算在这里正面冲突。冲突不是出在“频率”这个想法上，而是出在“积分”这个工具上：积分只肯为衰减得足够快的信号工作，而常数一点都不衰减。

于是整个系列真正要回答的问题是：

\begin{center}
\begin{tcolorbox}[colback=softblue,colframe=mainblue!50,boxrule=0pt,arc=2pt,width=0.93\textwidth]
\centering
\textbf{当一个信号的傅里叶积分按普通意义根本不存在时，\\[0.25em]
我们还能不能在“频率”这套语言里继续谈论它？\\[0.25em]
如果能，那句话的数学内容到底是什么？}
\end{tcolorbox}
\end{center}

后面六集的所有内容，都是这一个问题的展开。三个“主结果”是三条不同的进攻路线：

\begin{align}
\F[1]&=2\pi\delta(\omega), \label{eq:main1}\\
\F[t^{n}]&=2\pi i^{\,n}\delta^{(n)}(\omega),\qquad n=0,1,2,\dots \label{eq:main2}\\
\F[H]&=\pi\delta(\omega)+\PV\frac{1}{i\omega}. \label{eq:main3}
\end{align}

它们分别对应三类“不衰减”的信号：\textbf{完全不变}（常数）、\textbf{越涨越快}（多项式）、\textbf{突然跳一下然后停住}（阶跃）。这三类恰好是普通傅里叶积分最无能为力的三类。

\begin{notebox}[这三条式子为什么值得拍六集]
不是因为它们难算，而是因为它们各自暴露了一种不同的“失败方式”：
\begin{itemize}[leftmargin=1.4em,itemsep=0.15em,topsep=0.2em]
\item 常数暴露的是\textbf{振幅不衰减}：被积函数的模恒等于 $1$，振荡再剧烈也不会抵消。
\item 多项式暴露的是\textbf{增长更快}：$|t|^{n}$ 越来越大，连振荡都压不住。
\item 阶跃暴露的是\textbf{局部奇性}：积分在无穷远处勉强能救，却在 $\omega=0$ 处留下一个不可积的 $1/\omega$。
\end{itemize}
三种病，三种药。如果只用一种药，观众学不到东西。
\end{notebox}

\section{一张“病历表”：四类对象为什么都算不出来}

先做一个冷静的诊断。把定义式 $\fh(\omega)=\int_{\R}f(t)e^{-i\omega t}\dd t$ 拿去试，看看每一类对象是怎么失败的。

\begin{center}
\small
\renewcommand{\arraystretch}{1.5}
\begin{tabularx}{\textwidth}{@{}L{1.5cm}L{4.6cm}L{5.4cm}L{3.2cm}@{}}
\toprule
对象 & 代入定义后发生什么 & 卡在哪里 & 广义结果 \\
\midrule
$1$ & $\int_{-T}^{T}e^{-i\omega t}\dd t=\dfrac{2\sin(\omega T)}{\omega}$（$\omega\neq0$） & 振幅 $2/|\omega|$ 不随 $T$ 衰减，永远在振荡；$\omega=0$ 处则 $=2T\to\infty$ & $2\pi\delta$ \\
$t^{n}$ & 被积函数的模是 $|t|^{n}$，无界增长 & 比常数更糟：连“振荡抵消”的指望都没有 & $2\pi i^{n}\delta^{(n)}$ \\
$H$ & $\int_{0}^{T}e^{-i\omega t}\dd t=\dfrac{1-e^{-i\omega T}}{i\omega}$ & 端点项 $e^{-i\omega T}$ 不收敛；即使取极限，$1/\omega$ 在原点不可积 & $\pi\delta+\PV\dfrac{1}{i\omega}$ \\
$\delta$ & 根本不是函数，代不进去 & 没有“每点取值”这回事 & $1$ \\
\bottomrule
\end{tabularx}
\end{center}

这张表里最要紧的一行其实是最后一列的\textbf{共同点}：

\begin{itemize}[leftmargin=1.6em,itemsep=0.3em]
\item 这四类对象\textbf{都不属于} $\Lone$。所以“$\int|f|<\infty$ 则傅里叶积分绝对收敛”这条定理一律用不上。
\item 它们\textbf{也都不属于} $\Ltwo$。所以 Plancherel 定理那条“保能量”的退路也走不通。
\item 但它们\textbf{全都属于} $\Sv'$（温和分布）。在这个空间里，傅里叶变换是一个连续可逆的线性算子，三条式子同时成立。
\end{itemize}

\begin{ideabox}[这就是整个系列的“剧情大纲”]
第一幕（EP01--EP03）：\textbf{看见病}。用三个具体信号，让观众亲手撞上三堵墙。

第二幕（EP04）：\textbf{停下来问诊断标准}。“我们说的极限，到底是哪一种极限？”——从“算出有趣结果”转向“研究结果为何合理”。

第三幕（EP05--EP06）：\textbf{换一间更大的手术室}。定义分布，在 $\Sv'$ 上重做傅里叶变换，然后回头把前三集的账全部结清。
\end{ideabox}

\section{六集分工：谁提出问题，谁负责收账}

\begin{center}
\small
\renewcommand{\arraystretch}{1.45}
\begin{tabularx}{\textwidth}{@{}L{1.0cm}L{4.2cm}L{4.4cm}L{4.9cm}@{}}
\toprule
集 & 观众的问题 & 本集真正建立的东西 & 故意留给后面的 \\
\midrule
EP01 & 一条水平线的频谱为什么是 $\delta$？ & 零频率直觉；经典积分失败的现场；双边指数正则化；面积集中 & $\delta$ 的严格定义；极限的严格含义 \\
EP02 & 不断增长的多项式也有频谱？ & 时域乘 $t$ $\leftrightarrow$ 频域求导；$\delta$ 的导数靠“取样导数”理解 & 分布导数的合法性（借用规则） \\
EP03 & 阶跃的频谱为什么多一个 $\delta$？ & $H'=\delta$；主值；微分规则为什么会丢信息 & 主值极限的完整估计；歧义为何恰好是 $c\delta$ \\
EP04 & 点值都趋于零，积分为什么还在？ & 三种收敛；控制收敛；近似恒等核；质量集中 & 用“作用”定义对象的合法性 \\
EP05 & $\delta$ 到底是不是函数？ & 测试函数；连续线性泛函；嵌入；导数与光滑乘法 & 傅里叶变换怎么办 \\
EP06 & 发散的积分怎样被严格定义？ & $\Sv$ 与 $\Sv'$；对偶定义；连续性；统一证明 & （本系列不展开）PDE、卷积代数 \\
\bottomrule
\end{tabularx}
\end{center}

\subsection*{每一集在观众脑子里留下的那一句话}

把六句话连起来读，就是整个系列。

\begin{center}
\small
\renewcommand{\arraystretch}{1.4}
\begin{tabularx}{\textwidth}{@{}L{1.3cm}Y@{}}
\toprule
EP01 & 信号在时间上铺得越开，频谱就聚得越紧；常数的频谱全部堆在零频率，权重是 $2\pi$。 \\
EP02 & 在时间上乘一个 $t$，等于在频率上求一次导；所以多项式只有零频率，但它记录的是零点的各阶导数。 \\
EP03 & 微分规则只能定出“除掉零频率以外”的部分；零频率那一项被乘 $\omega$ 抹掉了，必须另外补回来。 \\
EP04 & “每一点都不为零”与“整体不为零”是两件事；核之所以变成 $\delta$，是因为它的质量向一点集中，而不是因为它变小。 \\
EP05 & 一个对象可以不由“每点取值”定义，而由“它怎样作用于探针”定义；$\delta$ 就是这样被定义的。 \\
EP06 & 只要把探针换成 $\Sv$，傅里叶变换就能整体搬到分布上，并且保持极限；前三集的所有计算在这一步同时获得合法性。 \\
\bottomrule
\end{tabularx}
\end{center}

\subsection*{为什么必须是这个顺序}

顺序不是按数学难度排的，而是按\textbf{悬念}排的。每一集结束时留下的问题，正好是下一集的开场。

\begin{enumerate}[leftmargin=2em,itemsep=0.45em]
\item \textbf{EP01 结尾}：常数对应 $2\pi\delta$。可是 $\delta$ 不是一个函数，凭什么能出现在公式右边？——先按下不表，因为观众此刻还没有非问不可的理由。
\item \textbf{EP02 结尾}：既然乘 $t$ 可以，那对 $t$ 求导会怎样？更尖锐的是，$\delta'$ 是什么？“对一个不是函数的东西求导”这句话本身就该被质疑。到这里，观众第一次\textbf{主动}想问 $\delta$ 是什么。
\item \textbf{EP03 结尾}：我们一直在说“趋于一个尖峰”“趋于主值”。这些“趋于”到底是哪一种趋近？为什么可以把极限搬进积分？——问题从“对象是什么”升级为“命题是什么意思”。
\item \textbf{EP04 结尾}：不同的尖峰形状给出同一个取样作用。既然如此，能不能干脆\textbf{用这个作用来定义}那个极限对象？
\item \textbf{EP05 结尾}：$\delta$ 定义好了，导数也定义好了。但傅里叶变换呢？把一个分布塞进普通积分显然不行。
\item \textbf{EP06 结尾}：全部结清，回到第一集那条水平线。
\end{enumerate}

\begin{scriptbox}[写作时可以直接用的一句总纲]
\textbf{用具体信号提出问题，用动画和计算建立理解，再逐步走向分布理论，解释那些普通傅里叶积分无法直接处理的结果。}
\end{scriptbox}

\section{三条暗线：观众看的是六个问题，你铺的是三条线索}

明线是六个具体信号。但真正让系列“成为一个系列”而不是“六个科普视频”的，是三条贯穿始终的暗线。写脚本时，每隔一集就该让其中一条线往前挪一格。

\subsection*{暗线一：正则化——先让积分收敛，再让参数消失}

这是全系列最常用的手法，五集里出现了五次，而且每次外形都不一样：

\begin{center}
\small
\renewcommand{\arraystretch}{1.45}
\begin{tabularx}{\textwidth}{@{}L{2.9cm}L{2.5cm}L{4.6cm}L{4.4cm}@{}}
\toprule
参数族 & 时域极限 & 频域极限 & 出现于 \\
\midrule
$e^{-\ep|t|}$ & $\to 1$（逐点） & $2\ep/(\ep^{2}+\omega^{2})\to 2\pi\delta$ & EP01 \\
$e^{-\ep t}H(t)$ & $\to H$（逐点） & $1/(\ep+i\omega)\to\pi\delta-i\PV(1/\omega)$ & EP03 \\
$e^{-\sigma t^{2}/2}$ & $\to 1$（逐点） & $\sqrt{2\pi/\sigma}\,e^{-\omega^{2}/(2\sigma)}\to 2\pi\delta$ & EP06 \\
$\mathbf{1}_{[-T,T]}$ & $\to 1$（逐点） & $2\sin(T\omega)/\omega\to 2\pi\delta$ & EP01 提过 / EP06 收尾 \\
$K_{\ep}=\frac{1}{\pi}\frac{\ep}{\ep^{2}+x^{2}}$ & （无普通极限） & $\to\delta$ & EP01 借用 / EP04 证明 \\
\bottomrule
\end{tabularx}
\end{center}

请盯住这张表的第二列和第三列：\textbf{五个完全不同的函数族，在频域收敛到同一个对象}（差一个 $2\pi$ 的比例）。指数核、高斯核、矩形窗，一个比一个不像，极限却是同一个 $\delta$。

\begin{ideabox}[这条暗线的意义]
它直接给出了“$\delta$ 不该由形状定义”的全部理由。如果五种形状都通向同一个终点，那么这个终点就不能是“某种特别的形状”，它只能是\textbf{一种共同的作用}。

这就是 EP04 结尾的问题，也是 EP05 开场的入口。观众不需要被告知“分布理论很重要”，他们会自己得出结论。
\end{ideabox}

\subsection*{暗线二：从“代入积分”到“配对”}

整部系列在偷偷更换“计算”这个词的含义：

\begin{center}
\small
\renewcommand{\arraystretch}{1.4}
\begin{tabularx}{\textwidth}{@{}L{1.5cm}L{4.4cm}L{9.0cm}@{}}
\toprule
集数 & 计算的形式 & 这句话在做什么 \\
\midrule
EP01 & $\displaystyle\int_{\R}K_{\ep}(\omega)\varphi(\omega)\dd\omega\to\varphi(0)$ & 借一个平滑探针 $\varphi$，观察核的作用（先说清楚，EP04 才证明）\\
EP04 & 同上，但给出证明 & 用“近处靠连续性、远处靠尾部”的估计把极限算出来 \\
EP05 & $\ip{\delta}{\varphi}=\varphi(0)$ & 把这个作用\textbf{提升为定义}：$\delta$ 就是“取 $\varphi(0)$”这个泛函 \\
EP06 & $\ip{\F T}{\varphi}:=\ip{T}{\F\varphi}$ & 把变换也搬到这个框架里：不知道 $T$ 是什么，但知道它对每个探针的响应 \\
\bottomrule
\end{tabularx}
\end{center}

从“算一个数”到“描述一种作用”，这是全系列唯一的观念跃迁。所有技术细节都是为它服务的。

\subsection*{暗线三：$\delta$ 的身份升级}

观众对 $\delta$ 的理解，每一集都要更新一次。注意这个升级是\textbf{单向}的：后面的理解包含前面的，但前面的理解不足以支撑后面的。

\begin{center}
\small
\renewcommand{\arraystretch}{1.4}
\begin{tabularx}{\textwidth}{@{}L{1.3cm}L{4.2cm}L{9.4cm}@{}}
\toprule
集数 & 此刻 $\delta$ 被当成 & 这一集的观众应该能做的判断 \\
\midrule
EP01 & 一族尖峰的极限 & “它不是某点很高，而是权重堆在一点。”\\
EP02 & 可以求导的对象 & “$\delta'$ 不是另一个尖峰，它有两叶、总积分为零，探测的是斜率。”\\
EP03 & 公式里的一个加项 & “零频率那一项是可以被 $c\delta$ 平移的，所以要额外定系数。”\\
EP05 & 一个连续线性泛函 & “$\delta$ 不是函数，但它有严格定义；它也\textbf{不能}和另一个分布随便相乘。”\\
EP06 & $\Sv'$ 的一个元素 & “它有傅里叶变换，而且那个变换就是常数 $1$。”\\
\bottomrule
\end{tabularx}
\end{center}

\begin{warnbox}[这条暗线的反面教材]
如果 EP01 就把 $\delta$ 定义成“原点无穷高、其余为零”的东西，后面五集全部崩盘：EP02 的 $\delta'$ 画不出来，EP03 的系数说不清，EP05 的分布定义变成多余，EP06 的对偶公式变成天书。

\textbf{EP01 必须把 $\delta$ 说成“一个极限的名字”，而不是“一个形状的名字”。}这是整个系列最重要的一次讲解设计。
\end{warnbox}

\section{欠账与还账：哪一集建立，哪一集证明}

这是全系列的结构性承诺，也是它区别于“六个独立科普视频”的地方：前三集\textbf{大量借用}规则，后三集\textbf{逐条追认}。

\begin{center}
\small
\renewcommand{\arraystretch}{1.5}
\begin{tabularx}{\textwidth}{@{}L{4.6cm}L{1.8cm}L{2.0cm}Y@{}}
\toprule
结论 & 何时出现 & 何时证明 & 证明依赖什么 \\
\midrule
$\F[1]=2\pi\delta$ & EP01 & EP06 & 对偶定义 + 反演公式（在 $\varphi(0)$ 处取值）\\
取样极限 $\int K_{\ep}\varphi\to\varphi(0)$ & EP01（借用） & EP04 & 换元 + 控制收敛；或“近处/远处”估计 \\
$\ip{\delta'}{\varphi}=-\varphi'(0)$ & EP02（借用） & EP05 & 分部积分把导数转移到探针 \\
$\F[t^{n}]=2\pi i^{n}\delta^{(n)}$ & EP02 & EP06 & 乘 $t$ 规则在 $\Sv'$ 上成立 \\
$H'=\delta$ & EP03 & EP05 & 分布导数的定义 \\
$\F[H]=\pi\delta+\PV(1/i\omega)$ & EP03 & EP06 & $H_{\ep}\to H$ 于 $\Sv'$ + 连续性 \\
$\ep/(\ep^{2}+\omega^{2})\to\pi\delta$ & EP03 & EP05、EP06 & 即 EP01 的核极限 \\
$\omega/(\ep^{2}+\omega^{2})\to\PV(1/\omega)$ & EP03 & EP06 & 控制收敛 + 中值估计 \\
\bottomrule
\end{tabularx}
\end{center}

\begin{ideabox}[讲解上的用法]
这张表可以直接变成叙事节奏。每当前三集借了一条规则，就明说一句：

\begin{center}
\itshape
“这条规则我们现在先用着，它的合法性要到后面某一集才能说清楚。”
\end{center}

这句话有两个作用：一是诚实，二是\textbf{制造欠账感}。观众会带着“到底欠了什么”继续看下去。等 EP05、EP06 把账还上时，那种“原来前面那些奇怪的记号真的有出处”的快感，是整个系列最大的回报。
\end{ideabox}

\section{统一约定：一套记号，一条命}

全系列（以及这本书）只用一套约定：

\begin{equation}
\fh(\omega)=\F[f](\omega)=\int_{\R}f(t)e^{-i\omega t}\dd t,
\qquad
f(t)=\frac{1}{2\pi}\int_{\R}\fh(\omega)e^{i\omega t}\dd\omega. \label{eq:conv}
\end{equation}

\begin{itemize}[leftmargin=1.6em,itemsep=0.3em]
\item $\omega$ 是\textbf{角频率}，单位是弧度/秒。时间轴一律写 $t$，频率轴一律写 $\omega$。
\item $2\pi$ 全部放在\textbf{逆变换}上。正变换前面没有系数。
\item 第二个等式（反演）需要条件；在 $\Sv$ 上它总成立，这是 EP06 的事。
\end{itemize}

\subsection*{三种常见约定的对照}

这是最容易出事的地方：不同教材的 $\delta$ 系数完全不一样，抄错一个 $2\pi$，整集公式全废。

\begin{center}
\small
\renewcommand{\arraystretch}{1.55}
\begin{tabularx}{\textwidth}{@{}L{2.6cm}L{4.3cm}L{3.4cm}L{3.6cm}@{}}
\toprule
约定 & 正变换 & $\F[1]$ & 常见于 \\
\midrule
角频率（\textbf{本书}） & $\displaystyle\int f(t)e^{-i\omega t}\dd t$ & $2\pi\delta(\omega)$ & 数学、物理 \\
对称型 & $\displaystyle\frac{1}{\sqrt{2\pi}}\int f(t)e^{-i\omega t}\dd t$ & $\sqrt{2\pi}\,\delta(\omega)$ & 泛函分析、量子力学 \\
频率型 & $\displaystyle\int f(t)e^{-2\pi i\nu t}\dd t$ & $\delta(\nu)$ & 信号处理、工程 \\
\bottomrule
\end{tabularx}
\end{center}

换算关系（用来快速自检）：设 $\F_\nu$ 是频率型、$\F_\omega$ 是本书的角频率型，则

\begin{equation}
\F_\nu[f](\nu)=\F_\omega[f](2\pi\nu),
\qquad
\delta(2\pi\nu)=\frac{1}{2\pi}\delta(\nu). \label{eq:conversion}
\end{equation}

验证：$\F_\nu[1](\nu)=\F_\omega[1](2\pi\nu)=2\pi\delta(2\pi\nu)=2\pi\cdot\frac{1}{2\pi}\delta(\nu)=\delta(\nu)$。两边对上了。

\begin{warnbox}[一句话自检法]
如果你换了一套约定，先算 $\F[1]$。它等于 $2\pi\delta$、$\sqrt{2\pi}\delta$ 还是 $\delta$，直接告诉你 $2\pi$ 被放在了哪里。这个方法在查别人给的公式时，比逐项核对快得多。
\end{warnbox}

\subsection*{五条硬规则}

这五条是全书通用的“物理学定律”，违反任何一条都会在某一集翻车。

\begin{enumerate}[leftmargin=2em,itemsep=0.45em]
\item \textbf{$\delta$ 是分布，不是函数。}它没有“在某点的取值”。任何“$\delta(0)=\infty$”的说法都要避免。
\item \textbf{箭头表示权重，不是高度。}$\F[1]=2\pi\delta$ 里的箭头，说的是“零频率处的权重是 $2\pi$”，不是说频谱在那里高到 $2\pi$。
\item \textbf{本例频谱曲线的面积不是能量。}EP01 里 $\int G_{\ep}=2\pi$ 是这条特定曲线的性质，不是“所有频谱都守恒”。真正守恒的是 $\|f\|_2^2=\frac{1}{2\pi}\|\fh\|_2^2$（Plancherel），而常数连 $\Ltwo$ 都不属于。
\item \textbf{出现 $\delta$、$\delta^{(n)}$、$\PV$ 时，必须声明“在广义函数意义下”。}这不是免责声明，这是数学内容的一部分。
\item \textbf{$2\pi$、$i$、负号三者联立。}改一个就要改全部。EP03 里 $\F[H]$ 的 $\pi$ 与 $\frac{1}{i}=-i$ 尤其容易错。
\end{enumerate}

\section{四个空间：一个对象该住在哪儿}

同一个“信号”，放在不同的空间里，能问的问题完全不同。这是整个系列的技术骨架。

\begin{center}
\small
\renewcommand{\arraystretch}{1.5}
\begin{tabularx}{\textwidth}{@{}L{1.7cm}L{4.4cm}L{4.2cm}Y@{}}
\toprule
空间 & 里的元素 & 傅里叶变换 & 关键性质 \\
\midrule
$\Lone$ & $\int|f|<\infty$ & 一定有，且 $\|\fh\|_\infty\le\|f\|_1$ & 积分绝对收敛；$\fh$ 连续且趋于零 \\
$\Ltwo$ & $\int|f|^{2}<\infty$ & 有（Plancherel 延拓） & 保能量：$\|f\|_2^{2}=\frac{1}{2\pi}\|\fh\|_2^{2}$ \\
$\D'$ & 分布（EP05） & \textbf{不一定有} & 可以求导、乘光滑函数；适合局部观察 \\
$\Sv'$ & 温和分布（EP06） & \textbf{一定有}，且连续可逆 & 本系列的工作空间 \\
\bottomrule
\end{tabularx}
\end{center}

\subsection*{把六集里出现的对象，一个个安置好}

\begin{center}
\small
\renewcommand{\arraystretch}{1.4}
\begin{tabularx}{\textwidth}{@{}L{3.4cm}L{1.5cm}L{1.5cm}L{1.5cm}L{1.5cm}L{4.3cm}@{}}
\toprule
对象 & $\Lone$ & $\Ltwo$ & $\D'$ & $\Sv'$ & 说明 \\
\midrule
$e^{-\ep|t|}$ & \checkmark & \checkmark & \checkmark & \checkmark & 最乖的对象，EP01 全靠它 \\
$1$ & $\times$ & $\times$ & \checkmark & \checkmark & 不衰减，但增长有界 \\
$t^{n}$ & $\times$ & $\times$ & \checkmark & \checkmark & 多项式增长 \\
$H(t)$ & $\times$ & $\times$ & \checkmark & \checkmark & 有界，但两端不衰减 \\
$1/(1+t^{2})$ & \checkmark & \checkmark & \checkmark & $\times$ & 可积但只按固定幂次衰减 \\
$\delta,\ \delta^{(n)}$ & $\times$ & $\times$ & \checkmark & \checkmark & 不是函数，是泛函 \\
$e^{t^{2}}$ & $\times$ & $\times$ & \checkmark & $\times$ & 有分布，但\textbf{没有}温和分布 \\
\bottomrule
\end{tabularx}
\end{center}

\begin{ideabox}[这张表最值得记住的两行]
\textbf{一行是 $1$：}它不在 $\Lone$，也不在 $\Ltwo$。这解释了为什么 EP01 的两条经典路线（普通积分、Plancherel）都救不了它。

\textbf{一行是 $e^{t^{2}}$：}它虽然“是函数”，却没有傅里叶变换（在这个框架内）。这提醒观众：能不能做变换，不是看“它像不像函数”，而是看它住在哪个空间。
\end{ideabox}

\begin{notebox}[换空间不是作弊]
初学者常有一个误解：觉得“先加衰减再取极限”是在绕开困难，甚至是在耍花招。

不是的。真实的逻辑是：\textbf{原来的工具只在一个小房间里够用}（$\Lone$ 或 $\Ltwo$），而你要处理的信号在房间外面。你有两条路：要么放弃，要么把房间扩大。数学选择了后者，而且扩大的方式非常克制——它不改变任何原有结论，只增加新的对象。

这一点在 EP06 结尾必须说清楚，否则观众会觉得整套理论是“为了自圆其说而发明出来的”。
\end{notebox}

\section{三张许可证：什么时候可以把极限搬进积分}

整个系列里，凡是出现“把极限、导数、求和搬进积分号”的地方，都必须掏出下面三张许可证之一。没有许可证的动作，就是错的。

\begin{center}
\small
\renewcommand{\arraystretch}{1.5}
\begin{tabularx}{\textwidth}{@{}L{2.2cm}L{5.8cm}Y@{}}
\toprule
许可证 & 内容 & 什么时候失效 \\
\midrule
控制收敛 & 若 $f_n\to f$，且存在\textbf{同一个}可积 $g$ 使 $|f_n|\le g$，则 $\int f_n\to\int f$ & 找不到公共控制函数时。EP04 会证明：对 $K_{\ep}$ 根本不存在这样的 $g$ \\
Fubini & 双重积分绝对收敛时，可以交换积分次序 & 只有条件收敛时。EP06 的反演证明专门用高斯因子制造绝对收敛 \\
分部积分 & $\int f'\varphi=-\int f\varphi'$，边界项必须为零 & 边界项不消失时。$f$ 有跳跃（EP03）、$f$ 不衰减（EP02 的分布版本）都要重新处理 \\
\bottomrule
\end{tabularx}
\end{center}

\subsection*{全书的“许可证使用记录”}

写脚本时如果记不清某个交换是否合法，查这张表。

\begin{center}
\small
\renewcommand{\arraystretch}{1.45}
\begin{tabularx}{\textwidth}{@{}L{5.4cm}L{3.6cm}Y@{}}
\toprule
正在做的事 & 需要哪张 & 条件 \\
\midrule
对 $\omega$ 求导搬进积分 & 控制收敛（局部） & $f\in\Sv$，故 $|tf|$ 可积 \\
$\ip{K_{\ep}}{\varphi}\to\varphi(0)$ & 控制收敛（换元后） & $\varphi$ 有界、在 $0$ 连续 \\
对偶公式的推导 & Fubini & $f,\varphi\in\Sv$（或都在 $\Lone$） \\
反演公式 & Fubini $+$ 近似恒等核 & $\varphi\in\Sv$ \\
$H'=\delta$ & 分部积分 & $\varphi\in\D$，边界项为零 \\
$\PV\frac{1}{\omega}$ 的存在性 & 中值估计 $+$ 尾部衰减 & $\varphi\in\Sv$ \\
$\F[H]$ 的两项极限 & 控制收敛 $+$ 上述全部 & $\varphi\in\Sv$ \\
\bottomrule
\end{tabularx}
\end{center}

\begin{warnbox}[全系列最大的一处陷阱]
\textbf{不要用还没证明的 $\delta$ 公式，去证明反演。}

这条循环论证非常诱人：既然 $\int e^{i\omega t}\dd\omega=2\pi\delta(t)$，那反演公式不就一眼看出来了？

问题在于：那个“等式”本身就是在分布意义下、并且依赖对偶定义才能写出来的。用它证反演，等于用结论证前提。EP06 里反演的证明必须走高斯正则化这条“干净”的路——先造一个绝对收敛的双重积分，再用近似恒等核。这条纪律必须守住。
\end{warnbox}

\section{六集速查卡}

把这一页打印出来贴在显示器边上，写脚本时随时对照。

\begin{center}
\small
\renewcommand{\arraystretch}{1.5}
\begin{tabularx}{\textwidth}{@{}L{1.0cm}L{1.6cm}L{3.9cm}Y L{3.8cm}@{}}
\toprule
集 & 时长 & 主公式 & 必须出现的画面 & 结尾留下的问题 \\
\midrule
01 & 8--10 & $\F[1]=2\pi\delta$ & 时频双面板，$\ep$ 同步变化 & 乘以 $t$ 会怎样？\\
02 & 10--12 & $\F[t^{n}]=2\pi i^{n}\delta^{(n)}$ & 正负两叶的 $\delta'$ 近似 & 求导对应什么？\\
03 & 10--14 & $\F[H]=\pi\delta+\PV\frac{1}{i\omega}$ & 实部成峰、虚部成奇 & 这些极限是什么？\\
04 & 12--15 & $\int K_{\ep}\varphi\to\varphi(0)$ & 三种核同时收缩 & 能否用作用定义对象？\\
05 & 12--15 & $\ip{\delta}{\varphi}=\varphi(0)$ & 探针进出、数值输出 & 怎么对分布做变换？\\
06 & 12--18 & $\ip{\F T}{\varphi}=\ip{T}{\F\varphi}$ & 三条公式并列回收 & 回到那条水平线 \\
\bottomrule
\end{tabularx}
\end{center}

\begin{quizbox}[读完这一章的自测]
不看上面任何一页，回答：
\begin{enumerate}[leftmargin=1.6em,itemsep=0.35em]
\item 为什么常数 $1$ 既不在 $\Lone$ 也不在 $\Ltwo$？这分别挡住了哪两条路？
\item 三条主公式各自暴露的是哪一类“失败方式”？
\item 哪一集是“从算结果转向问意义”的转折点？为什么它必须放在第三集之后？
\item 三条暗线分别是什么？各自在哪一集“往前挪一格”？
\item 如果有人在 EP01 就把 $\delta$ 说成“原点处无穷高的函数”，后面哪几集会出问题？各出什么问题？
\item 为什么不能用 $\int e^{i\omega t}\dd\omega=2\pi\delta(t)$ 去证明反演公式？
\end{enumerate}
\end{quizbox}
}{}
\IfFileExists{ep01.tex}{\chapter{EP01｜常数为什么对应 $\delta$？}

\begin{ideabox}[这一集的一句话]
\textbf{信号在时间上铺得越开，频谱就收得越紧；常数把这种收紧推到极限，于是剩下的只有零频率上的一份权重。}
\end{ideabox}

\section{观众的问题链}

这一集不要求观众懂分布，甚至不要求他见过 $\delta$。它只要求一件事：愿意相信“常数只包含零频率”这个直觉，然后亲眼看着这个直觉在积分面前失效。

问题要一层一层递进，每一层都由上一层逼出来：

\begin{center}
\small
\renewcommand{\arraystretch}{1.32}
\begin{tabularx}{\textwidth}{@{}L{0.8cm}Y L{4.4cm}@{}}
\toprule
层 & 问题 & 观众此刻的心理 \\
\midrule
1 & 一个完全不变的信号，包含频率吗？ & 当然，只有零频率。这还要问？\\
2 & 那就把它代进傅里叶积分算一算。 & 应该能算出来吧……\\
3 & 算不出来。振荡得再厉害也没有抵消。 & 那是不是我算错了？\\
4 & 换个办法：先让它衰减，再让衰减消失。 & 这算不算耍赖？\\
5 & 衰减越弱，频谱越窄越高，可是总面积不变。 & 它没有消失，只是挤到一点上去了。\\
\bottomrule
\end{tabularx}
\end{center}

\begin{center}
\small
\renewcommand{\arraystretch}{1.32}
\begin{tabularx}{\textwidth}{@{}L{3.0cm}Y@{}}
\toprule
观众带着什么进来 & 观众带着什么出去 \\
\midrule
“傅里叶变换 = 看看有哪些频率” & “不衰减的信号，普通积分一定失效；但失效的是积分，不是频率的想法。”\\
“常数就是直流，直流就是零频” & “常数的频谱全部压在 $\omega=0$，并且可以用一个叫 $\delta$ 的对象写出等式。”\\
把 $\delta$ 想成一根“很高的线” & “$\delta$ 描述的不是高度，是权重；它的意义在于它在积分里干什么。”\\
\bottomrule
\end{tabularx}
\end{center}

\section{工具：频率探针，以及它为什么要求衰减}

\subsection{为什么用 $e^{-i\omega t}$ 去乘}

傅里叶变换的核 $e^{-i\omega t}$ 是一个“频率探针”。把它乘到信号上：

\begin{equation}
f(t)\cdot e^{-i\omega t}.
\end{equation}

如果 $f$ 本身是频率为 $\omega_0$ 的纯振荡 $e^{i\omega_0 t}$，乘积就是

\begin{equation}
e^{i\omega_0 t}\cdot e^{-i\omega t}=e^{i(\omega_0-\omega)t}.
\end{equation}

\begin{itemize}[leftmargin=1.6em,itemsep=0.3em]
\item 当 $\omega=\omega_0$ 时，乘积等于 $1$：始终同相，积分不断累加，所以结果很大。
\item 当 $\omega\neq\omega_0$ 时，乘积是一个不停止的旋转，正负两半互相抵消，所以结果很小。
\end{itemize}

这就是“频率分析”的全部直觉：\textbf{用探针去和信号比对，对上了就积累，对不上就抵消}。

\begin{ideabox}[但“抵消”不是一个证明]
初学者最容易在这里滑过去的一句话是：“不匹配的时候，正负抵消，所以积分是零。”

这句话在\textbf{有限窗口}里是对的（而且只是近似），在\textbf{无限区间}里没有意义。原因很朴素：$|e^{-i\omega t}|=1$，被积函数从来没有变小过。“抵消”描述的是振荡的相位在转，而不是数值在衰减。要让它成为一个定理，必须另外给条件。

整个 EP01 的第一处张力就在这里，整个系列的第一处张力也在这里。
\end{ideabox}

\subsection{唯一的充分条件：绝对可积}

\begin{prop}[$\Lone$ 上的傅里叶变换]\label{prop:L1}
若 $f\in\Lone$，则对每个 $\omega$，积分 $\int_{\R}f(t)e^{-i\omega t}\dd t$ 绝对收敛，并且
\begin{equation}
|\fh(\omega)|\le \|f\|_1=\int_{\R}|f(t)|\dd t .
\end{equation}
此外 $\fh$ 在 $\R$ 上一致连续，且 $\fh(\omega)\to0$（$|\omega|\to\infty$）。
\end{prop}

证明只需一句：$|f(t)e^{-i\omega t}|=|f(t)|$，与 $\omega$ 无关，所以积分被一个与 $\omega$ 无关的可积函数控制。连续性与“趋于零”分别是控制收敛与 Riemann--Lebesgue 引理的推论。

\begin{notebox}[这两条性质在画面上的意义]
\begin{itemize}[leftmargin=1.4em,itemsep=0.2em]
\item “$|\fh|\le\|f\|_1$”：频谱有界。所以 EP01 后面出现的峰高 $2/\ep$ 越界增长，恰恰说明 $f_\ep$ 的 $\Lone$ 范数 $2/\ep$ 也在爆炸。两者是同一件事。
\item “$\fh\to0$ 于无穷远”：任何 $\Lone$ 信号的频谱在远处必然衰减到零。所以一条“处处不为零、远处还很亮”的频谱，注定不能来自 $\Lone$ 信号。
\end{itemize}
\end{notebox}

\subsection{常数不在 $\Lone$ 里，但这只说明定理用不上}

$\int_{\R}1\,\dd t=\infty$，所以常数 $f\equiv1$ 不属于 $\Lone$，命题 \ref{prop:L1} 用不上。

\begin{warnbox}[一个非常容易被滥用的推理]
“常数不在 $\Lone$，所以它的傅里叶变换不存在。”

这是\textbf{错的}推理。不在 $\Lone$ 只说明这条定理管不到它，不说明积分一定发散。要断言发散，必须\textbf{具体去算}。

这一集的价值，有一半在于让观众看见“具体去算”这四个字：不是引用定理，而是把 $\int_{-T}^{T}e^{-i\omega t}\dd t$ 真的积出来，然后看着它不收敛。

创作者要能分清这两句话：
\begin{itemize}[leftmargin=1.4em,itemsep=0.2em]
\item 对：常数不在 $\Lone$ 里，所以不能直接用绝对收敛的定理。
\item 错：常数不在 $\Lone$ 里，所以它的积分发散。
\end{itemize}
\end{warnbox}

\subsection{顺带一句：直流为什么是零频率}

$1=e^{i\cdot0\cdot t}$。零频率的意思是“相位不转”，也就是“不随时间变化”。所以“常数”和“零频率”在数学上是同一句话的两种说法。

这句话要在开场说清楚，因为后面所有的“集中在原点”都建立在这个等式上。

\section{直接积分：把失败算给人看}

\subsection{对称截断}

形式上把 $f(t)=1$ 代进定义：

\begin{equation}
\int_{\R}e^{-i\omega t}\dd t .
\end{equation}

为了给它一个机会，先做对称截断（这也是最“讲道理”的截断方式）：

\begin{equation}
I_T(\omega)=\int_{-T}^{T}e^{-i\omega t}\dd t
=\begin{cases}
\dfrac{2\sin(\omega T)}{\omega}, & \omega\neq0,\\[9pt]
2T, & \omega=0 .
\end{cases}\label{eq:IT}
\end{equation}

计算只有一行，但要讲清楚两件事：

\begin{itemize}[leftmargin=1.6em,itemsep=0.35em]
\item $\omega\neq0$ 时：$\left[\dfrac{e^{-i\omega t}}{-i\omega}\right]_{-T}^{T}=\dfrac{e^{-i\omega T}-e^{i\omega T}}{-i\omega}=\dfrac{2\sin(\omega T)}{\omega}$。
\item $\omega=0$ 时：被积函数是 $1$，长度是 $2T$，结果是 $2T$。
\item 两式在 $\omega=0$ 处衔接得上：$\lim_{\omega\to0}\frac{2\sin(\omega T)}{\omega}=2T$。所以这不是一个“有洞”的函数，它在 $\omega=0$ 处是连续的。
\end{itemize}

\begin{ideabox}[顺带记住这个函数]
$I_T(\omega)=\dfrac{2\sin(\omega T)}{\omega}$ 不是凭空出现的配角。它正是\textbf{矩形窗} $w_T=\mathbf{1}_{[-T,T]}$ 的傅里叶变换：
\begin{equation}
\widehat{w_T}(\omega)=\int_{-T}^{T}e^{-i\omega t}\dd t=I_T(\omega).
\end{equation}
这一点在 EP01 里只提一句，但它在 EP06 会变成一个漂亮的闭环：$w_T\to1$（分布意义），而 $\widehat{w_T}\to2\pi\delta$。所以第一集看似“失败”的那条曲线，其实已经在暗中指向答案。\textbf{不要在这一集把它讲成纯粹的废品。}
\end{ideabox}

\subsection{固定 $\omega\neq0$：振幅不变，永不收敛}

固定一个非零的 $\omega$，让 $T\to\infty$。此时

\begin{equation}
I_T(\omega)=\frac{2\sin(\omega T)}{\omega},
\end{equation}

而 $\sin(\omega T)$ 在 $-1$ 与 $1$ 之间永不停歇。振幅 $\dfrac{2}{|\omega|}$ 是一个与 $T$ 无关的常数。所以：

\begin{center}
\textbf{它不是“趋于零”，也不是“趋于无穷”，它是根本没有极限。}
\end{center}

想给观众一个无可争辩的证据，就取两列 $T$：

\begin{equation}
T_n=\frac{\pi/2+2\pi n}{|\omega|},
\qquad
S_n=\frac{3\pi/2+2\pi n}{|\omega|},
\qquad n=0,1,2,\dots
\end{equation}

于是

\begin{equation}
I_{T_n}(\omega)=\frac{2}{|\omega|},
\qquad
I_{S_n}(\omega)=-\frac{2}{|\omega|}.
\end{equation}

两列都趋于无穷，但结果分别停在两个不同的数上。极限不存在。

\begin{scriptbox}[动画要点 01-B]
左图：$I_T(\omega_*)$ 随 $T$ 变化的曲线，取定 $\omega_*=1$，标记点沿 $2\sin(T)$ 上下摆动，\textbf{永不收敛}。

右图：同一时刻的 $I_T(0)=2T$，箭头不断上升。

一句话旁白：\textbf{“非零频率处它一直在摆，零频率处它一直往上冲。这两种都不是‘收敛’。”}
\end{scriptbox}

\subsection{零频率处：不是“很大”，而是“没有上界”}

$\omega=0$ 时 $I_T(0)=2T\to\infty$。注意这里的措辞：不是“频谱在零频率处很高”，而是“随着截断变长，这个量无上界地增长”。

这个区别很重要。正因为它是“无上界地增长”，后面我们才需要一个\textbf{权重}的概念，而不是一个\textbf{高度}的概念。

\subsection{结论怎么写才准确}

\begin{warnbox}[三种说法，只有一种是对的]
\begin{itemize}[leftmargin=1.4em,itemsep=0.3em]
\item \textbf{错}：“非零频率处积分等于零。”——积分不存在，谈不上等于零。
\item \textbf{错}：“零频率处是无穷大，所以频谱是无穷大。”——$\infty$ 不是一个数，不能当频谱值。
\item \textbf{对}：“对固定 $\omega\neq0$，截断积分随 $T$ 持续振荡，没有极限；在 $\omega=0$ 处无上界增长。因此原式不能作为通常的反常积分求值。”
\end{itemize}
\end{warnbox}

\begin{notebox}[一句可以放在旁白里的比喻]
“振荡会抵消”就像说“两支方向相反的队伍拔河，所以绳子不动”。问题是：绳子上根本没有裁判喊停。只要两边一直用力，位置就一直在晃。\textbf{振荡是一种运动，不是一种收敛。}
\end{notebox}

\section{正则化：不是绕开，而是逼近}

\subsection{换一个问题的提法}

直接算 $1$ 失败了。但我们真正想要的不是“算出这个积分”，而是“说清楚常数在频率上的样子”。

于是换一个提法：\textbf{造一列能算的信号，每个都越来越像常数，看它们的频谱收敛到哪里。}

这在数学上叫\textbf{正则化}，在讲故事上叫“先给它一个尾巴，再让尾巴慢慢松开”。

\begin{ideabox}[为什么这不是耍赖]
观众（以及很多创作者自己）会有这个疑问：既然算不出来，为什么换个算得出来的东西算，答案就能算数？

关键区别在于：我们\textbf{没有}声称“$\int e^{-i\omega t}\dd t$ 这个积分等于某个数”。我们声称的是另一件事：\textbf{有一列函数 $f_\ep$，它们逐点趋于 $1$，它们的频谱在一种明确的收敛意义下趋于 $2\pi\delta$，所以我们把 $2\pi\delta$ 定义为 $1$ 的频谱。}

这个“明确的收敛意义”正是 EP04 要做的事。第一集只需要把这件事说成“一种新的、待定义的意义”，并诚实地说“严格定义留到后面”。
\end{ideabox}

\subsection{选哪个外壳}

取

\begin{equation}
f_\ep(t)=e^{-\ep|t|},\qquad \ep>0 .\label{eq:feps}
\end{equation}

三个理由，按重要性排列：

\begin{enumerate}[leftmargin=2em,itemsep=0.4em]
\item \textbf{它绝对可积。}$\|f_\ep\|_1=\int_{\R}e^{-\ep|t|}\dd t=\dfrac{2}{\ep}<\infty$。所以命题 \ref{prop:L1} 适用，积分真的收敛，可以放心算。
\item \textbf{它逐点趋于常数。}对每个固定的 $t$，$e^{-\ep|t|}\to1$（$\ep\to0^+$）。
\item \textbf{它是偶函数且非负。}偶性保证频谱是实偶函数（好画、好检查），非负性保证频谱非负（后面“面积集中”的图像才成立）。
\end{enumerate}

\begin{notebox}[为什么不用单边衰减 $e^{-\ep t}H(t)$]
因为那是“从零跳起来再衰减”，它的极限是阶跃 $H$，不是常数 $1$。这正好是 EP03 的主角。这里要的是\textbf{两边都衰减}，极限才是那条水平线。

顺便记住这个对照：EP01 用双边、EP03 用单边。两个外壳，两个不同的极限对象。
\end{notebox}

\begin{warnbox}[$\ep$ 越小越像常数，但永远不等于常数]
动画结束时千万别说“现在它等于 $1$ 了”。$\ep=0.1$ 时，$f_\ep(10)=e^{-1}\approx0.37$，还差得远。它只是在\textbf{越来越长的时间范围内}接近 $1$。这也是为什么“逐点收敛”必须限定“对每个固定的 $t$”：$t$ 越大，需要的 $\ep$ 越小。
\end{warnbox}

\section{唯一需要跟到底的计算}

这是全片唯一一段要求观众跟着走的代数。它必须拆成三到四步显现，不能一屏出完。

\subsection{第一步：为什么必须拆成两半}

$|t|$ 在 $t=0$ 处换了一张脸：左边是 $-t$，右边是 $t$。所以被积函数在左右半轴是不同的表达式，积分必须分开算。

\begin{equation}
\fh_\ep(\omega)=\underbrace{\int_{-\infty}^{0}e^{\ep t}e^{-i\omega t}\dd t}_{\text{左半轴：}|t|=-t}
+\underbrace{\int_{0}^{\infty}e^{-\ep t}e^{-i\omega t}\dd t}_{\text{右半轴：}|t|=t}.
\end{equation}

\begin{scriptbox}[动画要点 01-D]
左半轴用一种颜色，右半轴用另一种颜色；两种颜色分别“飞向”下面两项分式。合并时两种颜色合成一项。

观众只需要记住一句：\textbf{分母里那个 $\ep^{2}+\omega^{2}$，是左右两半加起来的结果。}
\end{scriptbox}

\subsection{第二步：两个半轴各自算}

\begin{align}
\int_{-\infty}^{0}e^{(\ep-i\omega)t}\dd t
&=\left[\frac{e^{(\ep-i\omega)t}}{\ep-i\omega}\right]_{-\infty}^{0}
=\frac{1}{\ep-i\omega},\\
\int_{0}^{\infty}e^{-(\ep+i\omega)t}\dd t
&=\left[\frac{e^{-(\ep+i\omega)t}}{-(\ep+i\omega)}\right]_{0}^{\infty}
=\frac{1}{\ep+i\omega}.
\end{align}

两处的边界项都要交代清楚，这是“为什么 $\ep>0$ 不能省”的地方：

\begin{itemize}[leftmargin=1.6em,itemsep=0.3em]
\item 左端 $t\to-\infty$：$|e^{(\ep-i\omega)t}|=e^{\ep t}\to0$，因为 $\ep>0$。
\item 右端 $t\to+\infty$：$|e^{-(\ep+i\omega)t}|=e^{-\ep t}\to0$，同样因为 $\ep>0$。
\end{itemize}

如果 $\ep=0$，两侧的模都恒等于 $1$，边界项根本不消失——这正是上一节失败的原因。\textbf{$\ep>0$ 不是技术性假设，它是整个计算能成立的原因。}

\subsection{第三步：通分合并}

\begin{equation}
\fh_\ep(\omega)=\frac{1}{\ep-i\omega}+\frac{1}{\ep+i\omega}
=\frac{(\ep+i\omega)+(\ep-i\omega)}{(\ep-i\omega)(\ep+i\omega)}
=\frac{2\ep}{\ep^{2}+\omega^{2}}.
\end{equation}

记

\begin{equation}
G_\ep(\omega):=\frac{2\ep}{\ep^{2}+\omega^{2}} .\label{eq:Geps}
\end{equation}

\subsection{第四步：检查结果是否说得出道理}

\begin{center}
\small
\renewcommand{\arraystretch}{1.3}
\begin{tabularx}{\textwidth}{@{}L{5.0cm}Y@{}}
\toprule
检查项 & 为什么应当如此 \\
\midrule
$G_\ep$ 是实函数 & $f_\ep$ 是实函数，所以 $\overline{\fh(\omega)}=\fh(-\omega)$；实偶信号的频谱是实偶函数 \\
$G_\ep$ 是偶函数 & $f_\ep$ 是偶函数 \\
$G_\ep\ge0$ & $f_\ep\ge0$ 且偶，频谱可以取到非负值（不是定理，但本例成立） \\
$G_\ep(0)=2/\ep$ & 直接代入；$\ep\to0$ 时爆炸，与“零频率处无上界”呼应 \\
\bottomrule
\end{tabularx}
\end{center}

\begin{ideabox}[一个能让画面变漂亮的视角（可选，不必讲）]
若把 $\omega$ 看成复变量，$G_\ep(\omega)=\dfrac{2\ep}{\ep^{2}+\omega^{2}}$ 有两个极点，位置在 $\omega=\pm i\ep$。

$\ep$ 越小，这两个极点越贴近实轴；当 $\ep\to0$，它们撞上实轴的原点——那正是 $\delta$ “出生”的位置。

这个图像对懂复变的观众非常有效，但\textbf{不要放在主视频里}：它会引入一条全新的解释线，和“面积集中”那条线抢戏。留作补充材料。
\end{ideabox}

\section{三个指标：变高、变窄、面积不变}

只算出一个公式是不够的。要让观众“看见”集中，必须给他三个可以量的数。

\begin{enumerate}[leftmargin=2em,itemsep=0.5em]
\item \textbf{峰高}：$G_\ep(0)=\dfrac{2}{\ep}$。$\ep$ 减半，峰高翻倍。
\item \textbf{半高全宽（FWHM）}：解 $\dfrac{2\ep}{\ep^{2}+\omega^{2}}=\dfrac{1}{2}\cdot\dfrac{2}{\ep}$，即 $\ep^{2}+\omega^{2}=2\ep^{2}$，得 $|\omega|=\ep$。所以宽度是 $2\ep$。$\ep$ 减半，宽度减半。
\item \textbf{总面积}：
\begin{equation}
\int_{\R}G_\ep(\omega)\dd\omega
=2\ep\int_{\R}\frac{\dd\omega}{\ep^{2}+\omega^{2}}
=2\ep\cdot\frac{1}{\ep}\left[\arctan\frac{\omega}{\ep}\right]_{-\infty}^{\infty}
=2\pi .
\end{equation}
\end{enumerate}

\begin{center}
\small
\renewcommand{\arraystretch}{1.35}
\begin{tabular}{@{}cccc@{}}
\toprule
$\ep$ & 峰高 $2/\ep$ & 半高全宽 $2\ep$ & 全实轴面积 \\
\midrule
$1$ & $2$ & $2$ & $2\pi$ \\
$0.5$ & $4$ & $1$ & $2\pi$ \\
$0.2$ & $10$ & $0.4$ & $2\pi$ \\
$0.1$ & $20$ & $0.2$ & $2\pi$ \\
$0.02$ & $100$ & $0.04$ & $2\pi$ \\
\bottomrule
\end{tabular}
\end{center}

\begin{figure}[htbp]
\centering
\begin{tikzpicture}
\begin{axis}[width=0.46\textwidth,height=4.3cm,
  domain=-4:4,samples=200,ymin=0,ymax=1.06,
  xlabel={$t$},ylabel={$f_\ep(t)$},
  title={\small 时域：$\ep$ 越小越平},
  legend style={font=\scriptsize,at={(0.03,0.97)},anchor=north west,draw=none},
  ticklabel style={font=\scriptsize},label style={font=\small}]
\addplot[thick,mainblue] {exp(-1*abs(x))};
\addplot[thick,warnred,dashed] {exp(-0.4*abs(x))};
\legend{$\ep=1$,$\ep=0.4$}
\end{axis}
\end{tikzpicture}\hfill
\begin{tikzpicture}
\begin{axis}[width=0.46\textwidth,height=4.3cm,
  domain=-4:4,samples=400,ymin=0,ymax=5.6,
  xlabel={$\omega$},ylabel={$G_\ep(\omega)$},
  title={\small 频域：$\ep$ 越小越尖},
  legend style={font=\scriptsize,at={(0.97,0.97)},anchor=north east,draw=none},
  ticklabel style={font=\scriptsize},label style={font=\small}]
\addplot[thick,mainblue] {2*1/(1+x^2)};
\addplot[thick,warnred,dashed] {2*0.4/(0.16+x^2)};
\legend{$\ep=1$,$\ep=0.4$}
\end{axis}
\end{tikzpicture}
\caption{同一组 $\ep$ 的时频对应。左图变平，右图变高变窄。注意两侧纵轴刻度不同：右图 $\ep=0.4$ 的峰高是 $5$。}
\end{figure}

\subsection{两个尺度是倒数关系}

时域的特征宽度是 $1/\ep$（$e^{-\ep|t|}$ 衰减到 $e^{-1}$ 所需的 $|t|$ 是 $1/\ep$），频域的特征宽度是 $\ep$。两者相乘等于 $1$。

这不是巧合，它来自尺度变换公式：若 $a>0$，

\begin{equation}
f\!\left(\frac{t}{a}\right)
\;\xleftrightarrow{\ \F\ }\;
a\,\fh(a\omega).\label{eq:scaling}
\end{equation}

验证：令 $t=as$，

\begin{equation}
\int_{\R}f\!\left(\frac{t}{a}\right)e^{-i\omega t}\dd t
=a\int_{\R}f(s)e^{-i\omega a s}\dd s
=a\fh(a\omega).
\end{equation}

取 $a=1/\ep$：时域被拉宽 $1/\ep$ 倍，频域就被压窄 $1/\ep$ 倍，同时幅度被压低 $1/\ep$ 倍。再看 $G_\ep$：

\begin{equation}
G_\ep(\omega)=\frac{2\ep}{\ep^{2}+\omega^{2}}
=\frac{1}{\ep}\cdot\frac{2}{1+(\omega/\ep)^{2}}
=\frac{1}{\ep}\,\widehat{g}\!\left(\frac{\omega}{\ep}\right),
\qquad g(t)=e^{-|t|}.
\end{equation}

与 \eqref{eq:scaling} 完全一致。\textbf{这条公式是给创作者的保险：所有“时间铺开、频率集中”的说法，都能用它一句话验证。}

\begin{warnbox}[不要把尺度对偶说成“不确定性原理”]
可以用“时宽 $\times$ 带宽 $=$ 常数”来打比方，但不要说成“测不准原理”，也不要声称“频谱面积守恒”。

真正的 Plancherel 守恒是 $\|f\|_2^2=\frac{1}{2\pi}\|\fh\|_2^2$，而常数连 $\Ltwo$ 都不属于。EP01 里那个“面积恒为 $2\pi$”，是 $G_\ep$ 这一族特定曲线的性质，不是对所有信号都成立的定律。
\end{warnbox}

\section{从“面积集中”到“取样作用”}

\subsection{把面积标准化}

动画里那条面积为 $1$ 的核，和实际频谱差一个 $2\pi$。把这件事说清楚，否则后面每个系数都会错。

\begin{equation}
K_\ep(\omega):=\frac{G_\ep(\omega)}{2\pi}
=\frac{1}{\pi}\frac{\ep}{\ep^{2}+\omega^{2}},
\qquad
\int_{\R}K_\ep(\omega)\dd\omega=1,
\qquad
G_\ep=2\pi K_\ep .\label{eq:Keps}
\end{equation}

记住这个配对：

\begin{center}
\small
\renewcommand{\arraystretch}{1.3}
\begin{tabularx}{\textwidth}{@{}L{5.4cm}L{1.8cm}Y@{}}
\toprule
对象 & 面积 & 用在哪里 \\
\midrule
$K_\ep$（Cauchy 核 / 归一化核） & $1$ & 动画中的“面积为 $1$ 的尖峰”；取样直觉 \\
$G_\ep$（\textbf{本例实际频谱}） & $2\pi$ & 公式里的 $\F[f_\ep]$ \\
\bottomrule
\end{tabularx}
\end{center}

\subsection{集中性：任意小的邻域，最终装下几乎全部质量}

固定任意 $a>0$（注意：$a$ 固定，然后让 $\ep\to0$）：

\begin{equation}
\int_{-a}^{a}K_\ep(\omega)\dd\omega
=\frac{1}{\pi}\left[\arctan\frac{\omega}{\ep}\right]_{-a}^{a}
=\frac{2}{\pi}\arctan\frac{a}{\ep}
\;\longrightarrow\;\frac{2}{\pi}\cdot\frac{\pi}{2}=1 .
\end{equation}

也就是说：\textbf{不管 $a$ 多小，只要 $\ep$ 足够小，原点这个小区间里就装下了几乎全部的权重。}

\begin{ideabox}[这句话要说得非常小心]
正确的说法：“对每个固定的 $a>0$，$\int_{|x|<a}K_\ep\to1$。”

错误的说法：“$K_\ep$ 除原点外处处趋于零，所以整体趋于零。”——前半句是错的（$K_\ep(0)=1/(\pi\ep)\to\infty$），后半句的推理也是错的。

“几乎每一点都趋于零”与“整体趋于零”是两回事。这个区别就是 EP04 的全部内容，这里只需要\textbf{埋下}它，不需要解释它。
\end{ideabox}

\subsection{取样：为什么会“提取原点值”}

取一个光滑、有界的函数 $\varphi$（后面叫它测试函数，这里只说“一条平滑的曲线”）。要算

\begin{equation}
\int_{\R}K_\ep(\omega)\varphi(\omega)\dd\omega .
\end{equation}

直觉分两段：

\begin{itemize}[leftmargin=1.6em,itemsep=0.35em]
\item \textbf{近处}（$|\omega|$ 与 $\ep$ 同量级）：$\varphi$ 变化很小，几乎等于 $\varphi(0)$；而核在近处的质量几乎等于 $1$。于是贡献约为 $1\cdot\varphi(0)$。
\item \textbf{远处}（$|\omega|\gg\ep$）：核的质量已经趋于 $0$，$\varphi$ 有界也翻不起浪。
\end{itemize}

所以

\begin{equation}
\int_{\R}K_\ep(\omega)\varphi(\omega)\dd\omega\;\longrightarrow\;\varphi(0).\label{eq:sample}
\end{equation}

\begin{scriptbox}[动画要点 01-F]
一层阴影从左右两侧向原点收缩，同时一条 $\varphi$ 曲线叠加在上面。阴影覆盖区域内，$\varphi$ 几乎是常数；最终读数显示 $\to\varphi(0)$。

标记要写“面积 $=1$”，但必须注明是\textbf{全实轴}的面积——可见画框里的阴影永远小于 $1$。
\end{scriptbox}

\subsection{用 $\delta$ 的记号把它写下来}

既然“在积分里提取原点值”就是这个对象唯一的作为，那就给它一个名字：

\begin{equation}
\int_{\R}\delta(\omega)\varphi(\omega)\dd\omega=\varphi(0),
\qquad\text{即}\qquad
\ip{\delta}{\varphi}=\varphi(0).\label{eq:deltadef}
\end{equation}

于是 \eqref{eq:sample} 就可以写成

\begin{equation}
K_\ep\longrightarrow\delta
\qquad(\ep\to0^+),
\end{equation}

而实际频谱是 $G_\ep=2\pi K_\ep$，所以

\begin{equation}
G_\ep=\frac{2\ep}{\ep^{2}+\omega^{2}}\longrightarrow2\pi\delta(\omega),
\qquad
\boxed{\ \F[1](\omega)=2\pi\delta(\omega)\ }.\label{eq:mainresult}
\end{equation}

\section{结论的准确含义}

这一节的每一句话，都是写旁白时的“标准答案”。

\subsection{能说的三句话}

\begin{enumerate}[leftmargin=2em,itemsep=0.5em]
\item \textbf{“等号是在广义函数意义下成立的。”}它不是说那个发散积分突然变成收敛积分；它说的是：这一族频谱在“对测试函数的作用”这个意义上收敛到 $2\pi\delta$。
\item \textbf{“常数的频谱集中在零频率，总权重是 $2\pi$。”}$2\pi$ 是一个权重（面积），不是一个高度。
\item \textbf{“为什么是 $2\pi$？因为逆变换里有一个 $1/(2\pi)$。”}\ 两个系数必须互相配套：正变换不带系数，逆变换带 $1/(2\pi)$，所以正变换在零频率上留下 $2\pi$。
\end{enumerate}

\subsection{不能说的三句话}

\begin{warnbox}[EP01 的三条硬禁令]
\begin{enumerate}[leftmargin=1.4em,itemsep=0.35em]
\item \textbf{不能说}“非零频率处原来的积分等于零”。它不存在。
\item \textbf{不能说}“$\delta$ 在原点无穷高”。它没有“取值”。
\item \textbf{不能说}“因为这族函数除了原点都趋于零，所以整体趋于零”。正好相反：整体不趋于零，它趋于一个集中到原点的对象。
\end{enumerate}
\end{warnbox}

\subsection{为什么“面积”不是“能量”}

有一个很容易犯的错：把 $\int G_\ep=2\pi$ 说成“能量守恒”，进而说“所以 $\delta$ 的能量是 $2\pi$”。

正确的图景是 Plancherel：

\begin{equation}
\|f\|_2^2=\frac{1}{2\pi}\|\fh\|_2^2 .
\end{equation}

对 $f_\ep$：左边 $=\int e^{-2\ep|t|}\dd t=\frac{1}{\ep}$，右边 $=\frac{1}{2\pi}\int G_\ep^2\dd\omega=\frac{1}{\ep}$。两边相等，没问题。

但是 $\ep\to0$ 时，$\|f_\ep\|_2^2=\frac{1}{\ep}\to\infty$。也就是说：\textbf{常数 $1$ 不是有限能量信号，它根本不住在 $\Ltwo$ 里。}所以“能量守恒”这条思路从一开始就不适用于它。

\begin{notebox}[要把这句话留到 EP06]
“$1$ 不在 $\Ltwo$，所以不能用 Plancherel；但它在 $\Sv'$，所以可以用对偶。”这两句要等到 EP06 才说得完整。EP01 只需要说：“这条曲线的面积不是信号能量，别把两者混为一谈。”
\end{notebox}

\subsection{一个独立的自检}

做完这一集，可以用另一条路验证 $\F[1]=2\pi\delta$：先算 $\F[\delta]$。

若 $\F[\delta]=1$，那么由反演（在 EP06 证明）应当有 $\F[1]=2\pi\delta$，因为反演说“变换两次基本回到自己，差一个反射和 $2\pi$”。

\begin{equation}
\ip{\F\delta}{\varphi}=\ip{\delta}{\ph}=\ph(0)=\int_{\R}\varphi(t)\dd t=\ip{1}{\varphi}
\quad\Longrightarrow\quad
\F[\delta]=1 .
\end{equation}

这条验证只需要一句 Fubini，非常干净。它同时也是 EP06 里“$\delta$ 与常数”那一段的预演。

\section{动画需要的关键数值}

\begin{center}
\small
\renewcommand{\arraystretch}{1.35}
\begin{tabularx}{\textwidth}{@{}L{1.4cm}L{2.2cm}L{2.2cm}L{2.4cm}Y@{}}
\toprule
$\ep$ & 峰高 $2/\ep$ & 半高宽 $2\ep$ & 面积 & 画面建议 \\
\midrule
$1$ & $2$ & $2$ & $2\pi$ & 纵轴到 $3$ 即可 \\
$0.5$ & $4$ & $1$ & $2\pi$ & 纵轴到 $5$ \\
$0.2$ & $10$ & $0.4$ & $2\pi$ & 纵轴到 $11$，注意采样加密 \\
$0.1$ & $20$ & $0.2$ & $2\pi$ & 纵轴到 $22$；\textbf{默认不是这个范围} \\
$0.02$ & $100$ & $0.04$ & $2\pi$ & 必须缩放或局部放大 \\
\bottomrule
\end{tabularx}
\end{center}

\begin{warnbox}[三个纯技术但会毁掉一集的坑]
\begin{enumerate}[leftmargin=1.4em,itemsep=0.35em]
\item \textbf{峰顶越界。}$\ep=0.1$ 时峰高 $20$。如果沿用 $\ep=1$ 时的纵轴（到 $3$），整个后半段动画的峰会被裁掉，观众看到的是“一条平线加一根竖线”，完全错误。
\item \textbf{采样不足。}$\ep=0.02$ 时半高宽只有 $0.04$。若在 $[-4,4]$ 上均匀取 $200$ 个点，峰顶根本采不到，画出来是个锯齿。必须用对数分布或在原点附近加密的参数曲线。
\item \textbf{面积标签。}可见画框里的阴影只是全实轴的一部分。标签必须写“全实轴面积 $=2\pi$”，否则观众会拿屏幕上的像素面积去核对，然后发现对不上。
\end{enumerate}
\end{warnbox}

\section{讲解设计：这一集要落空三次}

好的科普不是把正确的答案讲一遍，而是让观众的\textbf{预期}被依次打破。EP01 有三次：

\begin{enumerate}[leftmargin=2em,itemsep=0.6em]
\item \textbf{第一次落空}：直觉明明说“只有零频率”，可是积分算不出来。——观众开始怀疑“是不是我理解错了傅里叶变换”。
\item \textbf{第二次落空}：换了衰减外壳，算出很漂亮的峰；$\ep$ 变小，峰变高变窄，“除了原点都趋于零”。——观众以为这说明“频谱趋于零”，结果发现整体并不趋于零，而是集中。
\item \textbf{第三次落空}：以为 $\delta$ 是“一个特别高的尖峰”，结果发现它的意义根本不在高度，而在“它在积分里取走了 $\varphi(0)$”。
\end{enumerate}

\begin{ideabox}[三次落空各自对应一个动画]
\begin{itemize}[leftmargin=1.4em,itemsep=0.2em]
\item 第一次 $\to$ 01-B：截断积分不收敛。
\item 第二次 $\to$ 01-E：双面板，变高变窄，面积标签不变。
\item 第三次 $\to$ 01-F：阴影向原点集中，读数趋于 $\varphi(0)$。
\end{itemize}
如果某一集剪完发现只有“讲清楚了”，没有“落空感”，那多半是缺了其中一个动画。
\end{ideabox}

\section{讲解禁区（EP01）}

\begin{warnbox}[逐条对照，写旁白时逐句检查]
\begin{enumerate}[leftmargin=1.6em,itemsep=0.28em]
\item 不说“$\delta$ 是原点处无穷高的函数”。
\item 不说“非零频率处原积分 $=0$”。
\item 不说“$K_\ep$ 处处趋于零”（原点是例外，而且趋于无穷）。
\item 不说“因为除了原点都趋于零，所以整体趋于零”。
\item 不说“频谱面积守恒是普遍规律”；只说“本例这条曲线的全实轴面积恒为 $2\pi$”。
\item 不说“这就是信号的能量”；不写 $\|\delta\|^2$ 或 $\delta^2$。
\item 不说“$\ep=0$”或“现在它已经是常数了”。
\item 不混用 $K_\ep$（面积 $1$）与 $G_\ep$（面积 $2\pi$）的标注与纵轴。
\item 不在画面上只画一个频率峰来表示实余弦的频谱（实余弦有正负两个频率）。
\item 不说“把这族函数取极限就得到 $\delta$ 这个函数”——$\delta$ 不是函数，这句要等到 EP05 才说得完整。
\item 不声称“频率分量的大小就是峰高”；对 $\delta$ 而言是权重。
\item 如果缩放纵轴，必须让刻度变化可见，不能让屏幕像素面积冒充数学面积。
\end{enumerate}
\end{warnbox}

\section{与 EP02 的接口}

这一集结束在“$\F[1]=2\pi\delta$”。留下的钩子有两个，第二个才是 EP02 真正用的：

\begin{enumerate}[leftmargin=2em,itemsep=0.45em]
\item \textbf{弱的钩子}：“$\delta$ 到底是不是函数？”——这个问题要到 EP05 才正面回答，中间要吊三集。
\item \textbf{强的钩子}：“如果我把常数乘上 $t$，让它开始线性增长，频谱会怎样？”——这个问题下一集就能回答，而且答案会立刻引出新的疑问（$\delta$ 的导数是什么）。
\end{enumerate}

\begin{scriptbox}[最后 30 秒的旁白骨架]
“我们得到了一件很奇怪的东西：一个不是函数的符号，出现在一个等式的右边。

但更奇怪的是，它其实很好用。下一集，我们把这条水平线掰成一条斜线——然后看看会发生什么。”
\end{scriptbox}

\section{本集回顾卡}

\begin{center}
\small
\renewcommand{\arraystretch}{1.4}
\begin{tabularx}{\textwidth}{@{}L{3.6cm}Y@{}}
\toprule
项目 & 内容 \\
\midrule
主结果 & $\F[1](\omega)=2\pi\delta(\omega)$，广义函数意义 \\
核心计算 & $\F\big[e^{-\ep|t|}\big]=\dfrac{2\ep}{\ep^{2}+\omega^{2}}=G_\ep(\omega)$ \\
三个指标 & 峰高 $2/\ep$；半高全宽 $2\ep$；全实轴面积 $2\pi$ \\
归一化核 & $K_\ep=G_\ep/(2\pi)=\frac{1}{\pi}\frac{\ep}{\ep^{2}+\omega^{2}}$，面积 $1$ \\
安全话术 & “频谱集中在零频率，权重为 $2\pi$” \\
危险话术 & “$\delta$ 是原点无穷高”“非零频率处为零”“面积守恒是普遍规律” \\
下一集的钩子 & 把水平线掰成斜线：$\F[t]$ 会是什么？ \\
\bottomrule
\end{tabularx}
\end{center}

\begin{quizbox}[EP01 自测]
\begin{enumerate}[leftmargin=1.6em,itemsep=0.35em]
\item 从零开始，重算 $f_\ep=e^{-\ep|t|}$ 的傅里叶变换，并说明 $\ep>0$ 在哪里被用到。
\item $G_\ep$ 的半高全宽为什么是 $2\ep$？写出求解过程。
\item 证明 $\int_{\R}G_\ep\dd\omega=2\pi$，并解释为什么这个数不随 $\ep$ 变化。
\item 用 \eqref{eq:scaling} 验证 $G_\ep(\omega)=\frac{1}{\ep}\widehat{g}(\omega/\ep)$。
\item 如果某个动画把纵轴固定在 $[0,3]$，$\ep$ 减小到 $0.1$ 时画面会出什么问题？
\item 用一句话说明“面积不变”与“每点趋于零”为什么可以同时成立，而不矛盾。
\end{enumerate}
\end{quizbox}
}{}
\IfFileExists{ep02.tex}{\chapter{EP02｜多项式与 $\delta$ 的导数}

\begin{ideabox}[这一集的一句话]
\textbf{在时间上乘一个 $t$，等于在频率上求一次导。所以多项式只有零频率的支撑——但它记录的不是“一个直流值”，而是零点处的各阶导数。}
\end{ideabox}

\section{观众的问题链}

第一集留下的钩子是“把水平线掰成斜线”。这一集就从这里开始，但真正的高潮不是算出 $\F[t]$，而是被迫面对一个新问题：\textbf{对一个不是函数的东西求导，是什么意思？}

\begin{center}
\small
\renewcommand{\arraystretch}{1.32}
\begin{tabularx}{\textwidth}{@{}p{0.8cm}X p{4.4cm}@{}}
\toprule
层 & 问题 & 观众此刻的心理 \\
\midrule
1 & 常数有频谱，那 $t$ 呢？$t^{2}$ 呢？ & 越大越离谱，应该更算不出来。\\
2 & 直接代入，果然又发散了。 & 意料之中。\\
3 & 但能不能从 $\F[1]$ 推出来？ & 怎么推？\\
4 & 可以：时域乘 $t$ 对应频域求导。 & 这个规则很漂亮。\\
5 & 可是求导之后出现了 $\delta'$。 & 等等，$\delta'$ 是什么东西？\\
6 & 它不是尖峰，它有两叶，总积分为零。 & 那它到底“是什么”？\\
\bottomrule
\end{tabularx}
\end{center}

\begin{center}
\small
\renewcommand{\arraystretch}{1.32}
\begin{tabularx}{\textwidth}{@{}p{3.2cm}Y@{}}
\toprule
观众带着什么进来 & 观众带着什么出去 \\
\midrule
“$\delta$ 是一个极限的名字” & “$\delta$ 还可以被求导，而求导的含义由它对探针的作用决定。”\\
把“求导”理解成“看斜率” & “分布求导是把导数转移到探针上，代价是一个负号。”\\
以为“零频率”意味着“没有信息” & “零频率的 $\delta^{(n)}$ 恰好记录了 $\varphi$ 在零点的第 $n$ 阶导数。”\\
\bottomrule
\end{tabularx}
\end{center}

\section{先算一遍：直接代入为什么还是不行}

把 $f(t)=t$ 代进定义：

\begin{equation}
\int_{-T}^{T}t\,e^{-i\omega t}\dd t .
\end{equation}

分部积分（取 $u=t$，$\dd v=e^{-i\omega t}\dd t$）：

\begin{equation}
\int_{-T}^{T}t\,e^{-i\omega t}\dd t
=\left[\frac{t\,e^{-i\omega t}}{-i\omega}\right]_{-T}^{T}
+\frac{1}{i\omega}\int_{-T}^{T}e^{-i\omega t}\dd t
=\frac{2T\cos(\omega T)}{i\omega}+\frac{2\sin(\omega T)}{\omega^{2}} .
\end{equation}

第一项带着因子 $T$，$\cos(\omega T)$ 又不停振荡，所以它的振幅随 $T$ \textbf{线性增长}。这比常数的情况更糟：常数只是“不收敛”，$t$ 是“振荡着发散”。

\begin{center}
\small
\renewcommand{\arraystretch}{1.3}
\begin{tabularx}{\textwidth}{@{}p{2.2cm}L{5.0cm}Y@{}}
\toprule
对象 & 截断积分的量级 & 失败方式 \\
\midrule
$1$ & $\dfrac{2\sin(\omega T)}{\omega}$，振幅恒定 & 不收敛 \\
$t$ & $\dfrac{2T\cos(\omega T)}{i\omega}+\cdots$，振幅 $\sim T$ & 振荡着发散 \\
$t^{n}$ & 振幅 $\sim T^{n}$ & 发散得更快 \\
\bottomrule
\end{tabularx}
\end{center}

结论很清楚：\textbf{这条路走不通，一代比一代更不通。}必须换一条路——从已经得到的 $\F[1]$ 出发。

\section{换一条路：时域乘 $t$，频域求导}

\subsection{在一个足够好的函数上先验一遍}

先假设 $f$ 属于 Schwartz 空间 $\Sv$（光滑、各阶导数都快速衰减）。这样下面每一步都合法，不用边算边担心。对

\begin{equation}
\fh(\omega)=\int_{\R}f(t)e^{-i\omega t}\dd t
\end{equation}

关于 $\omega$ 求导。$e^{-i\omega t}$ 对 $\omega$ 的导数是 $-it\,e^{-i\omega t}$，于是

\begin{equation}
\frac{\dd}{\dd\omega}\fh(\omega)
=\int_{\R}f(t)\,(-it)\,e^{-i\omega t}\dd t
=-i\int_{\R}\big(tf(t)\big)e^{-i\omega t}\dd t
=-i\,\F[tf](\omega).
\end{equation}

整理一下：

\begin{equation}
\boxed{\ \F[tf](\omega)=i\,\frac{\dd}{\dd\omega}\fh(\omega)\ }\label{eq:mult-t}
\end{equation}

\begin{ideabox}[为什么必须写“$f\in\Sv$”]
如果把求导搬到积分号里面，需要两个条件：
\begin{itemize}[leftmargin=1.4em,itemsep=0.2em]
\item 被积函数对 $\omega$ 可导（这里显然），
\item 导数 $\big|(-it)f(t)e^{-i\omega t}\big|=|tf(t)|$ 能被一个\textbf{与 $\omega$ 无关}的可积函数控制。
\end{itemize}
$f\in\Sv$ 时 $|tf(t)|$ 可积，第二条自动满足。如果只假设 $f\in\Lone$，$tf$ 未必可积，这一步就可能是错的。
\end{ideabox}

\subsection{两条规则，方向相反，最容易记混}

\begin{center}
\small
\renewcommand{\arraystretch}{1.45}
\begin{tabularx}{\textwidth}{@{}L{5.2cm}L{5.4cm}Y@{}}
\toprule
时域做什么 & 频域做什么 & 用在哪一集 \\
\midrule
$t$ 乘上去 & 求导，再乘 $i$ & \textbf{EP02} \\
求导 & 乘 $i\omega$ & \textbf{EP03} \\
\bottomrule
\end{tabularx}
\end{center}

写成公式：

\begin{equation}
\F[tf]=i\,\partial_\omega\fh,
\qquad
\F[f']=i\omega\,\fh .\label{eq:two-rules}
\end{equation}

第二条的推导只要一次分部积分（边界项因为 $f$ 衰减而消失）：

\begin{equation}
\int_{\R}f'(t)e^{-i\omega t}\dd t
=\big[f(t)e^{-i\omega t}\big]_{-\infty}^{\infty}
+i\omega\int_{\R}f(t)e^{-i\omega t}\dd t
=i\omega\fh(\omega).
\end{equation}

\begin{notebox}[一个帮助记忆的对照]
\begin{itemize}[leftmargin=1.4em,itemsep=0.2em]
\item 时域\textbf{乘} $t$ $\longleftrightarrow$ 频域\textbf{求导}。
\item 时域\textbf{求导} $\longleftrightarrow$ 频域\textbf{乘} $i\omega$。
\end{itemize}
像照镜子：一边“乘”，另一边就“导”。两个 $i$ 的来历不同（一个来自 $e^{-i\omega t}$ 对 $\omega$ 求导，一个来自对 $t$ 分部积分），但都出现在同一位置，这就是它们容易混的原因。
\end{notebox}

\section{$\delta$ 的导数：负号是从哪里来的}

\subsection{从分部积分说起}

对普通的光滑函数 $f$ 和紧支撑的光滑函数 $\varphi$，

\begin{equation}
\int_{\R}f'(t)\varphi(t)\dd t
=-\int_{\R}f(t)\varphi'(t)\dd t .\label{eq:ibp}
\end{equation}

边界项消失了：$\varphi$ 紧支撑，所以两端都为 $0$。

这个等式本来只是“分部积分”，但它有一个了不起的结构：\textbf{右边的 $f$ 没有再被求导。}导数从 $f$ 身上转移到了 $\varphi$ 身上。

于是可以反过来\textbf{用右边来定义左边}：即使 $f$ 没有导数，甚至 $f$ 根本不是函数，只要它还是一个能作用于 $\varphi$ 的对象，右边就仍然有意义。

\begin{defn}[分布导数]\label{def:dist-deriv}
设 $T$ 是 $C_c^{\infty}(\R)$ 上的连续线性泛函。定义它的导数 $T'$ 为
\begin{equation}
\ip{T'}{\varphi}:=-\ip{T}{\varphi'},
\qquad
\ip{T^{(n)}}{\varphi}:=(-1)^{n}\ip{T}{\varphi^{(n)}} .
\end{equation}
\end{defn}

\begin{ideabox}[为什么这个定义是“对的”而不是“编的”]
三条理由，按重要性排列：
\begin{enumerate}[leftmargin=1.4em,itemsep=0.25em]
\item \textbf{它和经典导数兼容。}若 $T=T_f$ 且 $f$ 本来就可导，那么由 \eqref{eq:ibp}，新定义的 $T'$ 就是 $T_{f'}$。
\item \textbf{它不需要 $T$ 有逐点导数。}整个定义只用到 $\varphi$ 的光滑性，完全没问 $T$“在某点是多少”。
\item \textbf{它自动是连续的。}若 $T_n\to T$，则 $\ip{T_n'}{\varphi}=-\ip{T_n}{\varphi'}\to-\ip{T}{\varphi'}=\ip{T'}{\varphi}$。求导不会破坏收敛。
\end{enumerate}
\end{ideabox}

\subsection{$\delta'$ 的作用}

把 $T=\delta$ 代进定义：

\begin{equation}
\ip{\delta'}{\varphi}=-\ip{\delta}{\varphi'}=-\varphi'(0),
\qquad
\ip{\delta''}{\varphi}=(-1)^{2}\ip{\delta}{\varphi''}=\varphi''(0).
\end{equation}

一般地，

\begin{equation}
\boxed{\ \ip{\delta^{(n)}}{\varphi}=(-1)^{n}\varphi^{(n)}(0)\ }\label{eq:delta-n}
\end{equation}

\begin{warnbox}[负号不是笔误]
$\ip{\delta'}{\varphi}=-\varphi'(0)$ 里的负号，是分部积分的那个负号。它不会因为“对称”而消失。

很多人在这一步丢掉负号，然后整条递推全错：$\F[t]$ 会写成 $-2\pi i\delta'$，$\F[t^{2}]$ 的符号也会翻。

记住它来自哪里，比记住它在不在，更可靠。
\end{warnbox}

\subsection{一句话读懂每一阶}

\begin{center}
\small
\renewcommand{\arraystretch}{1.4}
\begin{tabularx}{\textwidth}{@{}L{2.6cm}L{4.4cm}Y@{}}
\toprule
对象 & 作用 & 用一句人话说 \\
\midrule
$\delta$ & $\varphi(0)$ & 读出探针在原点的\textbf{值} \\
$\delta'$ & $-\varphi'(0)$ & 读出探针在原点的\textbf{斜率}（带负号） \\
$\delta''$ & $\varphi''(0)$ & 读出探针在原点的\textbf{弯曲程度} \\
\bottomrule
\end{tabularx}
\end{center}

\begin{ideabox}[一句可以直接用的旁白]
“$\delta$ 关心的是你的函数在原点有多高，$\delta'$ 关心的是它在那里有多斜。它们不是‘一个尖峰’和‘一个更尖的尖峰’，而是\textbf{两种不同的提问方式}。”

这句话说清楚，EP05 的分布定义就有一半已经铺好了。
\end{ideabox}

\section{把 $\delta'$ 画出来：两叶，而且总积分为零}

\subsection{对第一集的核求导}

第一集的核

\begin{equation}
K_\ep(x)=\frac{1}{\pi}\frac{\ep}{x^{2}+\ep^{2}},
\qquad
\int_{\R}K_\ep\dd x=1
\end{equation}

对 $x$ 求导：

\begin{equation}
K_\ep'(x)=\frac{1}{\pi}\cdot\frac{-2\ep x}{(x^{2}+\ep^{2})^{2}}
=-\frac{2\ep x}{\pi\,(x^{2}+\ep^{2})^{2}} .\label{eq:Keps-prime}
\end{equation}

\begin{center}
\small
\renewcommand{\arraystretch}{1.35}
\begin{tabularx}{\textwidth}{@{}L{4.0cm}L{4.6cm}Y@{}}
\toprule
性质 & 结果 & 意味着什么 \\
\midrule
零点 & $x=0$ & 中间穿过横轴 \\
符号 & $x<0$ 为正，$x>0$ 为负 & 左叶朝上，右叶朝下 \\
极值位置 & $x=\pm\ep/\sqrt{3}$ & 两叶对称，越靠越近 \\
峰值大小 & $\dfrac{9}{8\sqrt{3}\,\pi\,\ep^{2}}\approx\dfrac{0.207}{\ep^{2}}$ & 按 $\ep^{-2}$ 增长，比 $K_\ep$ 快 \\
总积分 & $\displaystyle\int_{\R}K_\ep'\dd x=K_\ep(+\infty)-K_\ep(-\infty)=0$ & \textbf{正负面积抵消} \\
总变差 & $\displaystyle\int_{\R}|K_\ep'|\dd x=\frac{2}{\pi\ep}\to\infty$ & 不是测度 \\
\bottomrule
\end{tabularx}
\end{center}

\begin{figure}[htbp]
\centering
\begin{tikzpicture}
\begin{axis}[width=0.46\textwidth,height=4.4cm,
  domain=-2:2,samples=300,ymin=0,ymax=1.15,
  xlabel={$x$},ylabel={$K_\ep(x)$},
  title={\small 核本身：一个非负峰},
  ticklabel style={font=\scriptsize},label style={font=\small}]
\addplot[thick,mainblue] {0.5/(pi*(x^2+0.25))};
\addplot[only marks,mark=*,mark size=1.2pt,mainblue] coordinates {(0,0.6366)};
\end{axis}
\end{tikzpicture}\hfill
\begin{tikzpicture}
\begin{axis}[width=0.46\textwidth,height=4.4cm,
  domain=-2:2,samples=300,ymin=-1.1,ymax=1.1,
  xlabel={$x$},ylabel={$K_\ep'(x)$},
  title={\small 它的导数：左正右负},
  ticklabel style={font=\scriptsize},label style={font=\small},
  axis lines=middle]
\addplot[thick,warnred] {-2*0.5*x/(pi*(x^2+0.25)^2)};
\end{axis}
\end{tikzpicture}
\caption{$\ep=0.5$ 时的 $K_\ep$ 与 $K_\ep'$。左边那一叶的面积为正，右边为负，两者绝对值相等。\textbf{这不是 $\delta'$ 的点值图像，只是它的一族平滑近似。}}
\end{figure}

\subsection{为什么 $\delta'$ 不能用“面积”定义}

这是一个关键的观念转折，值得在视频里单独停一下。

$\delta$ 可以用“总质量 $=1$”加上“集中到原点”来描述。但 $\delta'$ 不行：

\begin{equation}
\int_{\R}K_\ep'(x)\dd x=0
\qquad\text{对每个 }\ep>0 .
\end{equation}

所以任何“总质量为 $1$”的说法在这里完全失效。$\delta'$ 的“信息量”不在总量里，而在\textbf{正负两叶的分布方式}里。

看它的作用就明白了：

\begin{equation}
\ip{K_\ep'}{\varphi}
=-\ip{K_\ep}{\varphi'}
\;\longrightarrow\;-\varphi'(0)=\ip{\delta'}{\varphi} .
\end{equation}

第一行是分部积分（这里 $K_\ep$ 光滑且衰减，合法）。第二行用了 EP01 的取样极限，只是把 $\varphi$ 换成了 $\varphi'$。

\begin{ideabox}[这就是“求导在分布意义下是连续的”]
$\ip{K_\ep'}{\varphi}=-\ip{K_\ep}{\varphi'}\to-\varphi'(0)$。

整个论证没有用到 $K_\ep'$ 的\textbf{点值}收敛。事实上 $K_\ep'$ 在每一点都不收敛到任何普通函数（在原点是 $0$，但峰值在别处爆炸）。

\textbf{结论：我们不需要“$\delta'$ 的点值图像”，也能严格说清它做什么。}这正是 EP05 要系统化的想法。
\end{ideabox}

\section{从 $1$ 走到所有多项式}

\subsection{逐步递推}

把 \eqref{eq:mult-t} 用在 $t\cdot1$ 上，并借用第一集的结果 $\F[1]=2\pi\delta$：

\begin{align}
\F[t]&=i\,\partial_\omega\F[1]=i\,\partial_\omega\big(2\pi\delta\big)=2\pi i\,\delta',\\[2pt]
\F[t^{2}]&=i\,\partial_\omega\F[t]=i\,\partial_\omega\big(2\pi i\,\delta'\big)=2\pi i^{2}\,\delta''=-2\pi\,\delta'',\\[2pt]
\F[t^{3}]&=i\,\partial_\omega\F[t^{2}]=i\,\partial_\omega\big(-2\pi\delta''\big)=-2\pi i\,\delta''' .
\end{align}

规律已经出来了：每乘一次 $t$，频域就求一次导，同时乘一个 $i$。归纳得到

\begin{equation}
\boxed{\ \F[t^{n}]=2\pi i^{\,n}\,\delta^{(n)}(\omega),\qquad n=0,1,2,\dots\ }\label{eq:tn}
\end{equation}

\begin{center}
\small
\renewcommand{\arraystretch}{1.4}
\begin{tabularx}{\textwidth}{@{}L{1.4cm}L{3.2cm}L{3.2cm}Y@{}}
\toprule
$n$ & $i^{n}$ & $\F[t^{n}]$ & 符号 \\
\midrule
$0$ & $1$ & $2\pi\delta$ & $+$ \\
$1$ & $i$ & $2\pi i\,\delta'$ & 纯虚 \\
$2$ & $-1$ & $-2\pi\,\delta''$ & 实 \\
$3$ & $-i$ & $-2\pi i\,\delta'''$ & 纯虚 \\
$4$ & $1$ & $2\pi\,\delta^{(4)}$ & 实 \\
\bottomrule
\end{tabularx}
\end{center}

由线性性，对任意多项式 $p(t)=\sum_{n=0}^{N}a_n t^{n}$，

\begin{equation}
\widehat{p}(\omega)=2\pi\sum_{n=0}^{N}a_n i^{\,n}\,\delta^{(n)}(\omega).\label{eq:poly}
\end{equation}

\begin{warnbox}[这不是“把多项式塞进一个收敛的积分”]
多项式不是 $\Lone$ 函数，$\int p(t)e^{-i\omega t}\dd t$ 依然发散。\eqref{eq:poly} 的合法性来自两件事：
\begin{itemize}[leftmargin=1.4em,itemsep=0.2em]
\item 多项式通过“对测试函数积分”成为一个温和分布（EP05--EP06 建立）；
\item 乘 $t$ 的规则 \eqref{eq:mult-t} 在这个更大的空间里仍然成立（EP06 证明）。
\end{itemize}
EP02 是\textbf{借用}这两件事。旁白必须说清楚“先借用、后补证”，否则观众会以为我们在偷偷把发散积分当成收敛积分。
\end{warnbox}

\subsection{一个立刻能做的小检查}

用 \eqref{eq:poly} 算 $\F[3-2t+t^{2}]$：

\begin{equation}
\F[3]-2\F[t]+\F[t^{2}]
=3\cdot2\pi\delta-2\cdot2\pi i\,\delta'+(-2\pi\,\delta'')
=6\pi\delta-4\pi i\,\delta'-2\pi\,\delta'' .
\end{equation}

检查虚部：只有 $-2t$ 这一项是奇函数，它贡献纯虚的项；$3+t^{2}$ 是偶的，贡献实的项。实虚分开，符合“实信号的频谱满足 $\overline{\fh(\omega)}=\fh(-\omega)$”。

\section{为什么频谱仍然支撑在零}

\subsection{“支撑”是什么意思}

先给一个临时定义（EP05 会正式化）：\textbf{一个分布在某点附近为零}，指的是它对所有支撑在该点附近的测试函数的作用都为零。\textbf{支撑}就是它不能这样“消掉”的区域。

对 $\delta^{(n)}$：若 $\varphi$ 在 $0$ 的某个邻域内恒为 $0$，则 $\varphi^{(n)}(0)=0$，于是

\begin{equation}
\ip{\delta^{(n)}}{\varphi}=(-1)^{n}\varphi^{(n)}(0)=0 .
\end{equation}

所以

\begin{equation}
\supp\delta^{(n)}=\{0\}.
\end{equation}

\subsection{反过来：支撑在零的对象是什么样}

这一步很多人会含糊过去，但它是本集最值得补上的一个小证明。

\begin{prop}\label{prop:support0}
若分布 $T$ 满足 $\supp T\subseteq\{0\}$，则存在整数 $N$ 和常数 $c_0,\dots,c_N$，使
\begin{equation}
T=\sum_{n=0}^{N}c_n\,\delta^{(n)} .
\end{equation}
\end{prop}

\begin{proof}[证明思路：泰勒展开把“零点附近的行为”全部抓出来]
取 $\chi\in C_c^{\infty}$，$\chi\equiv1$ 于 $|x|\le1$。对任意 $\varphi\in C_c^{\infty}$ 和任意 $N$，由带积分余项的泰勒公式，
\begin{equation}
\varphi(x)=\sum_{n=0}^{N}\frac{\varphi^{(n)}(0)}{n!}x^{n}+x^{N+1}\psi(x),
\end{equation}
其中 $\psi$ 是光滑函数（余项积分给出）。

把 $T$ 作用上去。第二项里 $T$ 作用在 $x^{N+1}\psi$ 上。由分布的连续性（每个紧集上存在一个有限阶的估计），存在有限的阶 $m$；取 $N+1>m$，则 $x^{N+1}\psi$ 与它的前 $m$ 阶导数在 $0$ 处都为零。把 $x^{N+1}\psi$ 在一小段区间外截断（截断函数在 $0$ 附近恒为 $1$），得到的仍是一个“零点附近行为相同”的函数，而 $T$ 只看得见 $0$ 的邻域，所以这一项为零。

剩下
\begin{equation}
\ip{T}{\varphi}=\sum_{n=0}^{N}\frac{\varphi^{(n)}(0)}{n!}\ip{T}{x^{n}}
=\sum_{n=0}^{N}\tilde c_n\,\varphi^{(n)}(0)
=\sum_{n=0}^{N}\big((-1)^{n}\tilde c_n\big)\ip{\delta^{(n)}}{\varphi}.
\end{equation}
\end{proof}

\begin{ideabox}[这个命题的真正用处]
它告诉我们：\textbf{“支撑在原点”不是一句空话，它精确地意味着“由 $\delta$ 及其各阶导数线性组合而成”。}

所以“多项式的频谱支撑在零”这句话，等价于“多项式的频谱完全由 $\delta,\delta',\delta'',\dots$ 构成”。这正是 \eqref{eq:poly} 说的事。
\end{ideabox}

\subsection{另一种检查：用微分方程}

$p$ 是 $N$ 次多项式，则 $p^{(N+1)}=0$。用 \eqref{eq:two-rules} 的第二条，

\begin{equation}
\F\big[p^{(N+1)}\big]=(i\omega)^{N+1}\widehat{p}=0 .
\end{equation}

即

\begin{equation}
\omega^{N+1}\,\widehat{p}=0 .
\end{equation}

由命题 \ref{prop:support0} 的同类论证（把那里的 $T$ 换成 $\widehat p$），$\widehat p$ 只能由 $\delta,\delta',\dots,\delta^{(N)}$ 组合而成。

\begin{warnbox}[但这个检查只能定“位置”，不能定“系数”]
$\omega^{N+1}\widehat p=0$ 只说明频谱住在原点，并不告诉你每个 $\delta^{(n)}$ 前面的系数是多少。系数必须靠乘 $t$ 的递推 \eqref{eq:tn} 一步步算出来。

\textbf{两种检查的关系：}微分方程给“支撑”，递推给“系数”。缺一不可。
\end{warnbox}

\subsection{“支撑在零”不等于“不随时间变化”}

这是本集最容易被误解的一句话。

\begin{center}
\small
\renewcommand{\arraystretch}{1.4}
\begin{tabularx}{\textwidth}{@{}L{3.0cm}L{5.4cm}Y@{}}
\toprule
说法 & 对不对 & 为什么 \\
\midrule
$\F[t]$ 只用到零频率 & 对 & 支撑是 $\{0\}$ \\
所以 $t$ 是一个“直流信号” & \textbf{错} & $t$ 随时间线性增长，一点都不“直流” \\
所以 $t$ 与常数在频域没区别 & \textbf{错} & 常数给 $2\pi\delta$，$t$ 给 $2\pi i\delta'$，作用完全不同 \\
\bottomrule
\end{tabularx}
\end{center}

正确的说法是：\textbf{多项式的“频率成分”全部堆在零频率这一点上，但它们以“$\delta$ 的哪一阶导数”的形式堆在那里。}阶数越高，记录的信息越“局部”、越“高阶”。

\begin{ideabox}[一个能救场的比喻]
把 $\delta^{(n)}$ 想成\textbf{不同的探针}：$\delta$ 问“值是多少”，$\delta'$ 问“斜率是多少”，$\delta''$ 问“弯得有多厉害”。

它们都只关心原点，所以都“支撑在零”。但它们问的问题完全不同，所以不能互相替换。就像同一个医生只用听诊器，也可以问出完全不同的三件事。
\end{ideabox}

\section{数值与动画}

\begin{center}
\small
\renewcommand{\arraystretch}{1.35}
\begin{tabularx}{\textwidth}{@{}L{1.4cm}L{3.4cm}L{3.2cm}L{3.0cm}Y@{}}
\toprule
$\ep$ & $K_\ep(0)=1/(\pi\ep)$ & $K_\ep'$ 峰值 $\approx0.207/\ep^{2}$ & 极值位置 $\ep/\sqrt3$ & 画面 \\
\midrule
$1$ & $0.318$ & $0.207$ & $0.577$ & 纵轴 $\pm0.3$ \\
$0.5$ & $0.637$ & $0.827$ & $0.289$ & 纵轴 $\pm1$ \\
$0.2$ & $1.59$ & $5.17$ & $0.115$ & 纵轴 $\pm6$，采样加密 \\
$0.1$ & $3.18$ & $20.7$ & $0.058$ & 纵轴 $\pm22$ \\
\bottomrule
\end{tabularx}
\end{center}

\begin{warnbox}[三个动画陷阱]
\begin{enumerate}[leftmargin=1.4em,itemsep=0.3em]
\item \textbf{用单箭头画 $\delta'$。}单箭头只能表达“权重为正”。$\delta'$ 有正负两叶、总积分为零，用单箭头画它等于把这一集的核心讲反了。
\item \textbf{把 $K_\ep'$ 说成 $\delta'$ 的图像。}它只是近似。$\ep\to0$ 时它不逐点收敛到任何函数；收敛只在“对测试函数的作用”这个意义下成立。
\item \textbf{把峰值当成“权重”。}$K_\ep'$ 的峰值按 $\ep^{-2}$ 爆炸，总积分却恒为 $0$。这两个数都不是 $\delta'$ 的“大小”。
\end{enumerate}
\end{warnbox}

\begin{scriptbox}[动画要点 02-D]
给测试函数 $\varphi$ 加一个可控的局部斜率（例如 $\varphi(x)=e^{-x^{2}}+\lambda x e^{-x^{2}}$），让 $\lambda$ 连续变化。

观众会看到：$\ip{K_\ep'}{\varphi}=-\varphi'(0)$ 的读数随 $\lambda$ 线性变化，而且\textbf{方向相反}（$\lambda$ 增大，读数变小）。

这一下就把“负号”和“读斜率”两件事同时讲清楚了。
\end{scriptbox}

\section{讲解禁区（EP02）}

\begin{warnbox}[逐条对照]
\begin{enumerate}[leftmargin=1.6em,itemsep=0.28em]
\item 不说“$\delta'$ 是另一个更尖的非负峰”。它有两叶，一正一负。
\item 不说“$\delta'$ 的面积是 $1$”。它的总积分是 $0$。
\item 不说“对 $\delta$ 求导就是看它的斜率”——$\delta$ 没有图像，谈不上斜率；要讲“把导数转移到探针上”。
\item 不丢 \eqref{eq:delta-n} 里的 $(-1)^{n}$。
\item 不说“多项式只有零频率，所以它是直流/不变的”。
\item 不说“$K_\ep'$ 逐点收敛到 $\delta'$”；只说分布意义收敛。
\item 不说“$\F[t^{n}]$ 是把这个多项式代进积分算出来的”。
\item 不把“乘 $t$”与“求导”两条规则的方向搞反：乘 $t$ $\to$ 频域求导；求导 $\to$ 频域乘 $i\omega$。
\item 不在没有说明“广义意义”的情况下直接写 $\F[t]=2\pi i\delta'$。
\item 不声称“任意分布都可以随便再求导又随便乘东西”——乘光滑函数可以，乘另一个分布不行（EP05 讲）。
\end{enumerate}
\end{warnbox}

\section{与前后集的接口}

\begin{enumerate}[leftmargin=2em,itemsep=0.45em]
\item \textbf{回头看}：这一集两次用到 EP01 的结果（$\F[1]=2\pi\delta$ 与 $K_\ep\to\delta$）。如果第一集的“$\delta$ 是极限的名字”没讲透，这里的求导会显得像魔术。
\item \textbf{往前看}：$\delta'$ 已经让观众不舒服了。但真正让他们\textbf{必须}要一个定义的时刻，是下一集——因为阶跃的频谱里会出现一个“除以 $\omega$”的麻烦，而微分规则给不出唯一答案。
\end{enumerate}

\begin{scriptbox}[最后 30 秒的旁白骨架]
“我们发现了一条很漂亮的规则：时间上乘 $t$，频率上就求导。

可是这条规则让我们写出了一个更奇怪的东西：$\delta$ 的导数。

而且它不像 $\delta$ 那样只有一根箭头——它有两叶，总面积为 $0$。

那么，‘对 $\delta$ 求导’到底是什么意思？

在回答这个问题之前，我们先看一个更麻烦的例子。”
\end{scriptbox}


\section{本集回顾卡}

\begin{center}
\small
\renewcommand{\arraystretch}{1.4}
\begin{tabularx}{\textwidth}{@{}L{3.4cm}Y@{}}
\toprule
项目 & 内容 \\
\midrule
主结果 & $\F[t^{\,n}]=2\pi i^{\,n}\delta^{(n)}(\omega)$ \\
核心规则 & $\F[tf]=i\,\partial_\omega\fh$；对照 $\F[f']=i\omega\fh$ \\
$\delta$ 的导数 & $\ip{\delta^{(n)}}{\varphi}=(-1)^{n}\varphi^{(n)}(0)$，负号来自分部积分 \\
可视化 & $K_\ep'$ 左正右负，极值在 $\pm\ep/\sqrt3$，总积分恒为 $0$ \\
安全话术 & “$\delta'$ 记录探针在原点的斜率” \\
危险话术 & “$\delta'$ 是更尖的峰”“$\delta'$ 的面积是 $1$” \\
下一集的钩子 & 时域求导对应频域做什么？能否由 $H'=\delta$ 反求 $\hf$？ \\
\bottomrule
\end{tabularx}
\end{center}

\begin{quizbox}[EP02 自测]
\begin{enumerate}[leftmargin=1.6em,itemsep=0.35em]
\item 取 $\varphi(x)=e^{-x^{2}}+xe^{-x^{2}}$，求 $\ip{\delta}{\varphi}$、$\ip{\delta'}{\varphi}$、$\ip{\delta''}{\varphi}$。
\item 求 $\F[3-2t+t^{2}]$，并检查实虚与奇偶。
\item 写出 $\F[t^{4}]$，说明 $i^{4}$ 与符号。
\item 证明 $\int_{\R}K_\ep'\dd x=0$，并解释为什么这使“面积”不能用来定义 $\delta'$。
\item 复述命题 \ref{prop:support0} 的证明思路：泰勒展开在哪一步起作用？为什么 $T$ 作用在 $x^{N+1}\psi$ 上会消失？
\item 为什么“$\omega^{N+1}\widehat p=0$”只能定支撑、不能定系数？
\item 用一句话向观众解释“$\delta'$ 记录斜率”这件事，不许出现“面积”这个词。
\end{enumerate}
\end{quizbox}
}{}
\IfFileExists{ep03.tex}{\chapter{EP03｜阶跃、主值与那个多出来的 $\delta$}

\begin{ideabox}[这一集的一句话]
\textbf{微分规则只能定出“除掉零频率以外”的那一部分频谱。零频率那一项在乘 $\omega$ 的时候被抹掉了，所以必须另外补回来。}
\end{ideabox}

\section{观众的问题链}

前两集里，$\delta$ 一直是“结论里出现的符号”。这一集它第一次变成“必须被解释的麻烦”：微分规则给了我们一个方程，而这个方程\textbf{有无穷多个解}。

\begin{center}
\small
\renewcommand{\arraystretch}{1.32}
\begin{tabularx}{\textwidth}{@{}p{0.8cm}X p{4.4cm}@{}}
\toprule
层 & 问题 & 观众此刻的心理 \\
\midrule
1 & 一个信号突然从 $0$ 跳到 $1$，它的导数是什么？ & 除了跳的那一下，别处都平，应该是零吧。\\
2 & 用平滑近似看，导数越来越集中在那一下。 & 所以导数是一个冲激？\\
3 & 既然如此，能不能用“时间求导 = 乘 $i\omega$”反推出阶跃的频谱？ & 这看起来很顺。\\
4 & 得到 $i\omega\widehat H=1$，只差除以 $i\omega$ 了。 & 那就除啊。\\
5 & 不能除。左边除以 $\omega$ 在原点有奇性，右边还要提防“乘 $\omega$ 会抹掉 $\delta$”。 & 等等，哪里来的 $\delta$？\\
6 & 因为 $\omega\delta=0$。所以这个方程的解不唯一。 & 那我们丢了什么信息？\\
\bottomrule
\end{tabularx}
\end{center}

\begin{center}
\small
\renewcommand{\arraystretch}{1.32}
\begin{tabularx}{\textwidth}{@{}p{3.2cm}Y@{}}
\toprule
观众带着什么进来 & 观众带着什么出去 \\
\midrule
“求导规则很好用” & “规则会丢信息；丢的那部分正好是零频率项。”\\
把 $\PV$ 看成“忽略奇点的小技巧” & “$\PV$ 是一种规定了对称方式的极限，它本身就是严格的对象。”\\
以为 $\F[H]$ 可以直接除出来 & “必须再做一次独立计算（正则化），才能定出那个 $\pi$。”\\
\bottomrule
\end{tabularx}
\end{center}

\section{阶跃与冲激：$H'=\delta$}

\subsection{阶跃的定义，以及为什么 $H(0)$ 无所谓}

\begin{equation}
H(t)=\begin{cases}0,&t<0,\\[2pt] 1,&t>0.\end{cases}
\end{equation}

$t=0$ 处的取值可以任选（$0$、$1/2$、$1$ 都行），因为它只影响一个点，而积分看不见一个点。

\begin{prop}\label{prop:Hprime}
作为分布，$H'=\delta$。
\end{prop}

\begin{proof}
对任意 $\varphi\in C_c^{\infty}$，由分布导数的定义与分部积分，
\begin{equation}
\ip{H'}{\varphi}
=-\ip{H}{\varphi'}
=-\int_{0}^{\infty}\varphi'(t)\dd t
=-\big[\varphi(t)\big]_{0}^{\infty}
=-\big(0-\varphi(0)\big)=\varphi(0)=\ip{\delta}{\varphi} .
\end{equation}
边界项消失是因为 $\varphi$ 紧支撑：$t\to\infty$ 时 $\varphi=0$；$t=0$ 处出现的是 $-\varphi(0)$，正好贡献出 $\delta$。
\end{proof}

\begin{ideabox}[这段证明在讲什么]
\textbf{一个跳跃，在求导之后变成一个点质量。}跳跃的“高度”就是 $\delta$ 的“权重”。

这不是近似，而是精确的分布等式。它的全部内容就是：如果你把 $\varphi$ 从 $0$ 积分到 $\infty$，得到的是 $-\varphi(0)$——而“$-\varphi(0)$”正是 $\delta$ 的招牌动作。
\end{ideabox}

\subsection{用平滑近似建立直觉}

取越来越陡的平滑阶跃（例如 $H_n(t)=\frac{1}{2}\big(1+\tanh(nt)\big)$，或 $\frac{1}{\pi}\arctan(nt)+\frac12$）。它们的导数

\begin{equation}
H_n'(t)=\frac{n}{2}\,\mathrm{sech}^{2}(nt)
\end{equation}

是一个非负的峰，峰值 $\sim n/2$，宽度 $\sim1/n$，总面积恒为 $1$。

\begin{scriptbox}[动画要点 03-A]
左边：阶跃从缓变陡。右边：它的导数从“缓坡”变成“窄峰”。

关键标注：\textbf{总面积始终为 $1$}。观众应当看到“跳变强度没有变，只是集中了”。

最后一句旁白：“跳变高度是 $1$，所以这个冲激的权重是 $1$。”
\end{scriptbox}

\subsection{一般规则（可以不说，但你必须知道）}

若 $f$ 分段 $C^{1}$，在点 $a_j$ 处有跳跃 $J_j=f(a_j+)-f(a_j-)$，则

\begin{equation}
Df=f'_{\text{经典}}+\sum_{j}J_j\,\delta_{a_j} .\label{eq:jump}
\end{equation}

也就是说：\textbf{分布导数 = 各段内部的光滑导数 + 每个跳跃贡献的一个 $\delta$。}

例：$|t|$ 在 $0$ 处跳跃的“导数”是符号函数，再求导得到 $2\delta$：

\begin{equation}
D|t|=\sgn t,
\qquad
D^{2}|t|=2\delta .
\end{equation}

这个例子很好用：它能一次性展示“分段光滑”“跳跃”“$\delta$ 的系数”三件事。

\section{能不能直接除以 $i\omega$}

\subsection{一步一步套规则}

由 \eqref{eq:two-rules} 的第二条，$\F[H']=i\omega\hf$。左边是 $\F[\delta]=1$。所以

\begin{equation}
i\omega\,\hf(\omega)=1 .\label{eq:iomegaH}
\end{equation}

看上去只差一步：

\begin{equation}
\hf(\omega)\overset{?}{=}\frac{1}{i\omega}.
\end{equation}

\begin{warnbox}[这一步有三个洞]
\begin{enumerate}[leftmargin=1.4em,itemsep=0.3em]
\item \textbf{$1/\omega$ 在原点不可积。}$\int_{0}^{1}\frac{\dd\omega}{\omega}=\infty$。所以“$\frac{1}{i\omega}$”根本不是一个局部可积函数，不能直接当分布用。
\item \textbf{方程有无穷多解。}因为 $\omega\delta=0$，所以若 $U$ 是解，$U+c\delta$ 也是解。
\item \textbf{“除”这个动作本身在分布里没有定义。}你不能对等式两边同时除以 $\omega$——$\omega$ 不是数，是乘子。
\end{enumerate}
\end{warnbox}

\subsection{关键观察：乘 $\omega$ 会抹掉 $\delta$}

\begin{equation}
\ip{\omega T}{\varphi}:=\ip{T}{\omega\varphi}.
\end{equation}

取 $T=\delta$：

\begin{equation}
\ip{\omega\delta}{\varphi}=\ip{\delta}{\omega\varphi}=0\cdot\varphi(0)=0 .
\end{equation}

所以

\begin{equation}
\boxed{\ \omega\,\delta=0\ }.\label{eq:omegadelta0}
\end{equation}

\begin{ideabox}[这就解释了“信息在哪里丢的”]
时域求导会丢掉常数：$(H+C)'=H'$。不管你加多大的常数，导数都一样。

频域乘 $\omega$ 会丢掉 $\delta$：$\omega(U+c\delta)=\omega U$。不管你加多大的 $\delta$，乘完都一样。

这两件事是同一件事的两面：\textbf{微分规则只能确定“非零频率”的部分，零频率那一项被乘 $\omega$ 消掉了。}

所以 $i\omega\hf=1$ 这句话，只能告诉我们 $\hf$ 在 $\omega\neq0$ 时的样子。至于 $\omega=0$ 上堆了多少权重，这个方程\textbf{完全不知情}。
\end{ideabox}

\section{主值：怎么给 $1/\omega$ 一个身份}

\subsection{为什么不能直接积分}

$1/\omega$ 在原点附近不是局部可积的：$\int_{0}^{1}|\frac{1}{\omega}|\dd\omega=\infty$。所以

\begin{equation}
\varphi\longmapsto\int_{\R}\frac{\varphi(\omega)}{\omega}\dd\omega
\end{equation}

对一般的 $\varphi$ 没有定义（正负两边的无穷大可能互相抵消，也可能不抵消，取决于你怎么取极限）。

\subsection{定义：对称地挖掉原点}

\begin{defn}[Cauchy 主值]\label{def:PV}
对 $\varphi\in C_c^{\infty}(\R)$，定义
\begin{equation}
\left\langle\PV\frac{1}{\omega},\varphi\right\rangle
:=\lim_{a\downarrow0}\int_{|\omega|>a}\frac{\varphi(\omega)}{\omega}\dd\omega
=\int_{0}^{\infty}\frac{\varphi(x)-\varphi(-x)}{x}\dd x .\label{eq:PVdef}
\end{equation}
\end{defn}

第二个等式值得单独解释，因为它是所有计算的起点：

\begin{equation}
\int_{|\omega|>a}\frac{\varphi(\omega)}{\omega}\dd\omega
=\int_{a}^{\infty}\frac{\varphi(x)}{x}\dd x+\int_{-\infty}^{-a}\frac{\varphi(\omega)}{\omega}\dd\omega
=\int_{a}^{\infty}\frac{\varphi(x)-\varphi(-x)}{x}\dd x .
\end{equation}

\begin{ideabox}[为什么要合并成正负一对]
因为 $\varphi(x)-\varphi(-x)$ 在 $x=0$ 处\textbf{一阶消失}：它是 $O(x)$。除以 $x$ 之后还是有界的。

而单独的 $\varphi(x)/x$ 是 $O(1/x)$，会发散。合并之后就没事了。

\textbf{主值的全部内容就是这一句：把正负两边凑成一对，让奇性自己抵消。}
\end{ideabox}

\subsection{存在性：两个估计}

\begin{prop}\label{prop:PVexists}
\eqref{eq:PVdef} 中的积分对任意 $\varphi\in\Sv$ 收敛，且 $\PV\frac1\omega$ 是 $\Sv$ 上的连续线性泛函。
\end{prop}

\begin{proof}[证明]
分两段。

\textbf{近处}（$0<x\le1$）：由中值定理，
\begin{equation}
|\varphi(x)-\varphi(-x)|\le 2x\sup_{|y|\le1}|\varphi'(y)|,
\end{equation}
所以
\begin{equation}
\left|\frac{\varphi(x)-\varphi(-x)}{x}\right|\le 2\sup_{|y|\le1}|\varphi'(y)| ,
\end{equation}
在 $(0,1]$ 上可积。

\textbf{远处}（$x\ge1$）：由 $\varphi$ 快速衰减（例如 $|\varphi(x)|\le C/x^{2}$），
\begin{equation}
\left|\frac{\varphi(x)-\varphi(-x)}{x}\right|\le\frac{C}{x^{3}},
\end{equation}
在 $[1,\infty)$ 上可积。

两段都收敛，故积分存在。连续性由上面两个估计给出的半范数控制直接得到。
\end{proof}

\subsection{对称性不能省}

这一点必须讲，否则观众会以为 $\PV$ 只是“随便忽略奇点”。

取不对称的截断：正半边挖掉 $(-\infty,-b)$ 之外的部分，负半边挖掉 $(-a,a)$，让 $a,b\to0$ 但比例不同。计算给出

\begin{equation}
\int_{a<|\omega|<1}\frac{\varphi(\omega)}{\omega}\dd\omega
=\int_{a}^{1}\frac{\varphi(x)-\varphi(-x)}{x}\dd x
+\int_{a}^{b}\frac{\varphi(-x)}{x}\dd x .
\end{equation}

第二项约为 $\varphi(0)\ln\frac{b}{a}$。于是：

\begin{itemize}[leftmargin=1.6em,itemsep=0.3em]
\item 若 $b/a\to1$（对称），第二项 $\to0$，回到主值。
\item 若 $b/a\to\lambda\neq1$，结果比主值多出 $\varphi(0)\ln\lambda$——\textbf{极限变了}。
\end{itemize}

\begin{warnbox}[结论]
$\PV$ 不是“忽略奇点”，而是“\textbf{规定了一种特定的、对称的逼近方式}”。换一种方式，得到的就不是同一个分布。
\end{warnbox}

\subsection{记号与性质}

\begin{equation}
\PV\frac{1}{i\omega}=\frac{1}{i}\,\PV\frac{1}{\omega}=-i\,\PV\frac{1}{\omega}.
\end{equation}

$\PV\frac{1}{\omega}$ 是\textbf{实、奇}的分布；乘 $-i$ 之后变成\textbf{纯虚、奇}的分布。这个对称性后面会用来独立检查 $\pi$ 这个系数。

\section{一次完整的计算：单边指数正则化}

\subsection{取外壳}

$H$ 从 $0$ 开始并停住，所以它的自然正则化是“在正半轴上加指数衰减”：

\begin{equation}
H_\ep(t)=H(t)\,e^{-\ep t},\qquad \ep>0 .
\end{equation}

它逐点趋于 $H$。

\begin{notebox}[为什么不用双边指数]
双边指数 $e^{-\ep|t|}$ 的极限是 $1$（EP01）。单边指数 $H(t)e^{-\ep t}$ 的极限是 $H$。

\textbf{衰减的“形状”决定了极限对象。}这句话是本集与上一集的分水岭。
\end{notebox}

\subsection{积分}

\begin{equation}
\F[H_\ep](\omega)=\int_{0}^{\infty}e^{-\ep t}e^{-i\omega t}\dd t
=\int_{0}^{\infty}e^{-(\ep+i\omega)t}\dd t
=\frac{1}{\ep+i\omega}.\label{eq:Heps}
\end{equation}

边界项消失是因为 $\ep>0$（上端）与积分从 $0$ 开始（下端）。

\subsection{拆成实部与虚部}

\begin{equation}
\frac{1}{\ep+i\omega}
=\frac{\ep-i\omega}{(\ep+i\omega)(\ep-i\omega)}
=\frac{\ep-i\omega}{\ep^{2}+\omega^{2}}
=\underbrace{\frac{\ep}{\ep^{2}+\omega^{2}}}_{\text{实部，偶}}
\;-\;i\underbrace{\frac{\omega}{\ep^{2}+\omega^{2}}}_{\text{虚部，奇}} .\label{eq:Hepssplit}
\end{equation}

\begin{scriptbox}[动画要点 03-D]
左面板：实部 $\dfrac{\ep}{\ep^{2}+\omega^{2}}$，一个非负的偶峰，面积 $\pi$。

右面板：虚部 $-\dfrac{\omega}{\ep^{2}+\omega^{2}}$，一条穿过原点的奇曲线。

两个面板共享 $\ep$。$\ep$ 变小时，实部变成尖峰，虚部在原点附近越来越陡。
\end{scriptbox}

\subsection{实部的极限}

注意

\begin{equation}
\frac{\ep}{\ep^{2}+\omega^{2}}=\pi\cdot\frac{1}{\pi}\frac{\ep}{\ep^{2}+\omega^{2}}=\pi K_\ep(\omega),
\end{equation}

而 $K_\ep\to\delta$（EP01 的结果，EP04 证明）。所以

\begin{equation}
\frac{\ep}{\ep^{2}+\omega^{2}}\longrightarrow\pi\delta(\omega).\label{eq:realpart}
\end{equation}

\begin{ideabox}[$\pi$ 是从哪里来的]
不是猜的，是\textbf{面积}算出来的：
\begin{equation}
\int_{\R}\frac{\ep}{\ep^{2}+\omega^{2}}\dd\omega
=\left[\arctan\frac{\omega}{\ep}\right]_{-\infty}^{\infty}=\pi .
\end{equation}
这一族函数的总面积恒为 $\pi$，所以它的极限是 $\pi\delta$。

\textbf{记住这个“面积 = 系数”的经验法则。}它会在 EP06 反复出现。
\end{ideabox}

\subsection{虚部的极限}

要证明

\begin{equation}
\frac{\omega}{\ep^{2}+\omega^{2}}\longrightarrow\PV\frac{1}{\omega}
\qquad(\ep\downarrow0).
\end{equation}

\begin{proof}[证明：合并正负半轴]
对 $\varphi\in\Sv$，
\begin{equation}
\int_{\R}\frac{\omega}{\ep^{2}+\omega^{2}}\varphi(\omega)\dd\omega
=\int_{0}^{\infty}\frac{x}{\ep^{2}+x^{2}}\big[\varphi(x)-\varphi(-x)\big]\dd x .
\end{equation}
（把 $\omega<0$ 的部分用 $\omega=-x$ 换元，与正半轴合并。）

现在验证控制收敛定理的条件。

\textbf{$0<x\le1$：}由 $x/(\ep^{2}+x^{2})\le 1/x$ 与中值估计，
\begin{equation}
\left|\frac{x}{\ep^{2}+x^{2}}\big[\varphi(x)-\varphi(-x)\big]\right|
\le\frac{1}{x}\cdot2x\sup_{|y|\le1}|\varphi'(y)|
=2\sup_{|y|\le1}|\varphi'(y)| .
\end{equation}
这是一个与 $\ep$ 无关的常数，在 $(0,1]$ 上可积。

\textbf{$x\ge1$：}由 $x/(\ep^{2}+x^{2})\le1/x$ 与 Schwartz 衰减，
\begin{equation}
\left|\frac{x}{\ep^{2}+x^{2}}\big[\varphi(x)-\varphi(-x)\big]\right|\le\frac{C}{x^{3}},
\end{equation}
在 $[1,\infty)$ 上可积，也与 $\ep$ 无关。

于是存在一个\textbf{与 $\ep$ 无关的}可积控制函数，而逐点地 $\dfrac{x}{\ep^{2}+x^{2}}\to\dfrac{1}{x}$。控制收敛给出
\begin{equation}
\int_{\R}\frac{\omega}{\ep^{2}+\omega^{2}}\varphi(\omega)\dd\omega
\longrightarrow\int_{0}^{\infty}\frac{\varphi(x)-\varphi(-x)}{x}\dd x
=\left\langle\PV\frac{1}{\omega},\varphi\right\rangle .
\end{equation}
\end{proof}

\begin{ideabox}[这个证明值得在视频里“演”一遍]
它是全系列唯一一处“标准的分析估计”，也是最值得让观众看一眼“数学长什么样”的地方。

不必讲全部细节，只需要展示那个\textbf{分段}的动作：
\begin{itemize}[leftmargin=1.4em,itemsep=0.2em]
\item 近处：用中值定理把 $\varphi(x)-\varphi(-x)$ 估成 $2x\|\varphi'\|_\infty$，与 $\ep$ 无关。
\item 远处：用快速衰减估成 $C/x^{3}$。
\item 中间：让 $\ep\to0$，$x/(\ep^{2}+x^{2})\to1/x$。
\end{itemize}
“近处靠连续性、远处靠衰减”——这句口诀在 EP04 还会出现，而且是同一个套路。
\end{ideabox}

\subsection{合起来}

把 \eqref{eq:realpart} 与上一小节的结果代回 \eqref{eq:Hepssplit}：

\begin{equation}
\F[H](\omega)=\pi\delta(\omega)-i\,\PV\frac{1}{\omega}
=\pi\delta(\omega)+\PV\frac{1}{i\omega} .\label{eq:FH}
\end{equation}

\begin{center}
\small
\renewcommand{\arraystretch}{1.4}
\begin{tabularx}{\textwidth}{@{}L{4.2cm}L{3.4cm}Y@{}}
\toprule
项 & 极限 & 来自时域的什么性质 \\
\midrule
实部 $\dfrac{\ep}{\ep^{2}+\omega^{2}}$ & $\pi\delta$ & $H$ 的“一直等于 $1$”那一半（偶部分）\\
虚部 $-\dfrac{\omega}{\ep^{2}+\omega^{2}}$ & $-i\PV\dfrac{1}{\omega}$ & $H$ 的“从 $0$ 跳到 $1$”那个奇性 \\
\bottomrule
\end{tabularx}
\end{center}

\section{用对称性独立检查 $\pi$}

这一步很重要，因为它是\textbf{不看计算}就能验证系数的方法。

\subsection{时域的对称关系}

作为分布，

\begin{equation}
H(t)+H(-t)=1 .\label{eq:Hsym}
\end{equation}

（$t\neq0$ 时逐点成立；$t=0$ 只影响一个点，无妨。）

\subsection{频域的对应}

反射性质：$\F[f(-t)](\omega)=\fh(-\omega)$。所以对 \eqref{eq:Hsym} 两边做变换：

\begin{equation}
\hf(\omega)+\hf(-\omega)=2\pi\delta(\omega).\label{eq:Hfreqsym}
\end{equation}

现在把候选答案写成

\begin{equation}
\hf(\omega)=-i\,\PV\frac{1}{\omega}+c\,\delta(\omega)
\end{equation}

（这里 $c$ 未知）。代进 \eqref{eq:Hfreqsym}：

\begin{itemize}[leftmargin=1.6em,itemsep=0.3em]
\item 奇的部分：$\PV\frac{1}{\omega}$ 是奇的，所以 $-i\PV\frac1\omega$ 与它的反射\textbf{正好抵消}。
\item 偶的部分：$\delta(-\omega)=\delta(\omega)$，所以 $c\delta$ 变成 $2c\delta$。
\end{itemize}

于是

\begin{equation}
2c\,\delta=2\pi\,\delta
\qquad\Longrightarrow\qquad
c=\pi .
\end{equation}

\begin{ideabox}[这一步为什么值得单独讲]
因为它把“$\pi$”从一个算出来的数字，变成一个\textbf{可以被逻辑逼出来}的必然结果。

观众会觉得：原来不是“算错了才凑出 $\pi$”，而是“对称性要求它必须是 $\pi$”。

这种“用另一条路验证”的段落，在全系列里应该至少出现三次（EP03 的 $\pi$、EP06 的 $2\pi$、EP06 的 $F^2=2\pi$ 反射）。
\end{ideabox}

\subsection{顺带得到符号函数的变换}

由 $H=\frac12+\frac12\sgn$，

\begin{equation}
\F[\sgn t]=2\Big(\hf-\tfrac12\F[1]\Big)=2\big(\pi\delta-i\PV\tfrac1\omega-\pi\delta\big)
=-2i\,\PV\frac{1}{\omega} .\label{eq:sgn}
\end{equation}

\section{为什么歧义恰好只有 $c\delta$}

上一节我们“假设”了答案形如 $-i\PV\frac1\omega+c\delta$。凭什么？

\begin{prop}\label{prop:ambiguity}
若 $T\in\Sv'$ 满足 $\omega T=0$，则 $T=c\,\delta$ 对某个常数 $c$。
\end{prop}

\begin{proof}[证明：把 $\varphi$ 拆成“常数部分 + 乘 $\omega$ 的部分”]
取 $\chi\in C_c^{\infty}$ 满足 $\chi\equiv1$ 于 $|x|\le1$。对任意 $\varphi\in\Sv$，令
\begin{equation}
\psi(\omega):=\frac{\varphi(\omega)-\varphi(0)\chi(\omega)}{\omega}.
\end{equation}

验证 $\psi\in\Sv$：
\begin{itemize}[leftmargin=1.4em,itemsep=0.2em]
\item 分子在 $\omega=0$ 处为零，且一阶消失（因为 $\chi(0)=1$，$\chi'(0)=0$ 可以取到），所以商可以光滑延拓到 $0$。
\item 远处 $\chi\equiv0$，$\psi=\varphi/\omega$，除以 $\omega$ 只会让衰减更快。
\end{itemize}

于是

\begin{equation}
\varphi(\omega)=\varphi(0)\chi(\omega)+\omega\,\psi(\omega),
\end{equation}

作用 $T$：

\begin{equation}
\ip{T}{\varphi}=\varphi(0)\ip{T}{\chi}+\ip{\omega T}{\psi}
=\varphi(0)\ip{T}{\chi}+0
=c\,\varphi(0)=c\ip{\delta}{\varphi},
\end{equation}

其中 $c:=\ip{T}{\chi}$。所以 $T=c\delta$。
\end{proof}

\begin{warnbox}[为什么不能再加 $\delta'$]
因为 $\omega\delta'=-\delta\neq0$。加一个 $\delta'$ 会破坏 $\omega T=0$ 这个条件。

这一句要讲：\textbf{能满足“乘 $\omega$ 之后为零”的对象，最小的一族就是 $c\delta$；$\delta'$ 太重了，乘 $\omega$ 消不掉。}（顺带由 $x\delta^{(n)}=-n\delta^{(n-1)}$，见 EP05。）
\end{warnbox}

\section{边界值公式（拓展，可放补充材料）}

式 \eqref{eq:Heps} 还给出一个在数学物理里极常用的写法：

\begin{equation}
\frac{1}{x\pm i0}:=\lim_{\ep\downarrow0}\frac{1}{x\pm i\ep}
=\PV\frac{1}{x}\mp i\pi\,\delta(x).\label{eq:boundary}
\end{equation}

这里 $\pm i0$ 是“$\ep\downarrow0$ 的分布极限”的记号，\textbf{不是}把复数零代进分母。

\begin{notebox}[符号核对：为什么是 $-\dfrac{i}{\omega-i0}$]
\begin{equation}
\frac{1}{\ep+i\omega}
=\frac{-i}{\omega-i\ep}
\;\longrightarrow\;
\frac{-i}{\omega-i0}
=-i\Big(\PV\frac{1}{\omega}+i\pi\delta\Big)
=\pi\delta-i\,\PV\frac{1}{\omega}.
\end{equation}
第一行用了 $\ep+i\omega=i(\omega-i\ep)$；最后一步用了 \eqref{eq:boundary} 中取“$-$”号的那一支。

\textbf{这个负号非常容易错。}如果你算出的结果是 $-\pi\delta$，多半就是在这里把 $i$ 与 $-i$ 弄反了。
\end{notebox}

\section{数值与动画}

\begin{center}
\small
\renewcommand{\arraystretch}{1.35}
\begin{tabularx}{\textwidth}{@{}L{1.4cm}L{3.6cm}L{3.8cm}L{2.6cm}Y@{}}
\toprule
$\ep$ & 实部峰值 $1/\ep$ & 实部面积 & 虚部极值 $\mp1/(2\ep)$ & 备注 \\
\midrule
$1$ & $1$ & $\pi$ & $\mp0.5$ & 纵轴 $\pm1$ \\
$0.5$ & $2$ & $\pi$ & $\mp1$ & 纵轴 $\pm2$ \\
$0.2$ & $5$ & $\pi$ & $\mp2.5$ & 纵轴 $\pm5$ \\
$0.1$ & $10$ & $\pi$ & $\mp5$ & 采样加密 \\
\bottomrule
\end{tabularx}
\end{center}

\begin{figure}[htbp]
\centering
\begin{tikzpicture}
\begin{axis}[width=0.46\textwidth,height=4.4cm,
  domain=-5:5,samples=400,ymin=0,ymax=2.2,
  xlabel={$\omega$},ylabel={实部},
  title={\small 实部：面积恒为 $\pi$},
  ticklabel style={font=\scriptsize},label style={font=\small}]
\addplot[thick,mainblue] {0.5/(0.25+x^2)};
\end{axis}
\end{tikzpicture}\hfill
\begin{tikzpicture}
\begin{axis}[width=0.46\textwidth,height=4.4cm,
  domain=-5:5,samples=400,ymin=-1.2,ymax=1.2,
  xlabel={$\omega$},ylabel={虚部},
  title={\small 虚部：奇的，面积为零},
  ticklabel style={font=\scriptsize},label style={font=\small},
  axis lines=middle]
\addplot[thick,warnred] {-x/(0.25+x^2)};
\end{axis}
\end{tikzpicture}
\caption{$\ep=0.5$ 时 $\F[H_\ep]$ 的实部与虚部。实部贡献 $\pi\delta$，虚部贡献 $-i\PV(1/\omega)$。}
\end{figure}

\begin{warnbox}[主值动画的两条纪律]
\begin{enumerate}[leftmargin=1.4em,itemsep=0.3em]
\item 必须\textbf{对称}地删去原点两侧的区间。画成“只挖掉正半边”或“两边挖得不一样宽”，观众学到的就是错的。
\item 不要把正负两侧的积分分别标成有限数再加起来。两侧各自都是发散的；有意义的是\textbf{合并之后}的极限。
\end{enumerate}
\end{warnbox}

\section{讲解禁区（EP03）}

\begin{warnbox}[逐条对照]
\begin{enumerate}[leftmargin=1.6em,itemsep=0.28em]
\item 不说“$H(0)$ 取哪个值决定了 $\delta$ 的系数”。
\item 不说“$1/(i\omega)$ 就是答案”，必须先处理原点奇性。
\item 不说“$\PV$ 就是忽略奇点”；要说“对称地删去原点邻域”。
\item 不把正负两个发散积分分别求值后相加。
\item 不说“$i\omega\widehat H=1$ 唯一确定 $\widehat H$”。
\item 不说“加了常数 $C$ 的阶跃频谱还是同一个”（时域平移/加常数会改零频项）。
\item 不把 $\PV(1/(i\omega))$ 的符号写成 $+i\PV(1/\omega)$（正确是 $-i$）。
\item 不说“$\pi$ 是凑出来的”；给出面积论证或对称性论证。
\item 不说“$H'=\delta$ 说明 $H$ 在通常意义下可导”。
\item 不在没有声明广义意义的情况下写 $\F[H]$ 的等式。
\item 不用“跳变处无穷大”这种描述代替“跳跃强度为 $1$ 的 $\delta$ 贡献”。
\end{enumerate}
\end{warnbox}

\section{与前后集的接口}

\begin{enumerate}[leftmargin=2em,itemsep=0.45em]
\item \textbf{回头看}：这一集把 EP01 的 $K_\ep$、EP02 的 $H'=\delta$ 全部用上了。它同时暴露了前两集一直含糊的那件事——“求极限”到底是哪种极限。
\item \textbf{往前看}：这里的每一句“趋于”都是\textbf{对测试函数的作用}在趋于。这已经是在用分布收敛了，只是还没给它名字。EP04 停下来问名字，EP05 给它定义。
\end{enumerate}

\begin{scriptbox}[最后 30 秒的旁白骨架]
“到这里，我们手里有三个结果，每一个都用到了一种‘趋于’。

常数的频谱趋于一个尖峰；阶跃的频谱趋于一个尖峰加一个奇怪的 $1/\omega$；$\delta$ 的导数趋于两叶的东西。

可我们一直在说‘趋于’，却从来没说清楚是哪一种趋于。

下一集，我们停下来，只问这一个问题。”
\end{scriptbox}


\section{本集回顾卡}

\begin{center}
\small
\renewcommand{\arraystretch}{1.4}
\begin{tabularx}{\textwidth}{@{}L{3.4cm}Y@{}}
\toprule
项目 & 内容 \\
\midrule
主结果 & $\F[H]=\pi\delta(\omega)+\PV\dfrac{1}{i\omega}$ \\
关键歧义 & $\omega\delta=0$，所以 $i\omega\widehat H=1$ 的解不唯一 \\
主值 & 对称挖去原点邻域；不对称截断会差 $\varphi(0)\ln(b/a)$ \\
正则化 & $\dfrac{1}{\ep+i\omega}=\dfrac{\ep}{\ep^{2}+\omega^{2}}-i\dfrac{\omega}{\ep^{2}+\omega^{2}}$ \\
安全话术 & “微分规则只确定非零频率那一部分” \\
危险话术 & “$\PV$ 就是忽略奇点”“除以 $i\omega$ 就完了” \\
下一集的钩子 & 我们说的这些“趋于”，究竟是哪一种？ \\
\bottomrule
\end{tabularx}
\end{center}

\begin{quizbox}[EP03 自测]
\begin{enumerate}[leftmargin=1.6em,itemsep=0.35em]
\item 把 $H(0)$ 改成 $1/2$，证明它的分布导数仍然是 $\delta$。
\item 验证 $i\omega\widehat H=1$：用 \eqref{eq:FH} 代进去，说明 $\omega\delta=0$ 与 $\omega\PV(1/\omega)=1$ 各起什么作用。
\item 为什么 $i\omega\widehat H=1$ 不能唯一确定 $\widehat H$？给出两个不同的解。
\item 用 \eqref{eq:PVdef} 计算 $\left\langle\PV\frac1\omega,\varphi\right\rangle$，其中 $\varphi(\omega)=(1+\omega)e^{-\omega^{2}}$。（提示：先算 $\varphi(x)-\varphi(-x)$；答案 $\sqrt{\pi}$。）
\item 写出不对称截断的极限，并说明它与 $\PV$ 相差 $\varphi(0)\ln(b/a)$。
\item 由 $\F[H]$ 推出 $\F[\sgn t]$，检查它是纯虚、奇。
\item 用 \eqref{eq:boundary} 与 $\dfrac{1}{\ep+i\omega}=\dfrac{-i}{\omega-i\ep}$ 重新验证 $\F[H]$ 的符号。
\end{enumerate}
\end{quizbox}
}{}
\IfFileExists{ep04.tex}{\chapter{EP04｜点值都趋于零，积分为什么还在}

\begin{ideabox}[这一集的一句话]
\textbf{一个有面积的尖峰变成 $\delta$，不是因为它变“小”了，而是因为它的质量向一个点集中了。把这两件事分开，就必须先分清“逐点收敛”和“积分意义下的收敛”。}
\end{ideabox}

\section{观众的问题链}

这一集是全系列唯一一次\textbf{主动停下来}，不推进结论，只审问前面用过的词。它也是从“算结果”转向“问意义”的转折点。

\begin{center}
\small
\renewcommand{\arraystretch}{1.32}
\begin{tabularx}{\textwidth}{@{}p{0.8cm}X p{4.4cm}@{}}
\toprule
层 & 问题 & 观众此刻的心理 \\
\midrule
1 & 回放那条核：固定一个非零的点，它的值确实越来越小。 & 那它应该趋于零啊。\\
2 & 可是它的积分一直是 $1$。 & 怎么可能？\\
3 & 因为质量全挤到原点附近去了。 & 那到底算趋于零还是不趋于零？\\
4 & 取决于你在问“点值”还是“面积”。 & 这两种“趋于”是不一样的？\\
5 & 不一样。而且我们以前学过的“把极限搬进积分”，需要许可证。 & 什么许可证？\\
6 & 控制收敛。而这条核族根本拿不到这张许可证。 & 那凭什么说它趋于 $\delta$？\\
\bottomrule
\end{tabularx}
\end{center}

\begin{center}
\small
\renewcommand{\arraystretch}{1.32}
\begin{tabularx}{\textwidth}{@{}p{3.2cm}Y@{}}
\toprule
观众带着什么进来 & 观众带着什么出去 \\
\midrule
“极限就是逐点看” & “同一个函数列可以在一种意义下趋于零、在另一种意义下趋于 $\delta$。”\\
“积分的极限等于极限的积分” & “这句话需要公共可积控制；没有它，一切皆可能。”\\
“$\delta$ 是某种特别尖的形状” & “任何满足‘质量守恒 + 尾部集中’的核都趋于 $\delta$，形状无所谓。”\\
\bottomrule
\end{tabularx}
\end{center}

\section{重看那个尖峰：两个读数，两种结论}

回到第一集的核：

\begin{equation}
K_\ep(x)=\frac{1}{\pi}\frac{\ep}{x^{2}+\ep^{2}},
\qquad
\int_{\R}K_\ep\dd x=1 .
\end{equation}

现在做两件不同的事。

\subsection{读点值}

固定 $x\neq0$：

\begin{equation}
K_\ep(x)=\frac{1}{\pi}\frac{\ep}{x^{2}+\ep^{2}}\;\longrightarrow\;0
\qquad(\ep\downarrow0).
\end{equation}

固定 $x=0$：

\begin{equation}
K_\ep(0)=\frac{1}{\pi\ep}\;\longrightarrow\;+\infty .
\end{equation}

所以：\textbf{除原点以外，每一点都趋于零；但原点本身趋于无穷。}

\subsection{读面积}

\begin{equation}
\int_{\R}K_\ep\dd x=1
\qquad\text{对每个 }\ep>0 .
\end{equation}

而且对每个固定的 $a>0$，

\begin{equation}
\int_{-a}^{a}K_\ep\dd x=\frac{2}{\pi}\arctan\frac{a}{\ep}\;\longrightarrow\;1 .
\end{equation}

所以：\textbf{全部质量都挤进了任意小的原点邻域，一点都没漏出去。}

\begin{ideabox}[这两个读数的冲突，就是这一集要解开的结]
“几乎每一点都趋于零”说的是函数值。

“总质量始终为 $1$ 且集中到一点”说的是积分。

这两句话\textbf{都对}，而且并不矛盾。矛盾只出现在你把它们混成一句“这个函数列趋于零”的时候。

\begin{center}
\textbf{正确的说法是：$K_\ep\to\delta$，在“对测试函数的作用”这个意义下。}
\end{center}
\end{ideabox}

\begin{warnbox}[一个必须避免的措辞]
不要说“$K_\ep$ 除了原点以外处处趋于零，所以极限是零，但面积不为零，所以出现了矛盾”。

这里没有矛盾，只有\textbf{用错了收敛方式}。原点是例外，而且正是它承载了全部质量；把它说成“除了一个点以外处处为零”，等于把这一集的主角删掉了。
\end{warnbox}

\section{三种收敛：它们分别在检查什么}

\begin{center}
\small
\renewcommand{\arraystretch}{1.5}
\begin{tabularx}{\textwidth}{@{}L{3.4cm}L{4.4cm}Y@{}}
\toprule
收敛方式 & 定义 & 它在检查什么 \\
\midrule
逐点 & 对每个固定的 $x$，$f_n(x)\to f(x)$ & 每个点各自的行为；最弱 \\
一致 & $\sup_x|f_n(x)-f(x)|\to0$ & 整体最大偏差；中等 \\
分布 & 对每个固定的 $\varphi\in\D$，$\ip{T_n}{\varphi}\to\ip{T}{\varphi}$ & 对每个探针的\textbf{读数}；与点值无关 \\
\bottomrule
\end{tabularx}
\end{center}

\subsection{反例一：一致收敛也不保证积分收敛}

\begin{equation}
f_n(x)=\frac{1}{n}\,\mathbf{1}_{[0,n]}(x).
\end{equation}

$\sup_x|f_n(x)|=\dfrac1n\to0$，所以 $f_n$\textbf{一致}趋于零。但

\begin{equation}
\int_{\R}f_n(x)\dd x=\frac{1}{n}\cdot n=1
\qquad\text{对每个 }n .
\end{equation}

\begin{ideabox}[这个反例在说什么]
一致收敛控制的是“高度”，不控制“长度”。$f_n$ 越来越矮，同时越来越长，面积恰好保持不变。

\textbf{在无界区间上，“一致收敛”不是把极限搬进积分的许可证。}这一点很多教材的入门部分会含糊过去，而它恰恰是本系列的关键。
\end{ideabox}

\subsection{反例二：逐点收敛更弱，而且可能收敛到 $\delta$}

\begin{equation}
g_n(x)=n\,\mathbf{1}_{(0,1/n)}(x).
\end{equation}

\begin{itemize}[leftmargin=1.6em,itemsep=0.3em]
\item 对每个固定的 $x\neq0$：只要 $n>1/x$，就有 $g_n(x)=0$，所以 $g_n(x)\to0$。
\item $\int_{\R}g_n=1$ 对每个 $n$。
\item 对任意 $\varphi\in C_c^{\infty}$（不必连续也可，只要在 $0$ 连续），
\begin{equation}
\ip{g_n}{\varphi}=n\int_{0}^{1/n}\varphi(x)\dd x\;\longrightarrow\;\varphi(0).
\end{equation}
\end{itemize}

所以

\begin{equation}
g_n\to0\ \text{（逐点）},
\qquad
g_n\to\delta\ \text{（分布意义）}.
\end{equation}

\begin{warnbox}[这两句话必须一起说]
如果只说“$g_n$ 逐点趋于零”，观众会以为它就是个零。

如果只说“$g_n$ 趋于 $\delta$”，观众会以为 $\delta$ 与零差不多。

\textbf{两句话同时为真，正是因为它们说的是两件不同的事。}这就是整集的核心，也是最难讲清的一处。
\end{warnbox}

\subsection{反例三：$K_\ep$ 连“处处逐点”都不是}

$K_\ep(0)\to\infty$，所以它不满足“每点趋于零”。它属于“\textbf{除原点外}逐点趋于零”。

这个细节决定了一个措辞：标题和旁白用“几乎每个点”，或直接说“除了原点”，\textbf{不要说“处处”}。

\begin{center}
\small
\renewcommand{\arraystretch}{1.4}
\begin{tabularx}{\textwidth}{@{}L{3.6cm}L{2.6cm}L{2.6cm}Y@{}}
\toprule
函数列 & 逐点 & 积分 & 备注 \\
\midrule
$\frac1n\mathbf{1}_{[0,n]}$ & $\to0$（一致） & $=1$ & 高度减、长度增 \\
$n\mathbf{1}_{(0,1/n)}$ & $\to0$ & $=1$ & 分布意义 $\to\delta$ \\
$K_\ep$ & $\to0$（除 $0$） & $=1$ & 原点趋于 $\infty$ \\
\bottomrule
\end{tabularx}
\end{center}

\section{控制收敛：许可证到底要什么}

\begin{thm}[控制收敛，本系列使用的版本]\label{thm:DCT}
设 $f_n\to f$ 几乎处处，且存在\textbf{同一个}可积函数 $g\in\Lone$，使
\begin{equation}
|f_n(x)|\le g(x)\qquad\text{几乎处处，对所有的 }n .
\end{equation}
则
\begin{equation}
\int_{\R}f_n\;\longrightarrow\;\int_{\R}f .
\end{equation}
\end{thm}

\subsection{重点在“同一个”三个字}

“每一项都小于某个可积函数”是不够的：那个函数必须对\textbf{所有} $n$ 同时有效。

\begin{prop}\label{prop:noControl}
核族 $\{K_\ep\}_{0<\ep<1}$ 不存在公共的 $\Lone$ 控制函数。
\end{prop}

\begin{proof}[反证]
假设存在 $g\in\Lone$ 使 $0\le K_\ep\le g$ 对一切 $\ep$ 成立。

由于 $K_\ep(x)\to0$ 对每个 $x\neq0$，由定理 \ref{thm:DCT}，
\begin{equation}
\int_{\R}K_\ep\dd x\longrightarrow\int_{\R}0\,\dd x=0 .
\end{equation}

但左边恒等于 $1$，矛盾。
\end{proof}

\begin{ideabox}[这个反证法本身就是一句很好的旁白]
“如果你能找到那样一个控制函数，那么控制收敛定理就会告诉你积分趋于零。可是积分明明是 $1$。所以，那样的控制函数根本不存在。”

一句话，把“为什么不能把极限搬进去”讲完了，而且没有用任何术语。
\end{ideabox}

\subsection{那 $e^{-\ep|t|}$ 呢？它明明被 $1$ 控制}

$|e^{-\ep|t|}|\le1$ 确实成立。但 $1\notin\Lone(\R)$，所以定理 \ref{thm:DCT} 用不上。

这就是 EP01 的困难在分析层面的精确位置：\textbf{你有控制，但控制函数不可积。}

\begin{ideabox}[测试函数是怎么救场的]
乘上一个 $\varphi\in\Sv$ 之后，
\begin{equation}
\big|e^{-\ep|t|}\varphi(t)\big|\le|\varphi(t)|\in\Lone .
\end{equation}
现在控制函数可积了，控制收敛立刻给出
\begin{equation}
\ip{e^{-\ep|t|}}{\varphi}=\int_{\R}e^{-\ep|t|}\varphi(t)\dd t
\;\longrightarrow\;\int_{\R}\varphi(t)\dd t=\ip{1}{\varphi}.
\end{equation}
也就是 $e^{-\ep|t|}\to1$ 于 $\Sv'$。

\textbf{“引入测试函数”不是形式主义，它是把不可积的控制换成可积的控制。}这句话要在 EP05 的开场再说一次。
\end{ideabox}

\section{近似恒等核：把“集中”写成条件}

\subsection{定义}

\begin{defn}[近似恒等核]\label{def:approx-id}
设 $k\in\Lone$，$\int_{\R}k=1$。令
\begin{equation}
k_\ep(x)=\frac{1}{\ep}\,k\!\left(\frac{x}{\ep}\right),
\qquad \ep>0 .\label{eq:dilate}
\end{equation}
称 $\{k_\ep\}$ 为近似恒等核（\textbf{approximate identity}）。
\end{defn}

两条性质立刻成立：

\begin{enumerate}[leftmargin=2em,itemsep=0.45em]
\item \textbf{质量守恒}：
\begin{equation}
\int_{\R}k_\ep(x)\dd x=\int_{\R}k(u)\dd u=1 .
\end{equation}
\item \textbf{尾部集中}：对每个 $a>0$，
\begin{equation}
\int_{|x|>a}k_\ep(x)\dd x=\int_{|u|>a/\ep}k(u)\dd u\;\longrightarrow\;0
\qquad(\ep\downarrow0).
\end{equation}
\end{enumerate}

\begin{ideabox}[这两条各自管什么]
\textbf{质量守恒}保证总权重不变，所以极限不会“变轻”。

\textbf{尾部集中}保证质量不会留在外面，所以极限不会“变宽”。

两者缺一不可。EP01 的 $K_\ep$、EP03 的 $\pi K_\ep$、EP06 的高斯核和 sinc 核，都是这两条的不同实现。
\end{ideabox}

\subsection{取样极限：定理与证明}

\begin{thm}[取样极限]\label{thm:sampling}
设 $k_\ep$ 如 \eqref{eq:dilate}，$\varphi$ 有界且在 $x=0$ 连续。则
\begin{equation}
\int_{\R}k_\ep(x)\varphi(x)\dd x\;\longrightarrow\;\varphi(0)
\qquad(\ep\downarrow0).
\end{equation}
\end{thm}

\begin{proof}
由总质量为 $1$，误差可以写成
\begin{equation}
\int_{\R}k_\ep(x)\varphi(x)\dd x-\varphi(0)
=\int_{\R}k_\ep(x)\big[\varphi(x)-\varphi(0)\big]\dd x .
\end{equation}

给定 $\eta>0$。由 $\varphi$ 在 $0$ 连续，存在 $a>0$ 使得 $|x|<a$ 时
\begin{equation}
|\varphi(x)-\varphi(0)|<\eta .
\end{equation}

把积分分成内外两部分：
\begin{align}
\left|\int_{\R}k_\ep(x)\big[\varphi(x)-\varphi(0)\big]\dd x\right|
&\le\int_{|x|<a}k_\ep(x)\,\eta\dd x
+\int_{|x|\ge a}k_\ep(x)\cdot2\|\varphi\|_\infty\dd x \notag\\
&\le \eta+2\|\varphi\|_\infty\int_{|x|\ge a}k_\ep(x)\dd x .
\end{align}

第一项用了 $\int_{|x|<a}k_\ep\le1$。

现在\textbf{先固定 $a$}，让 $\ep\downarrow0$：由尾部集中，第二项 $\to0$。于是
\begin{equation}
\limsup_{\ep\to0}\left|\int_{\R}k_\ep\varphi-\varphi(0)\right|\le\eta .
\end{equation}

最后让 $\eta\downarrow0$，得结论。
\end{proof}

\begin{ideabox}[这是全系列最值得掌握的一个估计]
它的结构可以背下来，因为后面还会用：

\begin{center}
\textbf{近处靠连续性}（$|\varphi(x)-\varphi(0)|<\eta$）\quad $+$\quad \textbf{远处靠尾部质量消失}（$\int_{|x|>a}k_\ep\to0$）。
\end{center}

“先固定 $a$，再让 $\ep\to0$，最后让 $\eta\to0$”这个三步顺序，是分析里非常典型的“两个小量不能同时放”的处理方式。讲的时候可以强调：\textbf{不要一次让两个东西同时趋于零。}
\end{ideabox}

\subsection{Cauchy 核的另一种更快的证明}

对 $K_\ep$，换元 $x=\ep u$：

\begin{equation}
\int_{\R}K_\ep(x)\varphi(x)\dd x
=\frac{1}{\pi}\int_{\R}\frac{\ep}{\ep^{2}+\ep^{2}u^{2}}\,\varphi(\ep u)\,\ep\dd u
=\frac{1}{\pi}\int_{\R}\frac{\varphi(\ep u)}{1+u^{2}}\dd u .
\end{equation}

此时被积函数被

\begin{equation}
\frac{\|\varphi\|_\infty}{\pi(1+u^{2})}
\end{equation}

控制，而\textbf{这个控制函数与 $\ep$ 无关且可积}。$\varphi(\ep u)\to\varphi(0)$ 对每个 $u$。于是控制收敛直接给出

\begin{equation}
\frac{1}{\pi}\int_{\R}\frac{\varphi(\ep u)}{1+u^{2}}\dd u
\longrightarrow\frac{\varphi(0)}{\pi}\int_{\R}\frac{\dd u}{1+u^{2}}
=\frac{\varphi(0)}{\pi}\cdot\pi=\varphi(0).
\end{equation}

\begin{ideabox}[换元做了什么]
它把“\textbf{不断变窄变高的核}”换成了“\textbf{固定不变的权重} $\frac{1}{\pi(1+u^{2})}$”，把变化全部转移到探针 $\varphi(\ep u)$ 身上。

于是“找不到公共控制”的困难消失了：现在控制函数就在眼前。

注意：我们没有对原来的 $K_\ep$ 声称存在公共控制（命题 \ref{prop:noControl} 说那不可能）。\textbf{我们换的是积分变量，不是收敛方式。}
\end{ideabox}

\section{为什么叫“近似恒等”}

\subsection{卷积}

定义卷积

\begin{equation}
(k*f)(x)=\int_{\R}k(x-y)f(y)\dd y=\int_{\R}k(y)f(x-y)\dd y .
\end{equation}

\begin{thm}[近似恒等]\label{thm:approx-id}
若 $f$ 有界且在 $x$ 连续，则
\begin{equation}
(k_\ep*f)(x)\longrightarrow f(x)\qquad(\ep\downarrow0).
\end{equation}
若 $f$ 有界且一致连续，则收敛是一致的。若 $f\in L^{p}$（$1\le p<\infty$），则 $k_\ep*f\to f$ 于 $L^{p}$。
\end{thm}

证明与定理 \ref{thm:sampling} 同一套路：把 $k_\ep(y)[f(x-y)-f(x)]$ 分内外估计。

\subsection{$\delta$ 就是卷积的单位元}

把 $k_\ep$ 换成 $\delta$：形式上有

\begin{equation}
(\delta*f)(x)=\int_{\R}\delta(y)f(x-y)\dd y=f(x).
\end{equation}

这正是“恒等”二字的来源：\textbf{和 $\delta$ 卷积，等于什么都不做。}

\begin{ideabox}[所以 $\delta$ 有两张脸]
\begin{itemize}[leftmargin=1.4em,itemsep=0.2em]
\item \textbf{频率的脸}：$\F[\delta]=1$，它在频域上“什么都不改变”。
\item \textbf{卷积的脸}：$\delta*f=f$，它在时域上“什么都不改变”。
\end{itemize}
这两张脸在 EP06 会合起来：$\delta$ 是傅里叶变换意义下的“单位”。
\end{ideabox}

\section{换一种核，为什么还是 $\delta$}

\subsection{三种核，同一个极限}

\begin{center}
\small
\renewcommand{\arraystretch}{1.45}
\begin{tabularx}{\textwidth}{@{}L{3.4cm}L{4.0cm}L{2.0cm}Y@{}}
\toprule
核 & 表达式 & 面积 & 与 $K_\ep$ 的关系 \\
\midrule
Cauchy（EP01） & $\dfrac{1}{\pi}\dfrac{\ep}{x^{2}+\ep^{2}}$ & $1$ & 基准 \\
高斯 & $\dfrac{1}{\sqrt{2\pi}\,\ep}e^{-x^{2}/(2\ep^{2})}$ & $1$ & 收敛更快，尾部指数衰减 \\
矩形 & $\dfrac{1}{2\ep}\mathbf{1}_{[-\ep,\ep]}(x)$ & $1$ & 不连续，但照样可以 \\
\bottomrule
\end{tabularx}
\end{center}

三者都满足定义 \ref{def:approx-id} 的两条性质，所以都趋于 $\delta$：

\begin{equation}
K_\ep\to\delta,\qquad
\rho_\ep\to\delta,\qquad
r_\ep\to\delta
\qquad(\text{于 }\Sv').
\end{equation}

\begin{figure}[htbp]
\centering
\begin{tikzpicture}
\begin{axis}[width=0.5\textwidth,height=4.6cm,
  domain=-1.5:1.5,samples=400,ymin=0,ymax=3.2,
  xlabel={$x$},ylabel={核},
  title={\small 三种核，$\ep=0.3$},
  legend style={font=\scriptsize,at={(0.97,0.97)},anchor=north east,draw=none},
  ticklabel style={font=\scriptsize},label style={font=\small}]
\addplot[thick,mainblue] {0.3/(pi*(x^2+0.09))};
\addplot[thick,warnred,dashed] {1/(sqrt(2*pi)*0.3)*exp(-x^2/(2*0.09))};
\addplot[thick,okgreen,dotted,domain=-0.3:0.3,samples=2] {1/(2*0.3)};
\legend{Cauchy,高斯,矩形}
\end{axis}
\end{tikzpicture}
\caption{三种形状完全不同、总面积都为 $1$、都在 $\ep\to0$ 时趋于同一个 $\delta$。\textbf{形状不是答案，作用是答案。}}
\end{figure}

\subsection{但只有面积是 $1$ 远远不够}

\begin{prop}[质量跑掉的反例]\label{prop:escaping}
令
\begin{equation}
k_n(x)=\frac{1}{2}\,\mathbf{1}_{[n-1,\,n+1]}(x).
\end{equation}
则 $\int_{\R}k_n=1$ 对每个 $n$，但
\begin{equation}
\ip{k_n}{\varphi}\longrightarrow0
\qquad\text{对每个 }\varphi\in\Sv,
\end{equation}
即 $k_n\to0$ 于 $\Sv'$，而不是 $\delta$。
\end{prop}

\begin{proof}
$\varphi\in\Sv$ 时，$\sup_{x\in[n-1,n+1]}|\varphi(x)|\le\dfrac{C}{n^{2}}$（快速衰减），于是
\begin{equation}
|\ip{k_n}{\varphi}|\le\frac{1}{2}\int_{n-1}^{n+1}|\varphi(x)|\dd x\le\frac{C}{n^{2}}\longrightarrow0 .
\end{equation}
\end{proof}

\begin{ideabox}[这个例子为什么让人不安，而且很重要]
它的\textbf{总面积}一直是 $1$，但\textbf{质量跑到了无穷远}。所以“总面积守恒”这个条件，在这里完全不起作用。

为什么？因为在 $\Sv'$ 里，能用来“测试”的函数都是快速衰减的。“总面积”相当于用常数 $1$ 去做测试函数——而 $1\notin\Sv$。

\textbf{结论：在 $\Sv'$ 的框架里，“总质量”根本不是一个连续的测试量。}你可以有质量，但如果它跑到无穷远，任何测试函数都摸不到它。

这句话在 EP06 会再次出现，并且是理解“为什么不能对任意 $\D'$ 做变换”的钥匙之一。
\end{ideabox}

\subsection{更一般的充分条件}

不带非负性，也可以做。只要：

\begin{enumerate}[leftmargin=2em,itemsep=0.4em]
\item $\displaystyle\int_{\R}k_\ep\dd x=1$（对每个 $\ep$）；
\item $\displaystyle\sup_{\ep}\|k_\ep\|_1<\infty$（一致有界）；
\item 对每个 $a>0$，$\displaystyle\int_{|x|>a}|k_\ep(x)|\dd x\to0$（尾部集中）。
\end{enumerate}

则 $\ip{k_\ep}{\varphi}\to\varphi(0)$。证明只需把定理 \ref{thm:sampling} 里的 $k_\ep$ 换成 $|k_\ep|$，并把 $\eta$ 项乘上一致界。

\begin{warnbox}[非负性是可选项，“集中”不是]
很多教材为了简化，只讲非负核。非负性确实让估计更干净（$\int_{|x|<a}k_\ep\le1$ 直接可用），但它不是本质条件。

本质条件是：\textbf{质量不能跑，而且必须往正确的点集中。}EP06 的 sinc 核就是“不非负、不逐点收敛，但依然分布收敛”的例子，那时候这个一般版本就派上用场了。
\end{warnbox}

\section{讲解设计与数值}

\begin{scriptbox}[动画要点 04-A / 04-B]
\textbf{04-A}：同一幅图里并排放两个读数。左边标“$K_\ep(x_0)$ 的数值”，右边标“全实轴面积”。前者趋零，后者恒为 $1$。让观众同时看到两个数字。

\textbf{04-B}：把数轴分成“原点附近 $|x|<a$”和“外面 $|x|\ge a$”两段，用两种颜色表示质量。内部的质量条越来越满（$\to1$），外部的越来越空（$\to0$）。

\textbf{04-F（可选反例）}：把核整体向右平移，做成 $k_n$ 那样跑到无穷远。面积标签仍是 $1$，但读数趋于零。这一幕能把“面积守恒不等于收敛”钉死。
\end{scriptbox}

\begin{center}
\small
\renewcommand{\arraystretch}{1.35}
\begin{tabularx}{\textwidth}{@{}L{1.3cm}L{3.2cm}L{3.2cm}L{3.0cm}Y@{}}
\toprule
$\ep$ & $\int_{|x|<\ep}K_\ep$ & $\int_{|x|<3\ep}K_\ep$ & 尾部 $\int_{|x|>10\ep}$ & 画面 \\
\midrule
$1$ & $0.50$ & $0.80$ & $0.063$ & 明显是“一条曲线” \\
$0.3$ & $0.50$ & $0.80$ & $0.063$ & 比例不变，视觉变尖 \\
$0.1$ & $0.50$ & $0.80$ & $0.063$ & 需要加密采样 \\
\bottomrule
\end{tabularx}
\end{center}

\begin{notebox}[这张表在说什么]
注意第二、三、四列\textbf{完全不随 $\ep$ 变}。这是 Cauchy 核的尺度不变性：所有比例都是固定的。

所以“集中”不能靠“看某个固定区间里的质量变多”来展示——固定区间 $[-a,a]$ 里，质量确实从 $0$ 涨到 $1$，但那是 $a$ 固定、$\ep$ 变小的结果。做动画时\textbf{一定要让区间固定、让 $\ep$ 变}，而不是反过来。
\end{notebox}

\section{讲解禁区（EP04）}

\begin{warnbox}[逐条对照]
\begin{enumerate}[leftmargin=1.6em,itemsep=0.28em]
\item 不说“$K_\ep$ 处处趋于零”。原点是例外。
\item 不说“一致收敛就能把极限搬进积分”。给出 $f_n=\frac1n\mathbf{1}_{[0,n]}$ 这个反例。
\item 不用“每一项都可积”代替“存在同一个可积控制函数”。
\item 不说“面积是 $1$，所以极限就是 $\delta$”。必须加“集中”。
\item 不说“$\delta$ 是某种特定形状的极限”。三种核已经否证了这一点。
\item 不说“$k_n$ 那个反例说明面积守恒是错的”。面积确实守恒，错的是“面积守恒足以决定分布极限”。
\item 不把 sinc 型核硬套进本节定理（它不非负、不绝对可积地集中）；留到 EP06。
\item 不说“$\Sv'$ 里也能测试总质量”。常数 $1\notin\Sv$。
\item 不把“分布收敛”讲成“某种更弱的逐点收敛”。它压根不是逐点概念。
\item 换元证明中，不要声称原来的 $K_\ep$ 存在公共控制函数。
\end{enumerate}
\end{warnbox}

\section{与前后集的接口}

\begin{enumerate}[leftmargin=2em,itemsep=0.45em]
\item \textbf{回头看}：EP01 借用的取样极限，这一集给了证明。EP03 里那个“合并正负半轴再用控制收敛”的估计，和这里是同一套路。
\item \textbf{往前看}：既然不同的核都给出同一个“作用”，那么能不能\textbf{直接把这个作用当作对象的定义}？——这就是 EP05 的全部动机。
\end{enumerate}

\begin{scriptbox}[最后 30 秒的旁白骨架]
“我们现在知道了：$\delta$ 不是某一种特别的尖峰。

指数核、高斯核、矩形核，形状差得很远，可是它们对每一个探针说的话，是一模一样的。

那么，为什么不干脆\textbf{用这句话来定义 $\delta$}？

下一集，我们就这么干。”
\end{scriptbox}


\section{本集回顾卡}

\begin{center}
\small
\renewcommand{\arraystretch}{1.4}
\begin{tabularx}{\textwidth}{@{}L{3.4cm}Y@{}}
\toprule
项目 & 内容 \\
\midrule
主结果 & $\displaystyle\int_{\R}K_\ep(x)\varphi(x)\dd x\to\varphi(0)$ \\
三种收敛 & 逐点、一致、分布；后者与点值无关 \\
许可证 & 控制收敛，且需要\textbf{同一个}可积控制函数 \\
两个反例 & $\frac1n\mathbf{1}_{[0,n]}$（一致趋零而积分不变）；$k_n=\frac12\mathbf{1}_{[n-1,n+1]}$（面积为 $1$ 却趋于 $0$）\\
安全话术 & “核没有变小，是质量集中了” \\
危险话术 & “面积是 $1$ 所以极限是 $\delta$” \\
下一集的钩子 & 能不能直接用“作用”来定义这个极限对象？ \\
\bottomrule
\end{tabularx}
\end{center}

\begin{quizbox}[EP04 自测]
\begin{enumerate}[leftmargin=1.6em,itemsep=0.35em]
\item 举出一个一致收敛但积分不收敛的函数列，并说明它为什么不违反控制收敛定理。
\item 证明命题 \ref{prop:noControl}（核族无公共可积控制）。
\item 用换元法证明取样极限 $\ip{K_\ep}{\varphi}\to\varphi(0)$，并指出控制函数是哪一行。
\item 为什么 $e^{-\ep|t|}\to1$ 于 $\Sv'$，但不能直接对全实轴用控制收敛？
\item 写出近似恒等核的三条充分条件，并说明每条各自防止什么。
\item 证明命题 \ref{prop:escaping}，并解释“总面积为 $1$ 却不趋于 $\delta$”为何不与任何定理矛盾。
\item 若把 $k_\ep$ 换成一个不非负的核，取样极限还成立吗？需要补上哪两条条件？请写出估计过程。
\end{enumerate}
\end{quizbox}
}{}
\IfFileExists{ep05.tex}{\chapter{EP05｜$\delta$ 到底是什么：分布理论的最低完整框架}

\begin{ideabox}[这一集的一句话]
\textbf{一个对象可以不由“它在每一点取什么值”来定义，而由“它怎样作用于每一个探针”来定义。$\delta$ 就是这样被定义的——这就是它“不是函数”却依然严格的全部原因。}
\end{ideabox}

\section{观众的问题链}

前四集一直在用 $\delta$、$\delta'$、$\PV$，却始终没有给它们下定义。这一集补上这个欠账，而且要用一种让人信服的方式补：\textbf{不是先抛定义再解释，而是先展示“作用”，再说明“这就是定义”。}

\begin{center}
\small
\renewcommand{\arraystretch}{1.32}
\begin{tabularx}{\textwidth}{@{}p{0.8cm}X p{4.4cm}@{}}
\toprule
层 & 问题 & 观众此刻的心理 \\
\midrule
1 & 三种完全不同的核，为什么都叫同一个 $\delta$？ & 它们唯一相同的地方是“对探针说的话”。\\
2 & 那能不能就用这个“说的话”来定义它？ & 听起来可行，但“说的话”要怎么写得严格？\\
3 & 需要一个探针的集合，以及“连续”的意思。 & 为什么要光滑？为什么要限制远处？\\
4 & 于是分布 = 连续线性泛函。 & 这个定义能装下普通函数吗？\\
5 & 能。$f\mapsto\int f\varphi$ 就是一个这样的泛函。 & 那 $\delta$ 呢？\\
6 & $\delta$ 不在其中——它不能由任何局部可积函数给出。 & 所以它是“新东西”，不是“旧东西的别名”。\\
\bottomrule
\end{tabularx}
\end{center}

\begin{center}
\small
\renewcommand{\arraystretch}{1.32}
\begin{tabularx}{\textwidth}{@{}p{3.2cm}Y@{}}
\toprule
观众带着什么进来 & 观众带着什么出去 \\
\midrule
“$\delta$ 是一个函数，只是很奇怪” & “$\delta$ 是一个泛函；普通函数只是它的一个子类。”\\
“分布是数学家造出来救场的” & “分布是把‘作用于探针’这件事本身当作对象；它是被逼出来的。”\\
“什么都能乘、什么都能求导” & “能求导、能乘光滑函数；但两个分布不能随便相乘。”\\
\bottomrule
\end{tabularx}
\end{center}

\section{从“作用”认识对象}

\subsection{普通函数也有“作用”}

对一个局部可积的函数 $f$，取一条光滑、紧支撑的曲线 $\varphi$，算

\begin{equation}
\varphi\longmapsto\int_{\R}f(x)\varphi(x)\dd x .
\end{equation}

这是一个规则：\textbf{输入一条探针，输出一个数。}

现在看 $\delta$：

\begin{equation}
\varphi\longmapsto\varphi(0).
\end{equation}

这也是一个规则：输入一条探针，输出一个数。

\begin{ideabox}[关键的一步观念转换]
平时我们说“$f$ 是函数”，是把它想成一张图：每个 $x$ 对应一个高度。

但换一个角度看，$f$ 也可以被完全刻画成“它对所有探针的响应”：知道 $\int f\varphi$ 对每个 $\varphi$ 是多少，就等于知道了 $f$（几乎处处）。

\textbf{那就干脆把这个“响应规则”本身当成对象。}这样，即使某个对象的“响应规则”不能由任何函数给出，它依然是一个合法的对象。

$\delta$ 就是这样的对象：它的响应规则极其简单——把探针在原点的值取出来。
\end{ideabox}

\subsection{一张对照表}

\begin{center}
\small
\renewcommand{\arraystretch}{1.55}
\begin{tabularx}{\textwidth}{@{}L{3.4cm}L{4.6cm}Y@{}}
\toprule
 & 普通函数 $f$ & 分布 $T$ \\
\midrule
怎样描述 & 每个 $x$ 对应一个值 $f(x)$ & 每条探针 $\varphi$ 对应一个数 $\ip{T}{\varphi}$ \\
能不能画图 & 能 & 一般不能（$\delta$ 就没有图） \\
怎样求导 & 差商 / 逐点 & 把导数转移到探针上（带负号） \\
能不能乘 & 一般能 & 只能乘光滑函数，不能乘另一个分布 \\
\bottomrule
\end{tabularx}
\end{center}

\section{测试函数：为什么既要光滑，又要限制远处}

\subsection{两个要求，各管一件事}

\begin{center}
\small
\renewcommand{\arraystretch}{1.5}
\begin{tabularx}{\textwidth}{@{}L{2.4cm}L{5.4cm}Y@{}}
\toprule
要求 & 数学说法 & 它负责什么 \\
\midrule
光滑 & 无限次可微 & 让“把导数搬过去”可以无限次进行；分部积分不产生边界项以外的麻烦 \\
紧支撑 & 在某个有限区间外恒为零 & 让分部积分的边界项消失；也让“局部观察”成为可能 \\
\bottomrule
\end{tabularx}
\end{center}

\begin{defn}[测试函数空间]\label{def:D}
\begin{equation}
\D(\R):=C_c^{\infty}(\R)=\{\varphi:\R\to\C\ ;\ \varphi\in C^{\infty},\ \supp\varphi\text{ 紧}\}.
\end{equation}
\end{defn}

\subsection{一个标准的非零例子}

\begin{equation}
\eta(x)=\begin{cases}
\exp\!\left(-\dfrac{1}{1-x^{2}}\right), & |x|<1,\\[8pt]
0, & |x|\ge1 .
\end{cases}\label{eq:bump}
\end{equation}

\begin{ideabox}[为什么这个函数是“光滑拼接”的典范]
当 $x\to\pm1$ 时，$1-x^{2}\to0^{+}$，所以 $-\frac{1}{1-x^{2}}\to-\infty$，于是 $e^{-\frac{1}{1-x^{2}}}\to0$。

关键在于\textbf{不只是函数值趋于零，所有阶导数都趋于零}。可以验证：$\eta$ 的每阶导数都含有因子 $e^{-\frac{1}{1-x^{2}}}$ 乘上一个关于 $\frac{1}{1-x^{2}}$ 的多项式，而指数衰减压倒多项式增长。

所以拼上 $|x|\ge1$ 处的零之后，$\eta$ 仍然是 $C^{\infty}$ 的。

\textbf{这一句要在视频里出现：}“它可以平白无故地在边界停下来，而且停得非常光滑，连斜率、曲率都一起停住。”普通的多项式做不到这一点——这正是 $\eta$ 值得单独介绍的原因。
\end{ideabox}

\subsection{高斯为什么不在 $\D$ 里}

$e^{-x^{2}}$ 光滑、衰减极快，但它在整个实轴上都不为零，\textbf{没有紧支撑}。所以

\begin{equation}
e^{-x^{2}}\in\Sv
\qquad\text{但}\qquad
e^{-x^{2}}\notin\D .
\end{equation}

\begin{notebox}[这个区别在 EP06 会变成主角]
$\D$ 适合“局部观察”，$\Sv$ 适合“整体变换”。

傅里叶变换会把紧支撑的东西变成不紧支撑的东西（矩形窗变 sinc），所以 $\D$ 在傅里叶变换下不封闭，而 $\Sv$ 封闭。这就是 EP06 必须换空间的原因。
\end{notebox}

\subsection{$\D$ 里的收敛}

$\varphi_j\to\varphi$ 于 $\D$，意思是：存在一个公共的紧集 $K$ 包含所有 $\supp\varphi_j$，并且对每个多重指标 $\alpha$，

\begin{equation}
\sup_{x\in K}\big|\varphi_j^{(\alpha)}(x)-\varphi^{(\alpha)}(x)\big|\longrightarrow0 .
\end{equation}

注意“公共紧集”这个要求：\textbf{不能让探针一边趋于零一边往无穷远跑}。

\section{分布的定义}

\begin{defn}[分布]\label{def:dist}
$\D'(\R)$ 中的元素 $T$ 是 $\D(\R)$ 上的连续线性泛函：给每个 $\varphi\in\D$ 分配一个复数 $\ip{T}{\varphi}$，并满足
\begin{enumerate}[leftmargin=1.6em,itemsep=0.25em]
\item \textbf{线性}：$\ip{T}{a\varphi+b\psi}=a\ip{T}{\varphi}+b\ip{T}{\psi}$；
\item \textbf{连续性}：若 $\varphi_j\to\varphi$ 于 $\D$，则 $\ip{T}{\varphi_j}\to\ip{T}{\varphi}$。
\end{enumerate}
\end{defn}

\subsection{线性不够，为什么还要连续}

\begin{ideabox}[一个“线性但不连续”的怪物]
在 $\D$ 上定义 $T$：把 $\varphi$ 展开成某种级数，然后取一个“不听话”的线性泛函。它可以线性，却对越来越小的探针给出越来越大的响应。

连续性就是用来禁止这种事的：\textbf{探针越小，响应必须越小。}
\end{ideabox}

连续性的一个等价刻画（局部估计）：

\begin{equation}
|\ip{T}{\varphi}|\le C_K\max_{0\le j\le m_K}\|\varphi^{(j)}\|_\infty,
\qquad \supp\varphi\subset K .\label{eq:local-est}
\end{equation}

也就是说：在每个紧集上，响应被\textbf{有限阶}导数的最大值控制。阶数 $m_K$ 可以依赖紧集，但必须是有限的。

\begin{notebox}[配对不是内积]
$\ip{T}{\varphi}$ 是\textbf{复线性}的：对 $\varphi$ 线性，对 $T$ 也线性。\textbf{没有复共轭。}

所以不要把它写成 $\langle T,\varphi\rangle=\int T\bar\varphi$ 那种 Hilbert 空间内积的形式。这一条在 EP06 定义傅里叶变换时会直接影响 $2\pi$ 与符号的位置。
\end{notebox}

\section{普通函数怎样住进这个新世界}

\subsection{嵌入}

\begin{prop}\label{prop:embed}
若 $f\in\Lonloc$，则
\begin{equation}
\ip{T_f}{\varphi}:=\int_{\R}f(x)\varphi(x)\dd x
\end{equation}
定义一个分布 $T_f$。而且 $T_f=T_g$ 当且仅当 $f=g$ 几乎处处。
\end{prop}

\begin{proof}[证明要点]
因为 $\varphi$ 紧支撑，积分只在一个有界区间上进行，所以由 $f$ 的局部可积性，积分存在。估计：
\begin{equation}
|\ip{T_f}{\varphi}|\le\|f\|_{L^{1}(K)}\|\varphi\|_\infty,
\qquad \supp\varphi\subset K .
\end{equation}
这正好是 \eqref{eq:local-est} 的形式（取 $m_K=0$），所以连续。
\end{proof}

\begin{ideabox}[这个嵌入解释了 EP03 的一个旧疑问]
$H(0)$ 取 $0$、$\frac12$ 还是 $1$？都不影响 $H$ 对应的分布，因为这些改值只发生在一个点上，而\textbf{一个点对积分没有贡献}。

所以“$H(0)$ 无关”不是含糊其辞，而是“分布只认几乎处处等价类”的直接后果。
\end{ideabox}

\subsection{$\delta$ 不在这个家族里}

\begin{prop}\label{prop:delta-not-L1loc}
不存在 $g\in\Lonloc$ 使 $T_g=\delta$。
\end{prop}

\begin{proof}[证明]
假设有。取任意支撑在 $\R\setminus\{0\}$ 内的 $\varphi\in\D$，则 $\ip{\delta}{\varphi}=0$，于是

\begin{equation}
\int_{\R}g(x)\varphi(x)\dd x=0
\qquad\text{对所有这样的 }\varphi .
\end{equation}

由局部可积函数的识别性（若 $\int g\varphi=0$ 对所有支撑在开集 $U$ 内的测试函数成立，则 $g=0$ 几乎处处于 $U$），得

\begin{equation}
g=0\ \text{几乎处处于}\ \R\setminus\{0\}.
\end{equation}

而 $\{0\}$ 是零测集，所以 $g=0$ 几乎处处于 $\R$，于是 $T_g=0\neq\delta$。矛盾。
\end{proof}

\begin{ideabox}[这一步的意义]
它证明 $\delta$ 是\textbf{真正的新对象}，不是“某个很奇怪但仍然是函数的函数”。

如果 $\delta$ 只是“在原点改一个值”，那么由上面的论证它必然几乎处处为零；但几乎处处为零的函数不可能满足 $\int g\varphi=\varphi(0)$。

\textbf{所以 $\delta$ 的出现不是技巧问题，是语言问题：旧的语言装不下它。}
\end{ideabox}

\subsection{但 $\delta$ 是测度}

$\delta$ 也可以理解为 Dirac 点质量：对集合 $A$，

\begin{equation}
\delta(A)=\begin{cases}1,&0\in A,\\0,&0\notin A .\end{cases}
\end{equation}

积分就写成 $\int\varphi\,\dd\delta=\varphi(0)$。

\begin{notebox}[“不是函数”$\neq$“没有严格定义”]
这是两句完全不同的话，写旁白时务必分开。
\begin{itemize}[leftmargin=1.4em,itemsep=0.2em]
\item $\delta$ 不是（局部可积）函数——这是对的，命题 \ref{prop:delta-not-L1loc} 证明了。
\item $\delta$ 没有严格定义——这是错的。它是测度，也是分布，有两种严格的解释。
\end{itemize}
\end{notebox}

\section{运算：能做什么}

\subsection{求导}

\begin{equation}
\ip{T'}{\varphi}=-\ip{T}{\varphi'} .
\end{equation}

定义良好的理由：$\varphi\in\D\Rightarrow\varphi'\in\D$；且 $\varphi_j\to\varphi$ 于 $\D$ 时 $\varphi_j'\to\varphi'$ 于 $\D$。

\begin{prop}\label{prop:deriv-cont}
若 $T_j\to T$ 于 $\D'$，则 $T_j'\to T'$ 于 $\D'$。
\end{prop}

\begin{proof}
$\ip{T_j'}{\varphi}=-\ip{T_j}{\varphi'}\to-\ip{T}{\varphi'}=\ip{T'}{\varphi}$。
\end{proof}

\begin{ideabox}[这条性质解释了 EP02 的做法]
EP02 里我们“先证明 $K_\ep\to\delta$，再直接对核求导，得到 $\delta'$ 的近似”。

这个操作之所以合法，正是命题 \ref{prop:deriv-cont}：\textbf{分布收敛对求导是封闭的。}不需要额外验证 $K_\ep'$ 的点值收敛——事实上它根本没有点值极限。
\end{ideabox}

\subsection{乘光滑函数}

对 $a\in C^{\infty}(\R)$，定义

\begin{equation}
\ip{aT}{\varphi}:=\ip{T}{a\varphi} .
\end{equation}

定义良好的理由：$a\varphi\in\D$。

由此立刻得到几条“代数式”：

\begin{align}
a\,\delta&=a(0)\,\delta, \label{eq:adelta}\\
x\,\delta&=0, \label{eq:xdelta}\\
x\,\delta'&=-\delta, \label{eq:xdelta'}\\
x\,\delta^{(n)}&=-n\,\delta^{(n-1)} . \label{eq:xdeltaN}
\end{align}

\begin{proof}[逐条验证]
式 \eqref{eq:adelta}：$\ip{a\delta}{\varphi}=\ip{\delta}{a\varphi}=a(0)\varphi(0)=a(0)\ip{\delta}{\varphi}$。

式 \eqref{eq:xdelta'}：$\ip{x\delta'}{\varphi}=\ip{\delta'}{x\varphi}=-(x\varphi)'(0)=-\big(\varphi(0)+0\big)=-\varphi(0)$。

式 \eqref{eq:xdeltaN}：因为 $(x\varphi)^{(n)}=x\varphi^{(n)}+n\varphi^{(n-1)}$，在 $0$ 处取值只剩 $n\varphi^{(n-1)}(0)$，于是
\begin{equation}
\ip{x\delta^{(n)}}{\varphi}=(-1)^{n}n\varphi^{(n-1)}(0)=-n\cdot(-1)^{n-1}\varphi^{(n-1)}(0)=-n\ip{\delta^{(n-1)}}{\varphi}.
\end{equation}
\end{proof}

\begin{ideabox}[$x\delta=0$ 与 $x\delta'=-\delta$ 是两道最好的练习题]
它们看起来像“用零乘一个无穷大”，其实是完全严格的。

先用 \eqref{eq:xdelta} 检查 EP03 的歧义：$\omega\delta=0$，所以加一项 $c\delta$ 不影响 $\omega\hat H=1$。这一步现在有了严格依据。

再用 \eqref{eq:xdeltaN} 检查 EP03 的“为什么不能再加 $\delta'$”：$x\delta'=-\delta\neq0$，所以 $\delta'$ 不满足“乘 $\omega$ 为零”。
\end{ideabox}

\subsection{Leibniz 规则}

\begin{equation}
(aT)'=a'T+aT' .
\end{equation}

证明：$\ip{(aT)'}{\varphi}=-\ip{aT}{\varphi'}=-\ip{T}{a\varphi'}=-\ip{T}{(a\varphi)'-a'\varphi}=\ip{T}{a'\varphi}+\ip{aT'}{\varphi}$。

\subsection{支撑}

\begin{defn}[在开集上为零]
若 $T$ 对一切 $\supp\varphi\subset U$（$U$ 开）的 $\varphi$ 都给出 $\ip{T}{\varphi}=0$，就说 $T$ 在 $U$ 上为零。$T$ 的\textbf{支撑} $\supp T$ 是不能这样“消掉”的最小闭集。
\end{defn}

\begin{itemize}[leftmargin=1.6em,itemsep=0.3em]
\item $\supp\delta=\{0\}$，$\supp\delta^{(n)}=\{0\}$。
\item $\supp H=[0,\infty)$：在 $(-\infty,0)$ 上 $H=0$，在 $0$ 的任何右邻域内都不能被消掉。
\item 若 $\varphi\equiv0$ 于 $0$ 的邻域，则 $\ip{\delta^{(n)}}{\varphi}=0$，所以支撑在 $\{0\}$。
\end{itemize}

\subsection{平移、缩放、复合}

平移：$\ip{\delta_a}{\varphi}=\varphi(a)$，常写 $\delta(x-a)$。

缩放（$a\neq0$）：

\begin{equation}
\delta(ax)=\frac{1}{|a|}\delta(x).\label{eq:scaling-delta}
\end{equation}

\begin{ideabox}[为什么是 $1/|a|$ 而不是 $1/a$]
因为这是换元的 Jacobian 因子，不是代数倒数。直观上：把自变量压缩 $|a|$ 倍，点质量要相应地“变浓” $|a|$ 倍，才能保持总质量 $1$。

\textbf{这是最容易写成 $1/a$ 的地方}，而 $a<0$ 时符号就错了。
\end{ideabox}

复合：若光滑函数 $g$ 的零点 $x_j$ 都是\textbf{简单零点}（$g'(x_j)\neq0$）且局部有限，则

\begin{equation}
\delta\big(g(x)\big)=\sum_{j}\frac{\delta(x-x_{j})}{|g'(x_{j})|} .\label{eq:delta-compose}
\end{equation}

\begin{warnbox}[简单零点这个条件不能省]
$g(x)=x^{2}$ 在 $0$ 处是二重零点，$g'(0)=0$，\eqref{eq:delta-compose} 的分母为零，公式失效。此时 $\delta(x^{2})$ 根本不是这个形式的对象。

讲这一条时，务必把“简单零点”说出口。
\end{warnbox}

\subsection{跳跃公式（回顾 EP03）}

若 $f$ 分段 $C^{1}$、在 $a_j$ 处有局部有限个跳跃 $J_j$，则

\begin{equation}
Df=f'_{\text{经典}}+\sum_{j}J_j\,\delta_{a_j} .
\end{equation}

例：$D|t|=\sgn t$，$D^{2}|t|=2\delta$。

\section{哪些运算不能随便做}

这是本集最“值钱”的一节，因为它决定了观众以后会不会乱用公式。

\begin{warnbox}[四条禁令]
\begin{enumerate}[leftmargin=1.6em,itemsep=0.4em]
\item \textbf{两个任意分布不能相乘。}$\delta^{2}$ 在经典分布理论中没有通用的规范定义。不要写 $|\delta|^{2}$，更不要算它的“能量”。
\item \textbf{$H\delta$ 没有内禀定义。}因为 $H$ 不是光滑乘子（它不连续），不同的正则化约定可能给出不同结果。用 $H(0)$ 的某个取值来“定义”它是靠不住的。
\item \textbf{卷积有存在条件。}紧支撑分布与任意分布可以卷积；两个任意分布未必可以。
\item \textbf{“分布”不是“概率分布”的同义词。}$\delta$ 是两者交集中的例子，但 $\delta'$ 不是正测度，$\PV\frac1\omega$ 也不是。
\end{enumerate}
\end{warnbox}

\begin{ideabox}[为什么要专门列“不能做什么”]
因为视频里最容易被观众模仿的，是公式的\textbf{形式}，而不是它的\textbf{条件}。

如果 EP02 出现了 $\F[t^{n}]=2\pi i^{n}\delta^{(n)}$，就会有人问：“那 $\F[t^{n}\cdot t^{m}]$ 是不是两个 $\delta$ 相乘？”

提前把“不能乘”说清楚，比事后解释便宜得多。
\end{ideabox}

\section{两个能检验理解的例子}

\subsection{例一：$x\delta'=-\delta$}

\begin{align}
\ip{x\delta'}{\varphi}
&=\ip{\delta'}{x\varphi}=-\big(x\varphi\big)'(0)
=-\big(\varphi(0)+0\cdot\varphi'(0)\big)=-\varphi(0)=-\ip{\delta}{\varphi}.
\end{align}

\subsection{例二：$(1+x)\delta'$}

\begin{align}
\ip{(1+x)\delta'}{\varphi}
&=\ip{\delta'}{(1+x)\varphi}
=-\big((1+x)\varphi\big)'(0)
=-\Big(\varphi(0)+(1+0)\varphi'(0)\Big)\notag\\
&=-\varphi(0)-\varphi'(0)
=\ip{-\delta+\delta'}{\varphi}.
\end{align}

所以 $(1+x)\delta'=\delta'-\delta$。这个结果可以用 Leibniz 检查：$(1+x)'\delta'+(1+x)\delta''=\delta'+\delta''+x\delta''=\delta'-\delta$（最后一步用了 $x\delta''=-\delta'$）。

\section{动画要点}

\begin{scriptbox}[动画要点 05-A / 05-C / 05-E]
\textbf{05-A}：一条测试函数曲线“走进去”，一个数字“出来”。同一个对象，换不同的探针，给不同的数。

\textbf{05-C}：把探针平移，或改变它在原点的值，看输出跟着 $\varphi(0)$ 变。这是 $\delta$ 的“指纹”。

\textbf{05-E}：分部积分的可视化——导数像接力棒一样从对象传递到探针，同时出现一个负号。

\textbf{05-F}：把第三集的跳变重新解释一次：$H$ 的跳跃强度为 $1$，所以在 $0$ 处留下权重为 $1$ 的 $\delta$。观众应当有“原来如此”的感觉。
\end{scriptbox}

\section{讲解禁区（EP05）}

\begin{warnbox}[逐条对照]
\begin{enumerate}[leftmargin=1.6em,itemsep=0.28em]
\item 不说“$\delta$ 是原点处无穷高、其余为零的函数”。
\item 不说“分布是一种更奇怪的函数”。它是泛函。
\item 不把 $\ip{T}{\varphi}$ 写成含复共轭的内积。
\item 不说“任意两个分布都能相乘/卷积”。
\item 不写 $\delta^{2}$，不写 $\|\delta\|^{2}$，不把 $\delta$ 当概率密度以外的“普通密度”。
\item 不说“局部可积函数都可以定义分布，所以 $\delta$ 也可以”。$\delta$ 恰恰是反例。
\item 不说“$H(0)$ 的取值影响了 $H$ 这个分布”。
\item 不把 $\delta(ax)$ 写成 $\frac1a\delta(x)$（漏绝对值）。
\item 不省略 $\delta(g(x))$ 公式中的“简单零点”条件。
\item 不在 $\D'$ 上随意做“乘一个不光滑的函数”的操作。
\item 不把“支撑在 $\{0\}$”讲成“$\delta$ 只影响原点”这种含糊说法；要说“对任何在原点邻域为零的探针，作用为零”。
\end{enumerate}
\end{warnbox}

\section{与前后集的接口}

\begin{enumerate}[leftmargin=2em,itemsep=0.45em]
\item \textbf{回头看}：EP01--EP04 里所有的“趋于”，现在有了名字：\textbf{分布收敛}。EP02 的“求导”现在有了定义。EP03 的 $\PV$ 现在可以写成一个分布。
\item \textbf{往前看}：还差最后一块——傅里叶变换。$\F$ 是在积分里定义的，而分布没有积分。怎么办？
\end{enumerate}

\begin{scriptbox}[最后 30 秒的旁白骨架]
“现在我们终于知道 $\delta$ 是什么了：它是一条规则，一条‘取出你函数在原点值’的规则。

也知道它能求导，能乘光滑函数，不能和别的分布相乘。

但还剩一个问题。

从第一集开始，我们一直在说 $\delta$ 的‘频谱’、说 $\F[1]=2\pi\delta$。可傅里叶变换是一个积分，而 $\delta$ 根本没有积分。

下一集，我们把这个变换也搬到这个世界里来。”
\end{scriptbox}


\section{本集回顾卡}

\begin{center}
\small
\renewcommand{\arraystretch}{1.4}
\begin{tabularx}{\textwidth}{@{}L{3.4cm}Y@{}}
\toprule
项目 & 内容 \\
\midrule
定义 & 分布 $=$ $\D$ 上的连续线性泛函；配对不含共轭 \\
嵌入 & 每个 $f\in\Lonloc$ 定义 $\ip{T_f}{\varphi}=\int f\varphi$ \\
$\delta$ 的身份 & $\ip{\delta}{\varphi}=\varphi(0)$；不由任何局部可积函数给出 \\
关键恒等式 & $x\delta=0$，$x\delta'=-\delta$，$x\delta^{(n)}=-n\delta^{(n-1)}$ \\
禁区 & 两个分布相乘、$\delta^{2}$、$H\delta$ 都没有通用定义 \\
安全话术 & “由它做什么来定义，不由它多高来定义” \\
下一集的钩子 & 傅里叶变换是一个积分，可它没有积分 \\
\bottomrule
\end{tabularx}
\end{center}

\begin{quizbox}[EP05 自测]
\begin{enumerate}[leftmargin=1.6em,itemsep=0.35em]
\item 独立验证 $x\delta'=-\delta$ 与 $x\delta''=-\delta'$。
\item 计算 $(2+x)\delta''$ 与 $x^{2}\delta''$。
\item 证明 $a\delta=a(0)\delta$，并解释为什么 $H\delta$ 不能这样处理。
\item 说明为什么 $\eta$ 在 \eqref{eq:bump} 中定义后属于 $\D$，而 $e^{-x^{2}}$ 不属于。
\item 写出 $\delta(3x-6)$ 的化简结果。
\item 用 \eqref{eq:local-est} 说明：为什么“连续性”排除了“探针越小、响应越大”的泛函。
\item 证明命题 \ref{prop:delta-not-L1loc}，并指出证明中哪一步用了“单点集是零测集”。
\end{enumerate}
\end{quizbox}
}{}
\IfFileExists{ep06.tex}{\chapter{EP06｜发散的积分怎样被严格定义}

\begin{ideabox}[这一集的一句话]
\textbf{只要把探针从“紧支撑光滑函数”换成 Schwartz 函数，傅里叶变换就能整体搬到分布上，并且保持极限。前三集的所有计算，在这一步同时获得合法性。}
\end{ideabox}

\section{观众的问题链}

这一集要还清全部欠账。它的难点不在计算，而在\textbf{为什么必须换一个空间}——这一点如果讲不透，整集就会变成“数学家又定义了一个新东西”。

\begin{center}
\small
\renewcommand{\arraystretch}{1.32}
\begin{tabularx}{\textwidth}{@{}p{0.8cm}X p{4.4cm}@{}}
\toprule
层 & 问题 & 观众此刻的心理 \\
\midrule
1 & 有了分布，为什么不直接定义 $\ip{\F T}{\varphi}=\ip{T}{\F\varphi}$？ & 看起来很自然啊。\\
2 & 因为 $\F\varphi$ 可能不再紧支撑。 & 那又怎样？\\
3 & 那 $\ip{T}{\F\varphi}$ 就没定义了——$T$ 只认紧支撑的探针。 & 哦，是个技术问题。\\
4 & 不是技术问题：矩形窗的变换是 sinc，它铺满整条实轴。 & 所以紧支撑这条路走不通。\\
5 & 换一个在变换下封闭的空间：Schwartz 空间。 & 这个空间凭什么刚好够用？\\
6 & 因为它对“求导”和“乘多项式”都封闭，而这两件事正好对应变换的两边。 & 那反演呢？\\
7 & 反演用高斯正则化来证，绕开循环论证。 & 于是三条公式全部落地。\\
\bottomrule
\end{tabularx}
\end{center}

\begin{center}
\small
\renewcommand{\arraystretch}{1.32}
\begin{tabularx}{\textwidth}{@{}p{3.2cm}Y@{}}
\toprule
观众带着什么进来 & 观众带着什么出去 \\
\midrule
“前三集的结果是靠技巧凑出来的” & “它们都是同一个定义的三条推论。”\\
“为什么不能对所有分布做变换？” & “因为 $\F$ 不保持 $\D$；它保持 $\Sv$。”\\
“$\delta$ 的变换是 $1$ 是怎么来的” & “由对偶定义一步得出，而且和 $2\pi\delta$ 互为反演。”\\
\bottomrule
\end{tabularx}
\end{center}

\section{为什么不能直接对 $\D'$ 照搬}

\subsection{想做的那件事}

最自然的想法是：对 $T\in\D'$，定义

\begin{equation}
\ip{\F T}{\varphi}:=\ip{T}{\F\varphi} .
\end{equation}

要让这个定义有意义，右边必须能算——也就是必须有 $\F\varphi\in\D$。

\subsection{但它不成立}

\begin{prop}\label{prop:F-not-D}
若 $\varphi\in\D$ 且 $\varphi\neq0$，则 $\F\varphi\notin\D$。
\end{prop}

\begin{proof}[证明：解析性]
对任意复数 $\zeta$，积分
\begin{equation}
\widehat{\varphi}(\zeta)=\int_{\R}\varphi(t)e^{-i\zeta t}\dd t
\end{equation}
绝对收敛（$\varphi$ 紧支撑），并且对 $\zeta$ 是\textbf{整函数}（在整个复平面上解析）。

整函数有一个硬性质：\textbf{若它在一个开集上为零，则处处为零。}

如果 $\F\varphi$ 有紧支撑，那么它在某个非空开集上为零，于是 $\widehat\varphi\equiv0$，再由反演得 $\varphi\equiv0$。矛盾。
\end{proof}

\begin{ideabox}[一个不用复分析的直观版本]
矩形窗 $w_T=\mathbf{1}_{[-T,T]}$ 的变换是
\begin{equation}
\widehat{w_T}(\omega)=\frac{2\sin(T\omega)}{\omega},
\end{equation}
它\textbf{在整条实轴上一直振荡下去}，根本没有紧支撑。

光滑化一点不会改变这件事：把矩形窗的边缘磨圆，得到的东西仍然“两头都拖得很长”。

\textbf{一句话：时间上被限制住，频率上就必然铺开。}这正是 EP01 的“时域铺开、频域集中”的镜像说法。
\end{ideabox}

\subsection{所以需要一个“变换下封闭”的探针空间}

要求很明确：

\begin{center}
\textbf{找到一个函数空间 $\mathcal{X}$，使得 $\F:\mathcal{X}\to\mathcal{X}$ 是连续可逆的，并且 $\D\subset\mathcal{X}$。}
\end{center}

这样就能对 $\mathcal{X}'$ 定义变换，而 $\mathcal{X}'$ 里包含了我们要的所有对象（常数、多项式、阶跃、$\delta$）。

这个 $\mathcal{X}$ 就是 Schwartz 空间。

\section{Schwartz 空间}

\begin{defn}[Schwartz 空间]\label{def:S}
$\Sv(\R)$ 由所有满足
\begin{equation}
p_{m,n}(\varphi):=\sup_{x\in\R}(1+|x|)^{m}\big|\varphi^{(n)}(x)\big|<\infty,
\qquad m,n=0,1,2,\dots
\end{equation}
的光滑函数 $\varphi$ 组成。
\end{defn}

换句话说：\textbf{每一阶导数，都比任何负幂衰减得更快。}

\subsection{谁在里面，谁不在}

\begin{center}
\small
\renewcommand{\arraystretch}{1.4}
\begin{tabularx}{\textwidth}{@{}L{4.2cm}L{2.0cm}Y@{}}
\toprule
函数 & 是否属于 $\Sv$ & 原因 \\
\midrule
$e^{-x^{2}}$ & 是 & 导数带多项式因子，被高斯压住 \\
$C_c^{\infty}$ 中的函数 & 是 & 紧支撑比快速衰减更强 \\
$1/(1+x^{2})$ & \textbf{否} & 可积，但只按固定幂次 $x^{-2}$ 衰减 \\
$1$、$x^{n}$ & \textbf{否} & 根本不衰减 \\
$e^{x^{2}}$ & \textbf{否} & 增长 \\
\bottomrule
\end{tabularx}
\end{center}

\begin{warnbox}[$1/(1+x^{2})$ 是最有教育意义的一个反例]
它\textbf{可积}，甚至属于 $\Lone$，但它不在 $\Sv$。

为什么？因为 $\Sv$ 要求的不是“可积”，而是“\textbf{每一阶导数都比任意负幂快}”。$1/(1+x^{2})\sim x^{-2}$ 在无穷远处衰减得太慢，而且它的高阶导数衰减得更慢。

\textbf{“可积”与“快速衰减”是两件不同的事。}这一点在 EP01（$\Lone$ 不够用）和这一集（$\Lone$ 里的东西也未必在 $\Sv$）各出现一次，方向相反。
\end{warnbox}

\subsection{$\Sv$ 里的收敛}

$\varphi_j\to\varphi$ 于 $\Sv$，意思是\textbf{所有半范数}都趋于零：

\begin{equation}
p_{m,n}(\varphi_j-\varphi)\longrightarrow0
\qquad\text{对所有的 }m,n .
\end{equation}

这比“函数值一致趋近”强得多：\textbf{每一阶导数、在每一个加权尺度上，都必须趋近。}

\section{为什么傅里叶变换保持 $\Sv$}

\subsection{两条关键恒等式}

第一条：对 $\omega$ 求导，等于在时域乘 $-it$。

\begin{equation}
\partial_\omega^{\,n}\ph(\omega)=\F\big[(-it)^{n}\varphi\big](\omega).\label{eq:d1}
\end{equation}

第二条：在频域乘 $\omega$，等于在时域求导。

\begin{equation}
\omega^{m}\,\widehat{g}(\omega)=i^{-m}\,\F\big[g^{(m)}\big](\omega).\label{eq:d2}
\end{equation}

（第二条来自 $\F[g^{(m)}]=(i\omega)^{m}\hat g$，两边整理即可。）

\subsection{估计}

把两条合起来：对 $\varphi\in\Sv$，

\begin{equation}
\omega^{m}\,\partial_\omega^{\,n}\ph(\omega)
=i^{-m}\int_{\R}\partial_t^{m}\big(t^{n}\varphi(t)\big)e^{-i\omega t}\dd t .
\end{equation}

取绝对值：

\begin{equation}
\big|\omega^{m}\partial_\omega^{\,n}\ph(\omega)\big|
\le\int_{\R}\big|\partial_t^{m}\big(t^{n}\varphi(t)\big)\big|\dd t
=\big\|\partial_t^{m}(t^{n}\varphi)\big\|_1 .\label{eq:Schwartz-bound}
\end{equation}

右边与 $\omega$ 无关。

\begin{prop}\label{prop:F-S}
$\F$ 把 $\Sv$ 连续地映到 $\Sv$。
\end{prop}

\begin{proof}[证明]
$\varphi\in\Sv$ 时 $t^{n}\varphi\in\Sv$，所以它的每阶导数都在 $\Lone$ 里（快速衰减），于是 \eqref{eq:Schwartz-bound} 右边有限。

对 $\omega$ 取上确界：
\begin{equation}
p_{m,n}(\ph)=\sup_\omega(1+|\omega|)^{m}\big|\partial^{n}\ph(\omega)\big|
\le C\sum_{m'\le m+1}\big\|\partial^{m'}(t^{n}\varphi)\big\|_1 .
\end{equation}
右边被有限个 $\Sv$ 半范数控制，所以 $\ph\in\Sv$，且映射连续。
\end{proof}

\begin{ideabox}[为什么这个证明“正好”能成]
注意两边的对称性：
\begin{itemize}[leftmargin=1.4em,itemsep=0.2em]
\item 频域求导 $\longleftrightarrow$ 时域乘 $t$；
\item 频域乘 $\omega$ $\longleftrightarrow$ 时域求导。
\end{itemize}
$\Sv$ 的定义恰好对这两件事\textbf{都封闭}：乘 $t$ 不破坏快速衰减，求导也不破坏。所以变换后的函数仍然满足所有半范数有界。

\textbf{$\Sv$ 不是随便挑的，它是被这两条恒等式“逼”出来的。}
\end{ideabox}

\section{反演：用高斯正则化绕开循环论证}

\subsection{不能走的那条捷径}

有一个极诱人的写法：

\begin{equation}
\int_{\R}e^{i\omega(t-s)}\dd\omega=2\pi\,\delta(t-s)
\quad\Longrightarrow\quad
\varphi(t)=\frac{1}{2\pi}\int_{\R}\ph(\omega)e^{i\omega t}\dd\omega .
\end{equation}

\begin{warnbox}[这是循环论证]
那个“等式”本身就是在分布意义下、并且要先用对偶定义才写得出来的。用它证反演，等于用结论证前提。

EP06 的证明必须走一条\textbf{不依赖 $\delta$} 的路。这条路就是高斯正则化。
\end{warnbox}

\subsection{高斯积分}

先用一个可以直接算的经典结果（证明见附录 A.1）：

\begin{equation}
\int_{\R}e^{-\eta\omega^{2}/2}e^{i\omega x}\dd\omega=\sqrt{\frac{2\pi}{\eta}}\;e^{-x^{2}/(2\eta)},
\qquad \eta>0 .\label{eq:gauss-int}
\end{equation}

\subsection{加一个高斯因子}

定义

\begin{equation}
I_\eta(t):=\frac{1}{2\pi}\int_{\R}\ph(\omega)e^{i\omega t}e^{-\eta\omega^{2}/2}\dd\omega .
\end{equation}

高斯因子 $e^{-\eta\omega^{2}/2}$ 的作用是：\textbf{把一切都变成绝对收敛的。}

\subsection{交换积分：这一步现在合法了}

把 $\ph(\omega)=\int\varphi(s)e^{-i\omega s}\dd s$ 代进去：

\begin{equation}
I_\eta(t)=\frac{1}{2\pi}\iint \varphi(s)\,e^{-i\omega s}e^{i\omega t}e^{-\eta\omega^{2}/2}\dd s\,\dd\omega .
\end{equation}

被积函数的绝对值是 $|\varphi(s)|e^{-\eta\omega^{2}/2}$，对 $(\omega,s)$ 双重可积（$\varphi\in\Sv\subset\Lone$，高斯在 $\omega$ 上可积）。所以 Fubini 定理适用，可以交换次序：

\begin{equation}
I_\eta(t)=\frac{1}{2\pi}\int_{\R}\varphi(s)\left[\int_{\R}e^{-\eta\omega^{2}/2}e^{i\omega(t-s)}\dd\omega\right]\dd s .
\end{equation}

用 \eqref{eq:gauss-int}：

\begin{equation}
I_\eta(t)=\int_{\R}\varphi(s)\,\frac{1}{\sqrt{2\pi\eta}}\,e^{-(t-s)^{2}/(2\eta)}\dd s
=\big(\rho_\eta*\varphi\big)(t),\label{eq:conv-gauss}
\end{equation}

其中 $\rho_\eta(u)=\frac{1}{\sqrt{2\pi\eta}}e^{-u^{2}/(2\eta)}$ 正是 EP04 的高斯近似恒等核。

\subsection{两边各取极限}

\textbf{左边：}因为 $\ph\in\Sv\subset\Lone$，且 $|e^{-\eta\omega^{2}/2}|\le1$，由控制收敛（控制函数 $|\ph|$ 可积且与 $\eta$ 无关），

\begin{equation}
I_\eta(t)\longrightarrow\frac{1}{2\pi}\int_{\R}\ph(\omega)e^{i\omega t}\dd\omega .
\end{equation}

\textbf{右边：}由 EP04 的近似恒等核定理，$\rho_\eta*\varphi\to\varphi$（逐点，且一致）。

于是

\begin{equation}
\boxed{\ \varphi(t)=\frac{1}{2\pi}\int_{\R}\ph(\omega)e^{i\omega t}\dd\omega\ }\label{eq:inversion}
\end{equation}

\begin{ideabox}[这个证明漂亮在哪里]
它\textbf{只用了三件已经证过的事}：
\begin{enumerate}[leftmargin=1.4em,itemsep=0.2em]
\item 高斯积分 \eqref{eq:gauss-int}（可以硬算）；
\item Fubini（由高斯因子保证绝对收敛）；
\item 近似恒等核（EP04 已经完整证明）。
\end{enumerate}

它\textbf{没有}用到 $\delta$、没有用到“$\int e^{i\omega t}\dd\omega$”，所以没有任何循环。

这是整个系列里技术含量最高、也最值得在视频里“给一个完整证明”的地方。
\end{ideabox}

\subsection{一次变换两次：反射}

由 \eqref{eq:inversion}，把 $t$ 换成 $-t$：

\begin{equation}
\F^{2}\varphi(t)=\int_{\R}\ph(\omega)e^{-i\omega t}\dd\omega
=2\pi\,\varphi(-t).\label{eq:F2}
\end{equation}

\begin{notebox}[记住这个“两次变换 = 反射 + $2\pi$”]
它是后面所有“反演”计算的发动机。

例如：$\F[\delta]=1$ 立刻给出 $\F[1]=2\pi\delta$。这就是 EP01 的计算不必再重复的原因。
\end{notebox}

\section{温和分布与对偶定义}

\begin{defn}[温和分布]\label{def:Sprime}
$\Sv'(\R)$ 是 $\Sv(\R)$ 上的连续线性泛函全体。连续性可写成：存在常数 $C$ 和有限的 $M,N$，使
\begin{equation}
|\ip{T}{\varphi}|\le C\sum_{m\le M,\ n\le N}p_{m,n}(\varphi)
\qquad\text{对一切 }\varphi\in\Sv .\label{eq:Sprime-est}
\end{equation}
\end{defn}

\subsection{谁在 $\Sv'$ 里}

\begin{center}
\small
\renewcommand{\arraystretch}{1.4}
\begin{tabularx}{\textwidth}{@{}L{3.6cm}L{4.4cm}Y@{}}
\toprule
对象 & 为什么温和 & 用哪个半范数控制 \\
\midrule
有界函数 & $|\ip{f}{\varphi}|\le\|f\|_\infty\|\varphi\|_1$ & $p_{2,0}$（因为 $(1+|x|)^{-2}$ 可积）\\
多项式增长函数 & $|f|\le C(1+|x|)^{N}$ & $p_{N+2,0}$ \\
$\delta^{(n)}$ & $\ip{\delta^{(n)}}{\varphi}=(-1)^{n}\varphi^{(n)}(0)$ & $p_{0,n}$ \\
$\PV\frac{1}{\omega}$ & 近处用一阶导估计，远处用衰减 & $p_{0,1}+p_{2,0}$ \\
\bottomrule
\end{tabularx}
\end{center}

\subsection{一个不温和的例子}

\begin{prop}\label{prop:ex2}
$e^{x^{2}}$ 是局部可积的，所以定义了一个分布；但它\textbf{不是}温和分布。
\end{prop}

\begin{proof}[证明思路]
取固定的非负 $0\neq\psi\in C_c^{\infty}$，令 $\psi_R(x)=\psi(x-R)$。

一方面，$\Sv$ 半范数 $p_{m,n}(\psi_R)$ 随 $R$ 至多按\textbf{多项式}增长（平移只让权重 $(1+|x|)^{m}$ 变大约 $R^{m}$ 倍）。

另一方面，
\begin{equation}
\ip{e^{x^{2}}}{\psi_R}=\int e^{x^{2}}\psi(x-R)\dd x\ge e^{(R-r)^{2}}\int\psi
\end{equation}
（其中 $r$ 是 $\psi$ 的支撑半径），随 $R$ 按\textbf{指数}增长。

指数增长压不过多项式增长，所以 \eqref{eq:Sprime-est} 不可能成立。
\end{proof}

\begin{ideabox}[这个反例的意义]
它回答了观众一个必然会问的问题：“是不是所有东西都能做傅里叶变换？”

\textbf{不是。}能做变换的不是“所有函数”，而是“所有温和分布”。$e^{x^{2}}$ 太大了，连 $\Sv$ 上的连续泛函都做不成。

这句话要说得准确：不是“$e^{x^{2}}$ 没有傅里叶变换”这种笼统说法，而是“在这个框架里它不在工作空间内”。
\end{ideabox}

\subsection{对偶定义}

先用普通函数把公式推出来。对 $f,\varphi\in\Sv$，由 Fubini（绝对收敛），

\begin{equation}
\ip{\F f}{\varphi}
=\iint f(t)e^{-i\omega t}\varphi(\omega)\dd t\dd\omega
=\int_{\R}f(t)\Big[\int_{\R}\varphi(\omega)e^{-i\omega t}\dd\omega\Big]\dd t
=\int_{\R}f(t)\ph(t)\dd t
=\ip{f}{\F\varphi} .
\end{equation}

这提示我们：

\begin{defn}[分布的傅里叶变换]\label{def:FTdist}
对 $T\in\Sv'$，定义 $\F T\in\Sv'$ 为
\begin{equation}
\boxed{\ \ip{\F T}{\varphi}:=\ip{T}{\F\varphi}\ ,\qquad \varphi\in\Sv .\ }\label{eq:dual}
\end{equation}
\end{defn}

\begin{ideabox}[四条必须同时说清的细节]
\begin{enumerate}[leftmargin=1.4em,itemsep=0.3em]
\item \textbf{右边有意义}：因为 $\varphi\in\Sv\Rightarrow\F\varphi\in\Sv$（命题 \ref{prop:F-S}），而 $T$ 能吃 $\Sv$。
\item \textbf{结果仍是温和分布}：$\varphi\mapsto\ip{T}{\F\varphi}$ 是 $\Sv$ 上的连续线性泛函（两次连续映射的复合）。
\item \textbf{和普通变换一致}：上面那段 Fubini 就是这个一致性。
\item \textbf{配对里没有复共轭，右边也没有额外的 $1/(2\pi)$}。\eqref{eq:dual} 是一个“干净”的公式；如果你见到的版本多了一个 $2\pi$，那是约定不同。
\end{enumerate}
\end{ideabox}

\subsection{逆变换}

由 \eqref{eq:F2} 与对偶定义，可以推出

\begin{equation}
\F^{-1}T=\frac{1}{2\pi}\,\F\big(T\circ(-1)\big),
\qquad\text{即}\qquad
\ip{\F^{-1}T}{\varphi}=\frac{1}{2\pi}\ip{T}{\F\varphi\circ(-1)} .
\end{equation}

在 $\Sv'$ 上，$\F$ 是一个连续线性同构。

\section{为什么变换保持极限}

\begin{prop}\label{prop:F-cont}
若 $T_j\to T$ 于 $\Sv'$，则 $\F T_j\to\F T$ 于 $\Sv'$。
\end{prop}

\begin{proof}
对每个固定的 $\varphi\in\Sv$，
\begin{equation}
\ip{\F T_j}{\varphi}=\ip{T_j}{\F\varphi}\longrightarrow\ip{T}{\F\varphi}=\ip{\F T}{\varphi} .
\end{equation}
\end{proof}

\begin{ideabox}[这三行补上了前三集缺的那一环]
第一集我们做了两件事：
\begin{enumerate}[leftmargin=1.4em,itemsep=0.2em]
\item 证明 $e^{-\ep|t|}\to1$ 于 $\Sv'$（用 $|\varphi|$ 控制，控制收敛）；
\item 算出一族频谱 $G_\ep\to2\pi\delta$。
\end{enumerate}

但“时域的极限”和“频域的极限”是两条独立的陈述，中间缺一座桥：\textbf{为什么时域极限的变换，等于变换的极限？}

命题 \ref{prop:F-cont} 就是这座桥。有了它：
\begin{equation}
e^{-\ep|t|}\to1\ \text{于}\ \Sv'
\quad\Longrightarrow\quad
\F\big[e^{-\ep|t|}\big]\to\F[1]
\quad\Longrightarrow\quad
2\pi\delta=\F[1] .
\end{equation}
（最后一步用了极限的唯一性。）

\textbf{第一集的结论，到这里才真正被证明。}
\end{ideabox}

\section{统一证明三个主结果}

\subsection{$\delta$ 与常数}

\begin{align}
\ip{\F\delta}{\varphi}&=\ip{\delta}{\ph}=\ph(0)=\int_{\R}\varphi(t)\dd t=\ip{1}{\varphi}
&&\Longrightarrow\quad \F[\delta]=1 .\label{eq:Fdelta}\\[4pt]
\ip{\F1}{\varphi}&=\ip{1}{\ph}=\int_{\R}\ph(t)\dd t=2\pi\varphi(0)=\ip{2\pi\delta}{\varphi}
&&\Longrightarrow\quad \F[1]=2\pi\delta .\label{eq:Fone}
\end{align}

\begin{ideabox}[注意第二个式子用了什么]
它用了 \eqref{eq:inversion} 在 $t=0$ 处的取值：$\varphi(0)=\frac{1}{2\pi}\int\ph(t)\dd t$。

\textbf{这就是为什么反演必须在前面证。}如果反演还没证，\eqref{eq:Fone} 就悬空了。

也是为什么不能用 \eqref{eq:Fone} 去证反演——顺序不能反。
\end{ideabox}

\subsection{两条微分规则（严格版）}

要证 $\F(tT)=i\partial_\omega\widehat T$。关键是先算清 $t\ph(t)$ 与 $\varphi$ 的关系。

由 $\F[\varphi']=i\omega\ph(\omega)$，把变量写成 $t$：

\begin{equation}
(\varphi')^{\wedge}(t)=it\,\ph(t)
\qquad\Longrightarrow\qquad
t\,\ph(t)=-i\,(\varphi')^{\wedge}(t).
\end{equation}

于是

\begin{align}
\ip{\F(tT)}{\varphi}
&=\ip{tT}{\ph}=\ip{T}{t\ph}
=-i\,\ip{T}{(\varphi')^{\wedge}}
=-i\,\ip{\widehat T}{\varphi'}
=i\,\ip{\widehat T'}{\varphi}
=\ip{i\,\partial_\omega\widehat T}{\varphi} .
\end{align}

所以

\begin{equation}
\F(tT)=i\,\partial_\omega\widehat T .\label{eq:mult-t-dist}
\end{equation}

同理，要证 $\F(T')=i\omega\widehat T$：

由 $\ph'(\omega)=-i\,(t\varphi)^{\wedge}(\omega)$，得

\begin{align}
\ip{\F(T')}{\varphi}
&=\ip{T'}{\ph}=-\ip{T}{\ph'}
=-\ip{T}{-i(t\varphi)^{\wedge}}
=i\,\ip{\widehat T}{t\varphi}
=i\,\ip{t\widehat T}{\varphi}
=\ip{i\omega\widehat T}{\varphi} .
\end{align}

所以

\begin{equation}
\F(T')=i\omega\,\widehat T .\label{eq:deriv-dist}
\end{equation}

\begin{ideabox}[这两段证明的“技巧”都在同一步]
把 $t\ph(t)$（或 $\ph'(\omega)$）\textbf{换成某个东西的傅里叶变换}，然后把它推给 $\widehat T$。

这一步依赖两条恒等式：
\begin{itemize}[leftmargin=1.4em,itemsep=0.2em]
\item $t\ph(t)=-i(\varphi')^{\wedge}(t)$；
\item $\ph'(\omega)=-i(t\varphi)^{\wedge}(\omega)$。
\end{itemize}
两个负号配合两个 $i$，最后给出正的 $i\partial_\omega\widehat T$ 与 $i\omega\widehat T$。

\textbf{如果符号算错，这两条规则会同时错。}建议用 $\varphi=e^{-x^{2}/2}$ 数值检查一次。
\end{ideabox}

\subsection{阶跃}

$H_\ep(t)=H(t)e^{-\ep t}\to H$ 于 $\Sv'$（因为 $|H_\ep\varphi|\le|\varphi|\in\Lone$，控制收敛）。

而 $\F[H_\ep]=\dfrac{1}{\ep+i\omega}\to\pi\delta-i\PV\dfrac{1}{\omega}$（EP03 的两条极限）。

由命题 \ref{prop:F-cont}，$\F[H_\ep]\to\F[H]$。由极限唯一性，

\begin{equation}
\F[H]=\pi\delta-i\,\PV\frac{1}{\omega}=\pi\delta+\PV\frac{1}{i\omega} .\label{eq:FH-final}
\end{equation}

\begin{center}
\small
\renewcommand{\arraystretch}{1.4}
\begin{tabularx}{\textwidth}{@{}L{4.4cm}L{4.0cm}Y@{}}
\toprule
结果 & 在哪一集“算出来” & 在哪一集“被证明” \\
\midrule
$\F[1]=2\pi\delta$ & EP01 & EP06（对偶 + 反演）\\
$\F[t^{n}]=2\pi i^{n}\delta^{(n)}$ & EP02 & EP06（\eqref{eq:mult-t-dist} 递推）\\
$\F[H]=\pi\delta+\PV\frac{1}{i\omega}$ & EP03 & EP06（连续性 + 极限唯一）\\
\bottomrule
\end{tabularx}
\end{center}

\section{换一种正则化，结果一样}

理论必须经得起“换一条路”的检验。这一节就是检验。

\subsection{高斯}

\begin{equation}
g_\sigma(t)=e^{-\sigma t^{2}/2},
\qquad
\widehat{g_\sigma}(\omega)=\sqrt{\frac{2\pi}{\sigma}}\;e^{-\omega^{2}/(2\sigma)} .
\end{equation}

（这个变换可以用 EP02 的微分方程法算，见附录 A.1。）

总面积：
\begin{equation}
\int_{\R}\widehat{g_\sigma}\dd\omega=\sqrt{\frac{2\pi}{\sigma}}\cdot\sqrt{2\pi\sigma}=2\pi .
\end{equation}

$\sigma\downarrow0$ 时，$g_\sigma\to1$，而 $\widehat{g_\sigma}$ 变成一个宽度 $\sim\sqrt\sigma$、高度 $\sim\sqrt{2\pi/\sigma}$、面积恒为 $2\pi$ 的尖峰，所以

\begin{equation}
\widehat{g_\sigma}\longrightarrow2\pi\delta .
\end{equation}

这条路线完全独立于 EP01 的指数核，结论一致。

\subsection{矩形窗与 sinc}

\begin{equation}
w_T=\mathbf{1}_{[-T,T]},
\qquad
\widehat{w_T}(\omega)=\frac{2\sin(T\omega)}{\omega} .
\end{equation}

$w_T\to1$ 于 $\Sv'$（控制函数 $|\varphi|$），所以

\begin{equation}
\widehat{w_T}\longrightarrow\F[1]=2\pi\delta
\qquad\text{于}\ \Sv' .
\end{equation}

\begin{warnbox}[但 sinc 不能套用 EP04 的定理]
$\dfrac{2\sin(T\omega)}{\omega}$：

\begin{itemize}[leftmargin=1.4em,itemsep=0.25em]
\item \textbf{不正}：在原点两侧无限次变号；
\item \textbf{不绝对可积地集中}：$\displaystyle\int_{\R}\Big|\frac{2\sin(T\omega)}{\omega}\Big|\dd\omega$ 按 $\ln T$ 增长，无界；
\item \textbf{不逐点收敛到 $\delta$}：它在每个 $\omega\neq0$ 处仍然振荡（这正是 EP01 的 $I_T$）。
\end{itemize}

所以 EP04 的“非负核 + 质量守恒 + 尾部集中”这套条件它一条都不满足。它之所以仍然收敛，靠的是\textbf{另一条逻辑}：$w_T\to1$ 于 $\Sv'$，加上 $\F$ 连续。

\textbf{这句对比很重要：}EP01 说“$I_T$ 不收敛”，EP06 说“$\widehat{w_T}$ 收敛到 $2\pi\delta$”——两句话不矛盾，因为前者说的是点值/截断积分，后者说的是分布。
\end{warnbox}

\begin{ideabox}[三条路，一个终点]
\begin{center}
\small
\begin{tabular}{@{}lll@{}}
\toprule
正则化 & 频域形状 & 共同极限 \\
\midrule
指数 $e^{-\ep|t|}$ & Cauchy 型，永远非负 & $2\pi\delta$ \\
高斯 $e^{-\sigma t^{2}/2}$ & 高斯型，永远非负 & $2\pi\delta$ \\
矩形 $\mathbf{1}_{[-T,T]}$ & sinc 型，无限变号 & $2\pi\delta$ \\
\bottomrule
\end{tabular}
\end{center}

三种完全不同的形状，同一个终点。\textbf{这就是“$\delta$ 由作用定义，不由形状定义”的最终证据。}第一集挖的坑，到这里填平。
\end{ideabox}

\section{收官：三个层次}

这一集结束时，要把整条路重新走一遍，但要快，而且要用观众已经掌握的语言。

\begin{center}
\small
\renewcommand{\arraystretch}{1.5}
\begin{tabularx}{\textwidth}{@{}L{3.2cm}L{5.0cm}Y@{}}
\toprule
层次 & 内容 & 一句话 \\
\midrule
普通积分 & $f\in\Lone$ 时 $\displaystyle\int f(t)e^{-i\omega t}\dd t$ 收敛 & “衰减够快，一切正常。”\\
测试函数的观察方式 & 用 $\ip{T}{\varphi}$ 代替“逐点取值” & “不知道它是什么，但知道它做什么。”\\
广义傅里叶变换 & 在 $\Sv'$ 上 $\ip{\F T}{\varphi}=\ip{T}{\F\varphi}$ & “同一个变换，更大的空间。”\\
\bottomrule
\end{tabularx}
\end{center}

\begin{ideabox}[最后要说的一句话]
我们\textbf{没有}把一个发散的积分硬塞一个值。

我们做的是：先把“信号”重新理解为“对探针的作用”，再在这个更大的空间里定义一个连续、可逆的傅里叶变换。

在这个框架里，常数、多项式、阶跃的公式\textbf{同时}成立，而且彼此不冲突。

\textbf{这不是给旧答案打补丁，这是换了一套更合适的语言。}
\end{ideabox}

\section{讲解禁区（EP06）}

\begin{warnbox}[逐条对照]
\begin{enumerate}[leftmargin=1.6em,itemsep=0.28em]
\item 不说“对任意分布都能做傅里叶变换”。只在 $\Sv'$ 上成立（$\D'$ 上未必）。
\item 不用 $\int e^{i\omega(t-s)}\dd\omega=2\pi\delta(t-s)$ 去证反演。
\item 说“$\Sv$ 与 $\Sv'$”时，不要把 $\Sv$（探针）和 $\Sv'$（对象）混用。
\item 不说“$\Sv$ 就是衰减很快的函数”，要强调\textbf{每阶导数}都快于任意负幂。
\item 不说“$1/(1+x^{2})$ 在 $\Sv$ 里”，它只是可积。
\item 不省略对偶公式里“无共轭、无 $1/(2\pi)$”的说明；也不要与含共轭的约定混搭。
\item 不说“$\F$ 是 $\Sv'$ 上的酉算子”。它是同构，但配对不是内积，别用“酉”这个词。
\item 不把 sinc 核当成非负近似核来证明；走“$w_T\to1$ + 连续性”这条路。
\item 不说“$e^{x^{2}}$ 也能做变换，只是结果很怪”；它不在 $\Sv'$ 里。
\item 不说“反演是显然的”；至少要给出高斯正则化的思路。
\item 不把 \eqref{eq:Fone} 与 \eqref{eq:inversion} 的证明顺序颠倒。
\end{enumerate}
\end{warnbox}

\begin{scriptbox}[最后 40 秒的旁白骨架]
“回到最开始那条水平线。

我们当时说，它的频谱算不出来。

现在我们知道：算不出来的，是那个积分；不是那个频谱。

换一种观察方式，换一个更大的空间，那条水平线依然有它的频谱——就在零频率上，权重是 $2\pi$。

而我们付出的代价，只是承认一件事：

一个东西是什么，可以不看它在每一点有多高，而看它对每一个问题怎么回答。”
\end{scriptbox}

\section{一页纸回顾：三条公式各靠什么成立}

\begin{center}
\small
\renewcommand{\arraystretch}{1.42}
\begin{tabularx}{\textwidth}{@{}L{1.0cm}L{4.6cm}L{4.4cm}Y@{}}
\toprule
步骤 & 内容 & 关键工具 & 缺了它会怎样 \\
\midrule
1 & 定义 $\Sv$ 与 $\Sv'$ & 半范数、连续线性 & 没有能容纳常数的空间 \\
2 & 证 $\F:\Sv\to\Sv$ 连续 & 两条微分恒等式 & 对偶定义右边无意义 \\
3 & 高斯正则化证反演 & 高斯积分 + Fubini + 近似核 & $\F[1]=2\pi\delta$ 悬空 \\
4 & 对偶定义 + 连续性 & $\ip{\F T}{\varphi}=\ip{T}{\F\varphi}$ & 无法把极限搬进变换 \\
5 & 三条主结果 & $1$--$4$ 的组合 & 只能逐条“凑” \\
\bottomrule
\end{tabularx}
\end{center}

\begin{notebox}[一句话总结这一集的技术路线]
\textbf{换探针（$\D\to\Sv$）$\Rightarrow$ 有定义（对偶）$\Rightarrow$ 有反演（高斯）$\Rightarrow$ 有连续性（保持极限）$\Rightarrow$ 三条结果同时落地。}

五个箭头，一条都不能少。
\end{notebox}

\begin{quizbox}[EP06 自测]
\begin{enumerate}[leftmargin=1.6em,itemsep=0.35em]
\item 证明命题 \ref{prop:F-not-D}，并说明为什么“整函数在开集上为零则处处为零”是关键。
\item 验证 $\omega^{m}\widehat{g}(\omega)=i^{-m}\F[g^{(m)}](\omega)$。
\item 用 \eqref{eq:Schwartz-bound} 说明 $\F$ 把 $\Sv$ 映入 $\Sv$。
\item 复述高斯正则化反演的四个步骤，指出 Fubini 用在哪里、近似恒等核用在哪里。
\item 由 $\F^{2}\varphi=2\pi\varphi(-t)$ 推出 $\F[\delta]=1$ 与 $\F[1]=2\pi\delta$ 的关系。
\item 直接由对偶定义证明 $\F[1]=2\pi\delta$（不用正则化核），并指出哪一步用了反演。
\item 证明 $e^{x^{2}}\notin\Sv'$。
\item 说明为什么 sinc 核不满足 EP04 的条件，却仍在 $\Sv'$ 中收敛。
\end{enumerate}
\end{quizbox}
}{}
\IfFileExists{ch7.tex}{\chapter{把六集缝起来：总回顾与制作落地}

\begin{ideabox}[这一章的作用]
前面六章是“分集精讲”。这一章把六集重新压缩成一张可以随时查阅的总图，并且把数学结论翻译成\textbf{制作动作}：哪一集必须出现哪个画面、哪句话必须说、哪个坑必须躲。
\end{ideabox}

\section{一张总表}

\begin{center}
\small
\renewcommand{\arraystretch}{1.45}
\begin{tabularx}{\textwidth}{@{}L{1.0cm}L{3.4cm}L{4.2cm}L{3.4cm}Y@{}}
\toprule
集 & 引入的新对象 & 本集真正建立的东西 & 借用了什么 & 留下什么 \\
\midrule
01 & $\delta$（作为极限的名字） & 经典积分失败的现场；双边指数正则化；面积集中与取样 & 取样极限（未证） & $\delta$ 是什么？\\
02 & $\delta'$、$\delta^{(n)}$ & 乘 $t$ 对应频域求导；分布导数的负号 & 分布导数规则 & 求导对应什么？\\
03 & $\PV\frac1\omega$ & 微分规则会丢零频项；主值；正则化 & $H'=\delta$ & 这些“趋于”是什么意思？\\
04 & 近似恒等核 & 三种收敛；控制收敛；取样极限的完整证明 & — & 能否用作用定义对象？\\
05 & 分布、测试函数 & 连续线性泛函；嵌入；运算与禁区 & — & 怎么对分布做变换？\\
06 & $\Sv$、$\Sv'$ & 对偶定义；反演；连续性；三结果统一 & — & （闭环）\\
\bottomrule
\end{tabularx}
\end{center}

\section{依赖树：谁支撑谁}

\begin{figure}[htbp]
\centering
\begin{tikzpicture}[
  every node/.style={font=\scriptsize,align=center},
  box/.style={draw=mainblue!55,fill=softblue,rounded corners=1.5pt,
              inner sep=3pt,minimum height=1.4cm,text width=2.25cm},
  ar/.style={-{Stealth[length=1.5mm]},draw=notegray,thick},
  back/.style={-{Stealth[length=1.6mm]},draw=warnred!80,dashed,thick}
]
\node[box] (e1) at (0,0)     {\textbf{EP01}\\[1pt]$e^{-\ep|t|}\!\to\!1$\\$G_\ep\to2\pi\delta$};
\node[box] (e2) at (2.70,0)  {\textbf{EP02}\\[1pt]$K_\ep'\!\to\!\delta'$\\$\F[t^{\,n}]$};
\node[box] (e3) at (5.40,0)  {\textbf{EP03}\\[1pt]$H_\ep\to H$\\$1/(\ep\!+\!i\omega)$};
\node[box] (e4) at (8.10,0)  {\textbf{EP04}\\[1pt]三种收敛\\取样极限};
\node[box] (e5) at (10.80,0) {\textbf{EP05}\\[1pt]定义 $\delta$\\定义 $\delta'$};
\node[box] (e6) at (13.50,0) {\textbf{EP06}\\[1pt]$\Sv$ 与 $\Sv'$\\对偶定义};

\draw[ar] (e1) -- (e2);
\draw[ar] (e2) -- (e3);
\draw[ar] (e3) -- (e4);
\draw[ar] (e4) -- (e5);
\draw[ar] (e5) -- (e6);

\draw[back] ($(e6.north)+(0,0.65)$) to[out=160,in=20]
   node[above,font=\scriptsize,text=warnred!80,pos=0.5] {EP06 回收并结清 EP01--EP03 的全部结论}
   ($(e1.north)+(0,0.65)$);
\end{tikzpicture}
\caption{观众看到的顺序（横向箭头）。而理论支撑的方向是\textbf{反的}：EP01--EP03 的结果要到 EP06 才被真正证明；EP02 的求导规则要到 EP05 才被定义。这条反向的虚线，就是整个系列的结构性承诺。}
\end{figure}

\begin{ideabox}[这张图最好的用法]
把它当成\textbf{写作时的良心检查}。

每当你写了一句“所以 $\F[t]=2\pi i\delta'$”，就在心里问：这句话的上游是哪一集？下游哪一集还账？
\begin{itemize}[leftmargin=1.4em,itemsep=0.2em]
\item 如果上游不存在（例如第一集就用反演），那就是硬伤。
\item 如果下游没说清楚（例如从未提到“后面会证”），观众会觉得你在糊弄。
\end{itemize}
\end{ideabox}

\section{每一集的三条底线}

所谓“底线”，指的是：\textbf{即使这一集剪短、删掉一半内容，这三件事也必须留下。}

\begin{center}
\small
\renewcommand{\arraystretch}{1.5}
\begin{tabularx}{\textwidth}{@{}L{1.1cm}L{4.6cm}L{5.2cm}Y@{}}
\toprule
集 & 底线一 & 底线二 & 底线三 \\
\midrule
01 & 常数 $=$ 零频率 & 普通积分确实算不出来（不是“懒得算”） & $\ep$ 变小时面积不变、峰变窄变高 \\
02 & 乘 $t$ $\leftrightarrow$ 频域求导 & $\delta'$ 的作用是 $-\varphi'(0)$ & $\delta'$ 不是单峰，不能画成一根箭头 \\
03 & $H'=\delta$ & $\omega\delta=0$ 导致解不唯一 & 主值 $=$ 对称挖掉原点 \\
04 & 逐点 $\neq$ 分布 & 控制收敛需要公共可积控制 & 形状不同、极限相同 \\
05 & 分布 $=$ 连续线性泛函 & $\delta$ 不是局部可积函数 & 两个分布不能随便相乘 \\
06 & $\F$ 不保持 $\D$，保持 $\Sv$ & $\ip{\F T}{\varphi}=\ip{T}{\F\varphi}$ & 三条主结果由此同时成立 \\
\bottomrule
\end{tabularx}
\end{center}

\section{伏笔与回收：叙事上的六个扣}

数学上讲通，不等于片子好看。片子的节奏靠“扣”：埋下去，吊着，最后解开。

\begin{center}
\small
\renewcommand{\arraystretch}{1.45}
\begin{tabularx}{\textwidth}{@{}L{1.2cm}L{4.4cm}L{4.4cm}Y@{}}
\toprule
扣 & 在哪一集埋 & 埋的时候说什么 & 在哪一集解 \\
\midrule
A & EP01 & “这个等号是广义意义下的，严格定义我们后面再说。” & EP05、EP06 \\
B & EP01 & “它没有消失，只是挤到了一起。” & EP04 \\
C & EP02 & “对一个不是函数的东西求导，是什么意思？” & EP05 \\
D & EP03 & “这个方程有无穷多个解，我们漏掉了什么？” & EP03 后半（部分）＋EP06（完整）\\
E & EP04 & “不同的核说了同一句话，为什么不直接用这句话定义它？” & EP05 \\
F & EP05 & “傅里叶变换是一个积分，可它没有积分。” & EP06 \\
\bottomrule
\end{tabularx}
\end{center}

\begin{warnbox}[关于“扣”的两条纪律]
\begin{enumerate}[leftmargin=1.4em,itemsep=0.3em]
\item \textbf{扣必须真的解开。}如果 EP06 只讲“对偶定义”而没有回头把 EP01 的 $2\pi$、EP02 的 $i^{n}$、EP03 的 $\pi$ 一个个对上，观众会觉得被吊了五集。
\item \textbf{扣不能太多。}六集六个扣已经是上限。再加就会变成“每集都在许诺，每集都不兑现”。
\end{enumerate}
\end{warnbox}

\section{从数学到分镜：一张转换表}

写分镜时最容易出的问题是：\textbf{把数学结构直接当成画面结构}。数学上的“一步”不一定是一个镜头；数学上的“显然”可能需要三十秒。

\begin{center}
\small
\renewcommand{\arraystretch}{1.45}
\begin{tabularx}{\textwidth}{@{}L{4.6cm}L{4.8cm}Y@{}}
\toprule
数学上的一步 & 画面上的处理 & 为什么 \\
\midrule
通分 $\frac{1}{\ep-i\omega}+\frac{1}{\ep+i\omega}$ & 3--4 次显现，配颜色对应左右半轴 & 这是全片唯一要求观众跟的代数 \\
$\int K_\ep\varphi\to\varphi(0)$ & 一个镜头 + 阴影收缩 + 读数 & 结论可以“看”，证明不必“看” \\
$\omega\delta=0$ & 单独停 5 秒，配一行字 & 这是整集的转折点，不能滑过去 \\
换元 $x=\ep u$ & 可省略，或只给一个静态对照 & 它是证明技巧，不是叙事节点 \\
$\Sv$ 的半范数 & 只给一句“每阶导数都比任何负幂快” & 半范数是工具，不是故事 \\
反演的高斯正则化 & 给四步流程，不逐行推导 & 观众需要“这条路是干净的”，不需要每个估计 \\
\bottomrule
\end{tabularx}
\end{center}

\subsection*{制作顺序（沿用你的既定流程，补上数学检查点）}

\begin{enumerate}[leftmargin=2em,itemsep=0.5em]
\item 用本书对应章节确认这一集的数学主线与底线。
\item 写旁白草稿，自己朗读一遍并录临时音轨。\textbf{检查点：}每出现一个“趋于”，旁边是否交代了是哪种趋于。
\item 把旁白拆成场景，按动画编号组织。\textbf{检查点：}每个“扣”是否有对应的画面，而不只是台词。
\item 低画质 Manim 预览。\textbf{检查点：}峰是否越界、采样是否够密、坐标缩放是否可见、面积标签是否写了“全实轴”。
\item 剪映里按旁白排画面，决定停顿。\textbf{检查点：}三次“预期落空”是否都留了呼吸。
\item 粗剪通过后高清渲染、正式录音、修字幕。\textbf{检查点：}$2\pi$、$i$、负号、$\PV$ 是否与本书约定一致。
\end{enumerate}

\section{四类检查清单}

每一集剪完，按四类各过一遍。前两类是“不能错”，后两类是“能更好”。

\subsection*{一、数学正确性}

\begin{itemize}[leftmargin=1.6em,itemsep=0.25em]
\item[$\square$] 变换约定统一（$\omega$，$2\pi$ 在逆变换上）。
\item[$\square$] $\delta$、$\delta^{(n)}$、$\PV$ 出现时都声明了广义意义。
\item[$\square$] $\delta^{(n)}$ 的 $(-1)^{n}$ 没丢。
\item[$\square$] $\F[H]$ 的 $\pi$ 与 $-i$ 符号正确。
\item[$\square$] 说“趋于”时指明了逐点 / 一致 / 分布。
\item[$\square$] 没有把“面积”当“能量”，没有写 $\delta^{2}$。
\item[$\square$] 没有用未证的结论去证前提（尤其反演）。
\end{itemize}

\subsection*{二、教学正确性}

\begin{itemize}[leftmargin=1.6em,itemsep=0.25em]
\item[$\square$] 全片有一个核心问题贯穿。
\item[$\square$] 每个新对象都由前面的具体困难引出，不是“现在我们来定义”。
\item[$\square$] 观众知道某条规则是“已证明”“直观解释”还是“暂时借用”。
\item[$\square$] 结尾的悬念来自本集真实留下的问题，不是硬造。
\item[$\square$] 至少有一次“预期落空”。
\end{itemize}

\subsection*{三、影像正确性}

\begin{itemize}[leftmargin=1.6em,itemsep=0.25em]
\item[$\square$] 尖峰采样足够密，峰顶没有越界。
\item[$\square$] 坐标缩放时，刻度变化可见。
\item[$\square$] 有限画框内的面积没有被标成全实轴总面积。
\item[$\square$] 普通曲线、分布示意、数值近似三者视觉上可区分。
\item[$\square$] 箭头只表示权重，不表示高度。
\item[$\square$] 实余弦的频谱画了正负两个频率。
\end{itemize}

\subsection*{四、节奏}

\begin{itemize}[leftmargin=1.6em,itemsep=0.25em]
\item[$\square$] 每 2--3 分钟有一次画面变化（不是连续念公式）。
\item[$\square$] 关键转折处留了停顿（$\omega\delta=0$、主值、对偶定义）。
\item[$\square$] 公式不是一屏出完，而是分层显现。
\item[$\square$] 时长达标（8--18 分钟区间），且没有为了凑时长而重复。
\end{itemize}

\section{如果要扩展这个系列}

当前六集已经形成闭环。若要继续，下面几条是自然的延伸，而且都能直接接在 EP06 的结尾：

\begin{center}
\small
\renewcommand{\arraystretch}{1.45}
\begin{tabularx}{\textwidth}{@{}L{3.4cm}L{5.4cm}Y@{}}
\toprule
方向 & 核心问题 & 与本书的接口 \\
\midrule
卷积与滤波 & 为什么时域卷积等于频域相乘？ & EP04 的 $\delta*f=f$ 与 EP06 的对偶 \\
L$^{2}$ 与 Plancherel & 为什么能量守恒？常数为什么不在里面？ & EP01 里“面积不是能量”的伏笔 \\
偏微分方程 & 用傅里叶变换把微分方程变成代数方程 & EP02、EP06 的两条微分规则 \\
格林函数 & 为什么“冲激响应”能决定一切？ & EP03 的 $H'=\delta$ 与 EP04 的 $\delta*f=f$ \\
\bottomrule
\end{tabularx}
\end{center}

\begin{notebox}[扩展的纪律]
不要在 EP06 的结尾突然展开这些。EP06 的任务是\textbf{关门}，不是\textbf{再开一扇门}。

如果真要做续集，把它做成独立的第二季，用一个全新的“第一集问题”开场（例如“为什么滤波器的输出是输入的卷积？”），而不是接在 $\Sv'$ 后面继续说。
\end{notebox}

\section{关键数字速查}

做动画时最常要临时确认的几个量，集中放在这里。

\begin{center}
\small
\renewcommand{\arraystretch}{1.4}
\begin{tabularx}{\textwidth}{@{}L{4.2cm}L{4.0cm}Y@{}}
\toprule
量 & 表达式 & 出现在哪一集 \\
\midrule
$G_\ep$ 的峰高 & $2/\ep$ & EP01，双面板纵轴 \\
$G_\ep$ 的半高全宽 & $2\ep$ & EP01，宽度对比 \\
$G_\ep$ 的全实轴面积 & $2\pi$ & EP01，面积标签 \\
$K_\ep$ 的原点值 & $1/(\pi\ep)$ & EP01，$\ep=0.1$ 时约 $3.18$ \\
$K_\ep'$ 的峰值 & $\approx0.207/\ep^{2}$ & EP02，左正右负两叶 \\
$K_\ep'$ 的极值位置 & $x=\pm\ep/\sqrt3$ & EP02，两叶收缩 \\
$K_\ep'$ 的总积分 & $0$ & EP02，“面积不能定义 $\delta'$” \\
$\F[H_\ep]$ 实部峰高 & $1/\ep$ & EP03，实部面板 \\
$\F[H_\ep]$ 实部面积 & $\pi$ & EP03，$\pi\delta$ 的来源 \\
$\F[H_\ep]$ 虚部极值 & $\mp1/(2\ep)$ & EP03，奇对称结构 \\
$\widehat{g_\sigma}$ 的峰高 & $\sqrt{2\pi/\sigma}$ & EP06，高斯路线的对照 \\
$\widehat{w_T}$ 的零点 & $\omega=k\pi/T$（$k\neq0$ 整数） & EP06，sinc 的振荡 \\
\bottomrule
\end{tabularx}
\end{center}

\begin{notebox}[怎么用这张表]
它的主要用途\textbf{不是背，而是检查}：
\begin{itemize}[leftmargin=1.4em,itemsep=0.2em]
\item 动画里出现的峰高，是否等于表里的数？
\item 面积标签写的数，是否与画的曲线一致？
\item 缩放纵轴之后，刻度有没有跟着变？
\end{itemize}
三个问题各过一次，图上八成的错就没了。
\end{notebox}

\section{最后：这一系列到底给了观众什么}

\begin{ideabox}[三句话]
\begin{enumerate}[leftmargin=1.6em,itemsep=0.35em]
\item \textbf{一个边界。}普通傅里叶积分只对衰减足够快的信号有效；越界之后，它给出的不是“错答案”，而是“没有答案”。
\item \textbf{一种看法。}一个对象可以不看“每点多高”，而看“对每个问题怎么回答”。这个看法一旦建立，$\delta$、$\delta'$、$\PV$ 都不再神秘。
\item \textbf{一次统一。}在 $\Sv'$ 上，傅里叶变换是一个连续可逆的整体；常数、多项式、阶跃的公式不是三套技巧，而是同一个定义的三条推论。
\end{enumerate}

如果观众看完只记住一句话，希望是这一句：

\begin{center}
\textbf{“算不出来”不等于“没有”，只等于“你用的尺子不够长”。}
\end{center}
\end{ideabox}

\vspace{0.6cm}
\begin{center}
\begin{tcolorbox}[colback=softblue,colframe=mainblue!45,boxrule=0pt,arc=2pt,width=0.90\textwidth]
\centering
\vspace{0.25em}
{\small\color{notegray} 到这里，六集的知识地图就画完了。}\\[0.45em]
{\normalsize 剩下的事只有一件：把它压进 8 到 18 分钟里。}
\vspace{0.25em}
\end{tcolorbox}
\end{center}
}{}
\IfFileExists{appendix.tex}{\appendix

\chapter{公式与变换总表}

\begin{ideabox}[使用说明]
本节所有公式\textbf{全部按本书约定}：$\fh(\omega)=\int f(t)e^{-i\omega t}\dd t$，逆变换带 $1/(2\pi)$，配对不含复共轭。

如果你在别处看到一个“差不多但差一个 $2\pi$”的公式，先回到第 0.6 节核对约定，不要直接改系数。
\end{ideabox}

\section{变换对总表}

\begin{center}
\small
\renewcommand{\arraystretch}{1.55}
\begin{tabularx}{\textwidth}{@{}L{4.6cm}L{4.6cm}Y@{}}
\toprule
时域对象 & 频域对象 & 出处 / 备注 \\
\midrule
$\delta(t)$ & $1$ & EP06，由对偶定义一步得出 \\
$1$ & $2\pi\delta(\omega)$ & EP01，主结果 \\
$t^{\,n}$ & $2\pi i^{\,n}\delta^{(n)}(\omega)$ & EP02，递推 \\
$\delta^{(n)}(t)$ & $(i\omega)^{n}$ & EP06，$\ip{\delta^{(n)}}{e^{-i\omega t}}$ \\
$\delta(t-a)$ & $e^{-ia\omega}$ & 平移 \\
$e^{i\omega_0 t}$ & $2\pi\delta(\omega-\omega_0)$ & 调制 \\
$\cos(\omega_0 t)$ & $\pi\big[\delta(\omega-\omega_0)+\delta(\omega+\omega_0)\big]$ & 两个频率，不能只画一个 \\
$\sin(\omega_0 t)$ & $(\pi/i)\big[\delta(\omega-\omega_0)-\delta(\omega+\omega_0)\big]$ & 见下方注 \\
$H(t)$ & $\pi\delta(\omega)-i\,\PV(1/\omega)$ & EP03，主结果 \\
$\sgn t$ & $-2i\,\PV(1/\omega)$ & EP03，由 $H=\frac12+\frac12\sgn$ \\
$\PV(1/t)$ & $-i\pi\,\sgn\omega$ & 与上一行互为对偶 \\
$e^{-\ep|t|}$ & $2\ep/(\ep^{2}+\omega^{2})$ & EP01，主计算 \\
$H(t)e^{-\ep t}$ & $1/(\ep+i\omega)$ & EP03，主计算 \\
$e^{-at^{2}/2}$ & $\sqrt{2\pi/a}\,e^{-\omega^{2}/(2a)}$ & $a>0$，附录 A.2 \\
\bottomrule
\end{tabularx}
\end{center}

\begin{notebox}[关于 $\sin$ 那一段的符号]
用 $\sin(\omega_0t)=\dfrac{1}{2i}\big(e^{i\omega_0t}-e^{-i\omega_0t}\big)$ 现推一遍最稳：

\begin{equation}
\F[\sin(\omega_0 t)]
=\frac{1}{2i}\Big[2\pi\delta(\omega-\omega_0)-2\pi\delta(\omega+\omega_0)\Big]
=\frac{\pi}{i}\big[\delta(\omega-\omega_0)-\delta(\omega+\omega_0)\big].
\end{equation}

有的表把它写成 $\pi i[\delta(\omega-\omega_0)-\delta(\omega+\omega_0)]$，那是与 $e^{+i\omega t}$ 约定配套的写法。\textbf{两种写法都“看着像”，符号却相反}，所以不要凭记忆抄，用 $\frac{1}{2i}(e^{i}-e^{-i})$ 现推。

同理，$\cos$ 与 $\sin$ 的系数都是 $\pi$（不是 $2\pi$），因为两者各是两个指数项的一半。
\end{notebox}

\section{运算规则}

\begin{center}
\small
\renewcommand{\arraystretch}{1.5}
\begin{tabularx}{\textwidth}{@{}L{6.4cm}L{5.4cm}Y@{}}
\toprule
规则 & 表达式 & 条件 / 出处 \\
\midrule
时间平移 & $\F[f(t-a)](\omega)=e^{-ia\omega}\fh(\omega)$ & 处处成立（在相应空间内） \\
频率平移 & $\F[e^{i\omega_0t}f(t)](\omega)=\fh(\omega-\omega_0)$ & 调制 \\
尺度 & $\F[f(at)](\omega)=\dfrac{1}{|a|}\fh\!\left(\dfrac{\omega}{a}\right)$ & $a\neq0$ \\
乘 $t$ & $\F[tf]=i\,\partial_\omega\fh$ & EP02；$f\in\Sv$ \\
求导 & $\F[f']=i\omega\fh$ & 分部积分，边界项消失 \\
分布的乘 $t$ & $\F(tT)=i\,\partial_\omega\widehat T$ & EP06，\eqref{eq:mult-t-dist} \\
分布的求导 & $\F(T')=i\omega\widehat T$ & EP06，\eqref{eq:deriv-dist} \\
两次变换 & $\F^{2}T=2\pi\,T(-\cdot)$ & EP06，\eqref{eq:F2} 的分布版 \\
反演 & $T=\dfrac{1}{2\pi}\F^{2}\big(T(-\cdot)\big)$ & $\Sv'$ 上恒成立 \\
\midrule
$\delta$ 的尺度 & $\delta(ax)=\dfrac{1}{|a|}\delta(x)$ & $a\neq0$，注意绝对值 \\
$\delta$ 的复合 & $\delta(g(x))=\sum_j\dfrac{\delta(x-x_j)}{|g'(x_j)|}$ & $g$ 只有简单零点 \\
$\delta$ 的乘法 & $a\,\delta=a(0)\,\delta$ & $a$ 光滑 \\
乘 $x$ & $x\,\delta^{(n)}=-n\,\delta^{(n-1)}$ & 特别 $x\delta=0$，$x\delta'=-\delta$ \\
Leibniz & $(aT)'=a'T+aT'$ & $a$ 光滑 \\
跳跃 & $Df=f'_{\text{经典}}+\sum_j J_j\delta_{a_j}$ & $f$ 分段 $C^1$ \\
\bottomrule
\end{tabularx}
\end{center}

\section{常用积分}

\begin{center}
\small
\renewcommand{\arraystretch}{1.5}
\begin{tabularx}{\textwidth}{@{}L{6.0cm}L{4.0cm}Y@{}}
\toprule
积分 & 结果 & 用途 \\
\midrule
$\displaystyle\int_{0}^{\infty}e^{-at}\dd t$ & $1/a$（$\Re a>0$） & 指数正则化 \\
$\displaystyle\int_{\R}\frac{\dd\omega}{\ep^{2}+\omega^{2}}$ & $\dfrac{\pi}{\ep}$ & EP01 面积 \\
$\displaystyle\int_{\R}\frac{\ep\,\dd\omega}{\ep^{2}+\omega^{2}}$ & $\pi$ & EP03 实部的面积 \\
$\displaystyle\int_{\R}\frac{2\ep\,\dd\omega}{\ep^{2}+\omega^{2}}$ & $2\pi$ & EP01 频谱面积 \\
$\displaystyle\int_{\R}K_\ep\dd\omega$ & $1$ & 归一化核 \\
$\displaystyle\int_{-a}^{a}K_\ep\dd\omega$ & $\dfrac{2}{\pi}\arctan\dfrac{a}{\ep}$ & 集中性 \\
$\displaystyle\int_{\R}e^{-x^{2}/2}\dd x$ & $\sqrt{2\pi}$ & 高斯 \\
$\displaystyle\int_{\R}e^{-\eta\omega^{2}/2}e^{i\omega x}\dd\omega$ & $\sqrt{\dfrac{2\pi}{\eta}}\,e^{-x^{2}/(2\eta)}$ & EP06 反演 \\
$\displaystyle\int_{-T}^{T}e^{-i\omega t}\dd t$ & $\dfrac{2\sin(\omega T)}{\omega}$（$\omega\neq0$） & EP01 失败现场 \\
$\displaystyle\int_{\R}\Big|\frac{\sin(T\omega)}{\omega}\Big|\dd\omega$ & $\sim\ln T$ & sinc 不可积 \\
\bottomrule
\end{tabularx}
\end{center}

\subsection*{高斯变换：两种算法}

\begin{enumerate}[leftmargin=2em,itemsep=0.5em]
\item \textbf{微分方程法。}设 $g_a(t)=e^{-at^{2}/2}$，$G_a=\F[g_a]$。由 $g_a'=-at\,g_a$，两边变换并用 $\F[g']=i\omega\hat g$ 与 $\F[tg]=i\hat g'$：
\begin{equation}
i\omega G_a=-a i\,G_a'
\qquad\Longrightarrow\qquad
G_a'=-\frac{\omega}{a}G_a .
\end{equation}
解这个常微分方程得 $G_a(\omega)=G_a(0)e^{-\omega^{2}/(2a)}$，而 $G_a(0)=\int e^{-at^{2}/2}\dd t=\sqrt{2\pi/a}$。所以
\begin{equation}
\F\big[e^{-at^{2}/2}\big](\omega)=\sqrt{\frac{2\pi}{a}}\,e^{-\omega^{2}/(2a)} .
\end{equation}
\item \textbf{配方法（直接算）。}把 $e^{-at^{2}/2-i\omega t}$ 的指数配方成 $-\frac{a}{2}\left(t+\frac{i\omega}{a}\right)^{2}-\frac{\omega^{2}}{2a}$，再用 $\int e^{-x^{2}/2}\dd x=\sqrt{2\pi}$。
\end{enumerate}

\begin{ideabox}[为什么这里可以放心用“经典”方法]
$g_a$ 光滑、快速衰减，所以 $\F[g_a']=i\omega G_a$、$\F[tg_a]=iG_a'$ 都是经典意义下成立的，\textbf{没有借用任何关于分布或反演的结论}。

这正是 EP06 反演证明能“不循环”的前提：高斯这一块是干净的。
\end{ideabox}

\section{对象住在哪个空间}

\begin{center}
\small
\renewcommand{\arraystretch}{1.35}
\begin{tabularx}{\textwidth}{@{}L{3.2cm}L{1.4cm}L{1.4cm}L{1.4cm}L{1.4cm}Y@{}}
\toprule
对象 & $\Lone$ & $\Ltwo$ & $\D'$ & $\Sv'$ & 说明 \\
\midrule
$e^{-\ep|t|}$ & \checkmark & \checkmark & \checkmark & \checkmark & 最乖 \\
$1$ & $\times$ & $\times$ & \checkmark & \checkmark & 有界不衰减 \\
$t^{\,n}$ & $\times$ & $\times$ & \checkmark & \checkmark & 多项式增长 \\
$H$ & $\times$ & $\times$ & \checkmark & \checkmark & 有界，两端不衰减 \\
$\delta,\ \delta^{(n)}$ & $\times$ & $\times$ & \checkmark & \checkmark & 泛函，不是函数 \\
$\PV\frac1\omega$ & $\times$ & $\times$ & \checkmark & \checkmark & 奇分布 \\
$1/(1+t^{2})$ & \checkmark & \checkmark & \checkmark & $\times$ & 可积但衰减不够快 \\
$e^{t^{2}}$ & $\times$ & $\times$ & \checkmark & $\times$ & 有分布，无傅里叶变换 \\
\bottomrule
\end{tabularx}
\end{center}

\section{三套约定换算速查}

\begin{center}
\small
\renewcommand{\arraystretch}{1.5}
\begin{tabularx}{\textwidth}{@{}L{2.4cm}L{5.2cm}L{3.6cm}Y@{}}
\toprule
约定 & 正变换 & $\F[1]$ & 与本书的关系 \\
\midrule
本书（角频率） & $\displaystyle\int f(t)e^{-i\omega t}\dd t$ & $2\pi\delta(\omega)$ & — \\
对称型 & $\displaystyle\frac{1}{\sqrt{2\pi}}\int f(t)e^{-i\omega t}\dd t$ & $\sqrt{2\pi}\,\delta(\omega)$ & 整体乘 $1/\sqrt{2\pi}$ \\
频率型 & $\displaystyle\int f(t)e^{-2\pi i\nu t}\dd t$ & $\delta(\nu)$ & 令 $\omega=2\pi\nu$ \\
\bottomrule
\end{tabularx}
\end{center}

换算时只需要记住两条：

\begin{equation}
\F_\nu[f](\nu)=\F_\omega[f](2\pi\nu),
\qquad
\delta(2\pi\nu)=\frac{1}{2\pi}\delta(\nu).
\end{equation}

\begin{warnbox}[最容易出错的三处]
\begin{enumerate}[leftmargin=1.4em,itemsep=0.3em]
\item \textbf{$\delta$ 的系数。}换约定之后 $\F[1]$ 不再是 $2\pi\delta$，而是 $\sqrt{2\pi}\delta$ 或 $\delta$。抄别人公式时先核对这一条。
\item \textbf{$\cos$ 与 $\sin$ 的系数。}本书里是 $\pi$（不是 $2\pi$），因为余弦是两个指数项的一半。
\item \textbf{$\sin$ 的符号。}见 A.1 末尾的说明：不同约定下符号会反，用 $\frac{1}{2i}(e^{i}-e^{-i})$ 现推最稳。
\end{enumerate}
\end{warnbox}

\chapter{讲解禁区：能说与不能说}

\begin{ideabox}[怎么用这一章]
写旁白时逐条对照。左列是\textbf{看起来很像对}的说法，中间是它错在哪，右列是可以直接念的说法。

如果时间只够检查三件事，检查这三条：
\begin{enumerate}[leftmargin=1.6em,itemsep=0.2em]
\item 有没有把 $\delta$ 说成“无穷高的函数”；
\item 有没有把“趋于”混成一种；
\item 有没有把 $2\pi$、$i$、负号互换过。
\end{enumerate}
\end{ideabox}

\section{关于 $\delta$ 与尖峰}

\begin{center}
\small
\renewcommand{\arraystretch}{1.45}
\begin{tabularx}{\textwidth}{@{}L{4.9cm}L{4.2cm}Y@{}}
\toprule
不能这样说 & 问题 & 可以这样说 \\
\midrule
“$\delta$ 是原点无穷高、其余为零的函数” & 它没有取值 & “$\delta$ 是一条规则：取出探针在原点的值” \\
“$\delta(0)=\infty$” & $\infty$ 不是数 & “$\delta$ 没有‘在某点的值’” \\
“$\delta$ 的面积是 $1$，所以它是一个峰” & 面积不决定形状 & “$K_\ep$ 的总质量为 $1$，且质量集中到原点” \\
“$K_\ep$ 处处趋于零” & 原点是例外 & “除原点外，每点趋于零；原点趋于无穷” \\
“所以整体趋于零” & 推理错误 & “整体不趋于零，它集中到一点” \\
“$\delta'$ 是一个更尖的峰” & 它有两叶 & “$\delta'$ 有正负两叶，总积分为零” \\
“$\delta'$ 的面积是 $1$” & 总积分为 $0$ & “$\delta'$ 靠正负分布方式起作用” \\
“$\delta$ 的能量是 $2\pi$” & $\delta^2$ 没有定义 & “零频率处的权重是 $2\pi$” \\
\bottomrule
\end{tabularx}
\end{center}

\section{关于极限与收敛}

\begin{center}
\small
\renewcommand{\arraystretch}{1.45}
\begin{tabularx}{\textwidth}{@{}L{4.9cm}L{4.2cm}Y@{}}
\toprule
不能这样说 & 问题 & 可以这样说 \\
\midrule
“非零频率处积分等于零” & 积分不存在 & “它持续振荡，没有极限” \\
“振荡会抵消，所以收敛” & 振荡不是收敛 & “振荡是运动，不是收敛；需要另外的条件” \\
“这族函数趋于 $\delta$，所以逐点也趋于 $\delta$” & 两种收敛不同 & “在分布意义下趋于 $\delta$” \\
“一致收敛就能交换极限与积分” & 无界区间反例 & “还需要公共可积控制” \\
“每一项都可积就够了” & 需要同一个控制 & “需要同一个可积控制函数” \\
“面积是 $1$，所以极限是 $\delta$” & 可能是 $\to0$ & “面积 $1$ 加尾部集中，才趋于 $\delta$” \\
\bottomrule
\end{tabularx}
\end{center}

\section{关于公式与约定}

\begin{center}
\small
\renewcommand{\arraystretch}{1.45}
\begin{tabularx}{\textwidth}{@{}L{4.9cm}L{4.2cm}Y@{}}
\toprule
不能这样说 & 问题 & 可以这样说 \\
\midrule
“$\F[1]=\delta$” & 漏了 $2\pi$ & “$\F[1]=2\pi\delta$（本约定）” \\
“$\F[t^{n}]=2\pi\delta^{(n)}$” & 漏了 $i^{n}$ & “$\F[t^{n}]=2\pi i^{n}\delta^{(n)}$” \\
“$\F[H]=\pi\delta-i\PV\frac{1}{\omega}$ 里符号无所谓” & 符号固定 & “按本约定，是 $\pi\delta-i\PV\frac1\omega$” \\
“$\delta(ax)=\frac1a\delta(x)$” & 漏绝对值 & “$\delta(ax)=\frac{1}{|a|}\delta(x)$” \\
“$\delta(g(x))$ 总能用简单零点公式” & 条件 & “只在 $g$ 的零点都是简单零点时成立” \\
“$\F$ 对所有分布都有定义” & 只在 $\Sv'$ & “在 $\Sv'$ 上成立；$\D'$ 上未必” \\
“$1$ 在 $\Sv$ 里” & 不衰减 & “$1\in\Sv'$，但 $1\notin\Sv$” \\
“$1/(1+x^{2})\in\Sv$” & 衰减不够快 & “它属于 $\Lone$，但不属于 $\Sv$” \\
\bottomrule
\end{tabularx}
\end{center}

\section{关于运算的合法性}

\begin{center}
\small
\renewcommand{\arraystretch}{1.45}
\begin{tabularx}{\textwidth}{@{}L{4.9cm}L{4.2cm}Y@{}}
\toprule
不能这样说 & 问题 & 可以这样说 \\
\midrule
“两个分布可以相乘” & 一般不行 & “可以乘光滑函数；两个分布相乘无通用定义” \\
“$H\delta$ 取 $H(0)=1/2$ 就定义了” & 依赖正则化 & “$H\delta$ 没有内禀定义” \\
“$\omega\delta=0$ 是错的” & 它是对的 & “$\ip{\omega\delta}{\varphi}=\ip{\delta}{\omega\varphi}=0$” \\
“$i\omega\hat H=1$ 能唯一确定 $\hat H$” & 可加 $c\delta$ & “只能定非零频率部分” \\
“$\PV$ 就是忽略奇点” & 规定了对称方式 & “对称地挖掉原点邻域” \\
\bottomrule
\end{tabularx}
\end{center}

\section{如果只查三件事}

\begin{center}
\begin{tcolorbox}[colback=softred,colframe=warnred!55,boxrule=0pt,arc=2pt,width=0.95\textwidth]
\small
\textbf{一查 $\delta$ 的身份。}全片有没有出现过“$\delta$ 在某点等于什么”？

\textbf{二查“趋于”的含义。}每一处“趋于”，是逐点、一致、$L^{p}$，还是分布？

\textbf{三查系数。}$2\pi$ 在哪一边？$i$ 的指数对不对？$(-1)^{n}$ 有没有丢？$\frac{1}{i}=-i$ 用对了吗？
\end{tcolorbox}
\end{center}

这三条过一遍，剩下的错基本都是笔误；这三条错一条，观众学到的就是错的。
\chapter{符号与术语}

\section{符号表}

\begin{center}
\small
\renewcommand{\arraystretch}{1.4}
\begin{tabularx}{\textwidth}{@{}L{3.4cm}Y@{}}
\toprule
记号 & 含义 \\
\midrule
$\F$, $\fh$, $\widehat{T}$ & 傅里叶变换；函数用帽子，分布也用帽子 \\
$\omega$ & 角频率（本书唯一使用的频率变量） \\
$\ep$ & 正则化参数，$\ep\downarrow0$ \\
$H(t)$ & 阶跃函数，$t>0$ 时等于 $1$ \\
$\sgn t$ & 符号函数 \\
$\delta$ & Dirac 分布，$\ip{\delta}{\varphi}=\varphi(0)$ \\
$\delta^{(n)}$ & 第 $n$ 阶分布导数，$\ip{\delta^{(n)}}{\varphi}=(-1)^{n}\varphi^{(n)}(0)$ \\
$\PV\frac{1}{\omega}$ & Cauchy 主值分布 \\
$K_\ep$ & 归一化 Cauchy 核，面积 $1$ \\
$G_\ep$ & 实际频谱 $\F[e^{-\ep|t|}]$，面积 $2\pi$ \\
$\D=C_c^{\infty}$ & 紧支撑光滑测试函数空间 \\
$\D'$ & 分布空间 \\
$\Sv$ & Schwartz 空间（各阶导数快速衰减） \\
$\Sv'$ & 温和分布空间（本系列的工作空间） \\
$\ip{T}{\varphi}$ & 分布对测试函数的复线性作用，\textbf{无共轭} \\
$p_{m,n}(\varphi)$ & $\Sv$ 的半范数，$\sup_x(1+|x|)^{m}|\varphi^{(n)}(x)|$ \\
$\supp T$ & 支撑 \\
$T_n\to T$ 于 $\Sv'$ & 对每个 $\varphi\in\Sv$，$\ip{T_n}{\varphi}\to\ip{T}{\varphi}$ \\
\bottomrule
\end{tabularx}
\end{center}

\section{术语对照}

\begin{center}
\small
\renewcommand{\arraystretch}{1.35}
\begin{tabularx}{\textwidth}{@{}L{3.6cm}L{4.6cm}Y@{}}
\toprule
中文 & 英文 & 一句话说明 \\
\midrule
正则化 & regularization & 先让积分收敛，再让参数消失 \\
近似恒等核 & approximate identity & 质量守恒 + 尾部集中的核族 \\
测试函数 & test function & 用来“提问”的光滑探针 \\
分布 & distribution & $\D$ 上的连续线性泛函 \\
温和分布 & tempered distribution & $\Sv$ 上的连续线性泛函 \\
主值 & Cauchy principal value & 对称挖掉奇点后取的极限 \\
支撑 & support & 不能被局部消掉的区域 \\
控制收敛 & dominated convergence & 交换极限与积分的许可证 \\
半范数 & seminorm & 衡量 $\Sv$ 中“有多小”的尺子 \\
快速衰减 & rapid decay & 比任何负幂都快 \\
\bottomrule
\end{tabularx}
\end{center}

\chapter{练习与详解}

\section{十二道题（覆盖六集）}

\begin{enumerate}[leftmargin=2em,itemsep=0.45em]
\item 对固定 $\omega\neq0$，构造两列 $T\to\infty$，说明 $\dfrac{2\sin(\omega T)}{\omega}$ 没有极限。
\item 计算 $G_\ep$ 在 $[-a,a]$ 上的面积，并解释“固定小区间最终容纳几乎全部面积”。
\item 取 $\varphi(x)=e^{-x^{2}}+xe^{-x^{2}}$，求 $\ip{\delta}{\varphi}$、$\ip{\delta'}{\varphi}$、$\ip{\delta''}{\varphi}$。
\item 求 $\F[3-2t+t^{2}]$，并检查奇偶性与实虚部。
\item 验证 $x\delta'=-\delta$ 与 $x\,\PV\frac1x=1$，并据此检验 $i\omega\widehat H=1$。
\item 证明把 $H(0)$ 改成任意有限值不影响它对应的分布。
\item 为什么 $K_\ep$ 不能被同一个 $\Lone$ 函数控制？换元之后为什么又能用控制收敛？
\item 给出一个面积为 $1$ 但分布极限为零的核族；这与“质量守恒”矛盾吗？
\item 用对偶定义证明 $\F[\delta]=1$ 与 $\F[1]=2\pi\delta$。
\item 求 $\F[H(t-a)]$，并解释为什么 $\delta$ 项不受相位因子影响。
\item 解释 $\F(tT)=i\partial_\omega\widehat T$ 中的 $i$，与分布导数定义里的负号怎样配合。
\item 选一条你想放进视频的等式，逐项说明它在哪个空间成立、有没有隐藏的交换极限步骤。
\end{enumerate}

\section{详解}

\subsection*{1. 两列趋于无穷的 $T$}

设 $a=|\omega|>0$。取

\begin{equation}
T_n=\frac{\pi/2+2\pi n}{a},
\qquad
S_n=\frac{3\pi/2+2\pi n}{a},
\qquad n=0,1,2,\dots
\end{equation}

则 $aT_n=\frac{\pi}{2}+2\pi n$，$aT_n'=\frac{3\pi}{2}+2\pi n$，于是 $\sin(aT_n)=1$，$\sin(aS_n)=-1$。而 $\sin(\omega T)=\pm\sin(aT)$，故

\begin{equation}
\frac{2\sin(\omega T_n)}{\omega}=\pm\frac{2}{a},
\qquad
\frac{2\sin(\omega S_n)}{\omega}=\mp\frac{2}{a}
\end{equation}

（符号取决于 $\omega$ 的正负）。两列都趋于无穷，极限点却不同，所以极限不存在。

\subsection*{2. 固定小区间里的面积}

\begin{equation}
\int_{-a}^{a}G_\ep(\omega)\dd\omega
=2\ep\cdot\frac{1}{\ep}\Big[\arctan\frac{\omega}{\ep}\Big]_{-a}^{a}
=4\arctan\frac{a}{\ep}\;\longrightarrow\;4\cdot\frac{\pi}{2}=2\pi .
\end{equation}

关键在\textbf{次序}：先固定 $a$，再让 $\ep\downarrow0$。这样 $\arctan(a/\ep)\to\pi/2$，整个区间装下全部面积。

若反过来先让 $a\to\infty$，则对每个固定 $\ep$ 也得到 $2\pi$——但那说明的是“全实轴面积”，与“集中”无关。

\subsection*{3. $\delta$ 及其导数的取样}

$\varphi(0)=1$。求导：

\begin{equation}
\varphi'(x)=e^{-x^{2}}(-2x)+e^{-x^{2}}+x e^{-x^{2}}(-2x)=e^{-x^{2}}\big(1-2x-2x^{2}\big),
\end{equation}

所以 $\varphi'(0)=1$。再求导并在 $0$ 处取值，得 $\varphi''(0)=-2$。

因此

\begin{equation}
\ip{\delta}{\varphi}=1,
\qquad
\ip{\delta'}{\varphi}=-\varphi'(0)=-1,
\qquad
\ip{\delta''}{\varphi}=\varphi''(0)=-2 .
\end{equation}

\subsection*{4. $\F[3-2t+t^{2}]$}

由线性性：

\begin{equation}
\F[3]-2\F[t]+\F[t^{2}]
=3\cdot2\pi\delta-2\cdot2\pi i\,\delta'+\big(-2\pi\delta''\big)
=6\pi\delta-4\pi i\,\delta'-2\pi\delta'' .
\end{equation}

检查：$3+t^{2}$ 是实偶的，对应实的偶分布 $6\pi\delta-2\pi\delta''$；$-2t$ 是实奇的，对应纯虚的奇分布 $-4\pi i\delta'$。

\subsection*{5. 两条乘法恒等式与阶跃}

\begin{align}
\ip{x\delta'}{\varphi}&=\ip{\delta'}{x\varphi}=-(x\varphi)'(0)=-\varphi(0)=-\ip{\delta}{\varphi},\\
\left\langle x\,\PV\tfrac1x,\varphi\right\rangle
&=\left\langle \PV\tfrac1x,x\varphi\right\rangle
=\int_{0}^{\infty}\frac{x\varphi(x)-(-x)\varphi(-x)}{x}\dd x
=\int_{0}^{\infty}\big(\varphi(x)+\varphi(-x)\big)\dd x=\int_{\R}\varphi .
\end{align}

所以 $x\delta'=-\delta$，$x\PV\frac1x=1$。于是

\begin{equation}
i\omega\Big[\pi\delta-i\,\PV\frac{1}{\omega}\Big]
=i\pi\,\underbrace{\omega\delta}_{=0}+\underbrace{\omega\,\PV\frac{1}{\omega}}_{=1}=1 .
\end{equation}

这验证了 $i\omega\widehat H=1$，同时说明为什么它\textbf{不能}唯一确定 $\widehat H$：任何 $\widehat H+c\delta$ 都给同一个结果。

\subsection*{6. $H(0)$ 无关}

$H_1$ 与 $H_2$ 只在 $t=0$ 处不同（例如 $0$ 与 $1$）。对任意 $\varphi\in\D$，

\begin{equation}
\ip{H_1}{\varphi}-\ip{H_2}{\varphi}=\int_{\R}\big(H_1-H_2\big)\varphi\dd t=0,
\end{equation}

因为差函数只在单点上非零，而单点集的 Lebesgue 测度为零，积分看不见它。所以两个函数定义同一个分布。

\subsection*{7. 为什么没有公共控制}

反证：若存在 $g\in\Lone$ 使 $K_\ep\le g$ 对一切 $\ep$，则由 $K_\ep(x)\to0$（$x\neq0$）与控制收敛，

\begin{equation}
1=\int K_\ep\longrightarrow\int 0=0,
\end{equation}

矛盾。

换元 $x=\ep u$ 之后：

\begin{equation}
\int_{\R}K_\ep(x)\varphi(x)\dd x=\frac{1}{\pi}\int_{\R}\frac{\varphi(\ep u)}{1+u^{2}}\dd u,
\end{equation}

被积函数被 $\|\varphi\|_\infty/(\pi(1+u^{2}))$ 控制，\textbf{这个控制与 $\ep$ 无关且可积}。所以控制收敛可以使用——注意控制的对象已经换成了 $\varphi(\ep u)$，而不是原来的 $K_\ep$。

\subsection*{8. 面积 $1$ 但极限为零}

\begin{equation}
k_n=\frac{1}{2}\mathbf{1}_{[n-1,\,n+1]}.
\end{equation}

$\int k_n=1$，但对 $\varphi\in\Sv$，

\begin{equation}
\ip{k_n}{\varphi}=\frac{1}{2}\int_{n-1}^{n+1}\varphi(x)\dd x=O\!\left(\frac{1}{n^{2}}\right)\longrightarrow0 .
\end{equation}

所以 $k_n\to0$ 于 $\Sv'$。

不矛盾：\textbf{“总质量为 $1$”这个条件在 $\Sv'$ 里不是连续的测试量}，因为常数函数 $1\notin\Sv$。质量守恒了，但它跑到了无穷远，任何快速衰减的探针都摸不到它。

\subsection*{9. 对偶定义证两个结果}

\begin{align}
\ip{\F\delta}{\varphi}&=\ip{\delta}{\ph}=\ph(0)=\int_{\R}\varphi(t)\dd t=\ip{1}{\varphi}
&&\Longrightarrow \F\delta=1;\\
\ip{\F1}{\varphi}&=\ip{1}{\ph}=\int_{\R}\ph(t)\dd t=2\pi\varphi(0)=\ip{2\pi\delta}{\varphi}
&&\Longrightarrow \F1=2\pi\delta .
\end{align}

第二个式子用到了\textbf{反演}在 $t=0$ 处的取值：$\varphi(0)=\frac{1}{2\pi}\int\ph\dd t$。\textbf{这是整个证明唯一的“外部输入”}，也是为什么反演必须先证。

\subsection*{10. 平移后的阶跃}

由时间平移规则 $\F[f(t-a)]=e^{-ia\omega}\fh(\omega)$，

\begin{equation}
\F[H(t-a)](\omega)=e^{-ia\omega}\Big[\pi\delta(\omega)-i\,\PV\frac{1}{\omega}\Big]
=\pi\delta(\omega)-i\,e^{-ia\omega}\PV\frac{1}{\omega}.
\end{equation}

$\delta$ 项不受相位影响，因为 $e^{-ia\omega}$ 是光滑乘子，而

\begin{equation}
e^{-ia\omega}\,\delta(\omega)=e^{-ia\cdot0}\,\delta(\omega)=1\cdot\delta(\omega)=\delta(\omega).
\end{equation}

一般规律：\textbf{光滑函数乘 $\delta$，只取它在原点的值。}主值项则要整体保留相位因子，因为 $\PV\frac1\omega$ 不是“集中在一点”的对象。

\subsection*{11. 那个 $i$ 与那个负号}

从三步看：

\begin{enumerate}[leftmargin=1.8em,itemsep=0.4em]
\item 由 $\F[\varphi']=i\omega\ph(\omega)$，把变量写成 $t$：$(\varphi')^{\wedge}(t)=it\,\ph(t)$，所以
\begin{equation}
t\,\ph(t)=\frac{1}{i}(\varphi')^{\wedge}(t)=-i\,(\varphi')^{\wedge}(t).
\end{equation}
这里出现了第一个负号（来自 $1/i=-i$）。

\item 分布导数的定义带来第二个负号：
\begin{equation}
\ip{\widehat T'}{\varphi}=-\ip{\widehat T}{\varphi'} .
\end{equation}

\item 两个负号相乘得正：
\begin{equation}
\ip{\F(tT)}{\varphi}=\ip{T}{t\ph}
=-i\,\ip{T}{(\varphi')^{\wedge}}
=-i\,\ip{\widehat T}{\varphi'}
=i\,\ip{\widehat T'}{\varphi}.
\end{equation}
\end{enumerate}

所以最终是 $+i\,\partial_\omega\widehat T$，不是 $-i$。

\subsection*{12. 一条等式的“体检报告”}

以 $\F[H]=\pi\delta+\PV\frac{1}{i\omega}$ 为例，逐项交代：

\begin{center}
\small
\renewcommand{\arraystretch}{1.45}
\begin{tabularx}{\textwidth}{@{}L{2.4cm}L{3.6cm}Y@{}}
\toprule
要素 & 内容 & 说明 \\
\midrule
等号在哪成立 & $\Sv'$ & 两边都是温和分布 \\
左边的来历 & $H_\ep\to H$ 于 $\Sv'$ & 控制函数 $|\varphi|$ \\
右边的来历 & $1/(\ep+i\omega)$ 的两项极限 & 实部 $\to\pi\delta$，虚部 $\to-i\PV$ \\
隐藏的交换 & 极限与变换交换 & 由 $\F:\Sv'\to\Sv'$ 连续（命题 \ref{prop:F-cont}）\\
隐藏的唯一性 & 两个极限相等 & 分布极限唯一 \\
不能做的事 & 不能逐点把 $\frac{1}{\ep+i\omega}$ 换成 $\frac{1}{i\omega}$ & 后者在原点不可积，不是函数 \\
\bottomrule
\end{tabularx}
\end{center}

\section{各章自测的答案要点}

\subsection*{EP01}

\begin{enumerate}[leftmargin=1.8em,itemsep=0.3em]
\item 见 1.5 节；$\ep>0$ 用在两处边界项消失（$e^{\ep t}\to0$ 于 $-\infty$，$e^{-\ep t}\to0$ 于 $+\infty$）。
\item 解 $\frac{2\ep}{\ep^{2}+\omega^{2}}=\frac{1}{\ep}$ 得 $\omega^{2}=\ep^{2}$，故半高全宽 $2\ep$。
\item 用 $\arctan$，得 $2\pi$；不随 $\ep$ 变是因为积分只依赖比值 $\omega/\ep$。
\item $G_\ep(\omega)=\frac{1}{\ep}\hat g(\omega/\ep)$，其中 $\hat g(\omega)=2/(1+\omega^{2})$。
\item 峰高 $20$ 超出纵轴上限，峰被裁掉，观众会误以为“没有峰”。
\item 面积是积分、点值是函数值；这两件事在 $\ep\to0$ 时的行为本来就可以不同。
\end{enumerate}

\subsection*{EP02}

\begin{enumerate}[leftmargin=1.8em,itemsep=0.3em]
\item 依次为 $1$、$-1$、$-2$。
\item $6\pi\delta-4\pi i\delta'-2\pi\delta''$。
\item $i^{4}=1$，故 $\F[t^{4}]=2\pi\delta^{(4)}$。
\item $\int K_\ep'=K_\ep(\infty)-K_\ep(-\infty)=0$；总积分为零，故“面积”不能刻画它。
\item 泰勒展开把 $\varphi$ 拆成“零点值 + 各阶导数项 + $x^{N+1}\psi$”；$T$ 作用在 $x^{N+1}\psi$ 上为零，因为该项在原点邻域内有足够高阶的零点而 $T$ 支撑在 $\{0\}$。
\item $\omega^{N+1}\widehat p=0$ 只说明频谱住在原点（由支撑命题），不给出各 $\delta^{(n)}$ 的系数。
\item 参考：“$\delta$ 关心值，$\delta'$ 关心斜率；它们都只看原点，但问的不是同一件事。”
\end{enumerate}

\subsection*{EP03}

\begin{enumerate}[leftmargin=1.8em,itemsep=0.3em]
\item 改一个点不改变分布（见详解第 6 题）。
\item $i\pi\omega\delta+\omega\PV\frac1\omega=0+1=1$。
\item 两个解：$\pi\delta-i\PV\frac1\omega$ 与它加 $c\delta$。
\item $\varphi(x)-\varphi(-x)=2xe^{-x^{2}}$，故积分 $=\int_0^\infty2e^{-x^{2}}\dd x=\sqrt\pi$。
\item 差 $\varphi(0)\ln(b/a)$；只有 $b/a\to1$ 才给出主值。
\item $\F[\sgn t]=2(\widehat H-\pi\delta)=-2i\PV\frac1\omega$。
\item $\frac{1}{\ep+i\omega}=\frac{-i}{\omega-i\ep}\to\frac{-i}{\omega-i0}=-i\big(\PV\frac1\omega+i\pi\delta\big)=\pi\delta-i\PV\frac1\omega$。
\end{enumerate}

\subsection*{EP04}

\begin{enumerate}[leftmargin=1.8em,itemsep=0.3em]
\item $f_n=\frac1n\mathbf{1}_{[0,n]}$：一致趋零，积分恒为 $1$。不违反控制收敛，因为不存在公共可积控制（否则会推出积分趋于零）。
\item 见详解第 7 题。
\item 控制函数是 $\|\varphi\|_\infty/(\pi(1+u^{2}))$。
\item $1\notin\Lone$，控制不可积；乘 $\varphi$ 后控制变成 $|\varphi|\in\Lone$。
\item 质量 $1$、一致 $\Lone$ 有界、尾部集中；分别防止“变轻”“失控”“变宽”。
\item 见详解第 8 题；不矛盾，因为 $1\notin\Sv$。
\item 成立；需要 $\int k_\ep=1$、$\sup_\ep\|k_\ep\|_1<\infty$、尾部集中。估计时把 $k_\ep$ 换成 $|k_\ep|$。
\end{enumerate}

\subsection*{EP05}

\begin{enumerate}[leftmargin=1.8em,itemsep=0.3em]
\item 见详解第 5 题；同理 $x\delta''=-\delta'$。
\item $(2+x)\delta''=2\delta''-2\delta'$；$x^{2}\delta''=2\delta$。
\item 见 5.5 节；$H$ 不光滑，$H\delta$ 依赖正则化，没有内禀定义。
\item $\eta$ 紧支撑且所有导数在 $\pm1$ 处趋于零；$e^{-x^{2}}$ 无紧支撑。
\item $\delta(3x-6)=\delta(3(x-2))=\frac13\delta(x-2)$。
\item 若探针以“所有低阶导数一致趋零”的方式变小，而响应不趋于零，就违反 \eqref{eq:local-est} 给出的半范数估计。
\item 见 5.4 节；最后一步用了“单点集测度为零，故 $g=0$ 几乎处处于 $\R$”。
\end{enumerate}

\subsection*{EP06}

\begin{enumerate}[leftmargin=1.8em,itemsep=0.3em]
\item 整函数在开集上为零则处处为零；若 $\F\varphi$ 紧支撑则 $\varphi\equiv0$。
\item 由 $\F[g^{(m)}]=(i\omega)^{m}\hat g$ 直接整理。
\item 由 \eqref{eq:Schwartz-bound}，$p_{m,n}(\ph)$ 被有限个 $\|t^{k}\partial^{j}\varphi\|_1$ 控制。
\item 四步：高斯积分 \eqref{eq:gauss-int}、Fubini（高斯因子保证绝对收敛）、写成卷积 \eqref{eq:conv-gauss}、两边分别取极限（左边用控制收敛，右边用近似恒等核）。
\item $\F^{2}\varphi=2\pi\varphi(-\cdot)$ 就是反演的变形；$\F\delta=1$ 由对偶一步得到，两者结合给出 $\F1=2\pi\delta$。
\item $\ip{\F1}{\varphi}=\int\ph\dd t=2\pi\varphi(0)$；这一步用的是反演在 $0$ 处的值。
\item 见命题 \ref{prop:ex2}：指数增长压不过多项式增长的半范数。
\item sinc 不正、不绝对可积地集中；它靠 $w_T\to1$ 于 $\Sv'$ 加上 $\F$ 连续而收敛。
\end{enumerate}

\section{延伸阅读}

\begin{ideabox}[与本书框架最接近的两份讲义]
\begin{enumerate}[leftmargin=1.8em,itemsep=0.3em]
\item S. Dyatlov, \emph{Lecture Notes for 18.155: Distributions, Elliptic Regularity, and Applications to PDEs}, MIT. 重点看“分布与运算”“傅里叶变换”两部分。\\
{\small\ttfamily https://math.mit.edu/\~{}dyatlov/18.155/155-notes.pdf}
\item G. Hörmann, R. Steinbauer, \emph{Lecture Notes on the Theory of Distributions}, University of Vienna, 2009. 测试函数、分布导数、Schwartz 空间；对阶跃变换给了两种推导。\\
{\small\ttfamily https://www.mat.univie.ac.at/\~{}stein/teaching/courses/SoSem09/distrvo.pdf}
\end{enumerate}
\end{ideabox}

\noindent 想再往深处走：Rudin, \emph{Functional Analysis}（严格陈述）；Strichartz, \emph{A Guide to Distribution Theory and Fourier Transforms}（从直觉走向分布）；Gel'fand--Shilov, \emph{Generalized Functions} Vol.~1（正则化思路）；Lighthill, \emph{Introduction to Fourier Analysis and Generalised Functions}（篇幅小、物理味重）。

\begin{warnbox}[引用前必做的一件事]
各资料的 $2\pi$ 归一化、$\delta$ 系数、$\sin$ 符号可能不同。\textbf{引用前先算一遍该资料的 $\F[1]$}，用它对齐约定（见 0.6 节）。差一个常数时，先怀疑约定，再怀疑内容。
\end{warnbox}
}{}

\end{document}
